% !TeX root = ../main.tex
\chapter{Intro}
\underline{Riemann-Stieltjes Integral}: If $f$ is continuous and $g$ is of \underline{finite variation}, then $\int_{a}^{b} f(x) ~ \mathrm{d}g(x)$ is well-defined, and is given by
$$
    \int_{a}^{b} f(x) ~ \mathrm{d}g(x):=\lim_{n\to \infty} \sum_{i=0}^{n-1}f(\xi_{i})\left(g(x_{i+1}) - g(x_{i})\right)
$$
where $a \leq x_{0} < x_{1} < \cdots < x_{n} \leq b$ and $\xi_{i} \in [x_{i}, x_{i+1}]$.

\underline{Main Task}: define the It\^o integral
$$
    (H\cdot X)_{t} = \int_{0}^{t} H_{s} ~\mathrm{d}X_{s}.
$$

The most important case is when $X$ is Brownian motion.

\thmnn{$\displaystyle\sum_{k=1}^{2^{n}}\left(B_{\frac{k}{2^{n}}} - B_{\frac{k-1}{2^{n}}}\right)^{2} \xlongrightarrow{a.s.} 1 .$}

Recall that the finite variation of a function $g:[a,b] \to \mathbb{R}$ is
$$
    V_{g}([a,b]) = \int_{a}^{b} \abs{ \mathrm{d}g(x)} = \sup_{n}\sum_{i=1}^{n}\abs{g(x_{i+1}) - g(x_{i})}
$$

\cornnp{$V_{B}([0,1]) = \infty$ a.s.}{
    \begin{equation*}
        \begin{aligned}
                    \textcolor{red}{(\star)}\sum_{k=1}^{2^{n}}\left( B_{\frac{k}{2^{n}}} - B_{\frac{k-1}{2^{n}}}\right)^{2} &\leq \sup_{1 \leq k \leq 2^{n}} \abs{B_{\frac{k}{2^{n}}} - B_{\frac{k-1}{2^{n}}}}\cdot \sum_{k=1}^{2^{n}}\abs{B_{\frac{k}{2^{n}}} - B_{\frac{k-1}{2^{n}}}}\\
        & \leq V_{B}([0,1]) \cdot \sup_{1 \leq k \leq 2^{n}}\abs{B_{\frac{k}{2^{n}}} - B_{\frac{k-1}{2^{n}}}}
        \end{aligned}
    \end{equation*}
    If $V_{B}([0,1]) < \infty$, then by the continuity of Brownian paths,
    $$
        \sup_{1 \leq k \leq 2^{n}}\abs{B_{\frac{k}{2^{n}}} - B_{\frac{k-1}{2^{n}}}} \to 0,
    $$
    and hence $\textcolor{red}{(\star)} \to 0$.
    
    This contradicts $\textcolor{red}{(\star)} \to 1$. Therefore, with probability one,
    $$
        V_{B}([0,1]) = \infty.
    $$
}

This shows that Brownian motion is too irregular to be treated by the classical Riemann--Stieltjes theory. Therefore, to define
$$
    (H\cdot X)_{t} = \int_{0}^{t} H_{s}~\mathrm{d}X_{s},
$$
we need a new framework, where typically $H$ is previsible (predictable) and $X$ is a semimartingale.

\ex{}{Let $H_{n} \in \mathcal{F}_{n-1}$ and let $(M_{n})$ be a martingale.

Then $(H\cdot M)_{n} = \sum_{k=1}^{n}H_{k}(M_{k} - M_{k-1})$ is a martingale.}

\section{Set-up}

Throughout, let $(\Omega, \mathcal{F}, \mathsf{P})$ be a probability space. A stochastic process $X=(X_{t})_{t \geq 0}$ is a family of random variables taking values in $(\mathbb{R}^{d},\mathcal{B}(\mathbb{R}^{d}))$:
$$
\begin{aligned}
    X\colon [0,\infty) \times \Omega &\to \mathbb{R}^{d}\\
   (t,\omega) &\mapsto X_{t}(\omega).
\end{aligned}
$$

For each fixed $\omega$, the map $t \mapsto X_{t}(\omega)$ is called the sample path of $X$ corresponding to $\omega$.

\defnn{For two stochastic processes $X$ and $Y$, we say that $Y$ is a modification of $X$ if
$$
    \forall~ t \geq 0,\quad \mathsf{P}(X_{t}= Y_{t}) = 1.
$$}

\defnn{We say that $X$ and $Y$ are indistinguishable if
$$
    \mathsf{P}(X_{t} = Y_{t},~ \forall~ t \geq 0) = 1.
$$
In this case, $Y$ is also called a version of $X$.}

\ex{}{Let $Y$ be a modification of $X$, and suppose that both $X$ and $Y$ have a.s. right-continuous sample paths. Then $X$ and $Y$ are indistinguishable.}

\defnn{A function $f$ is right-continuous if
$$
    f(x)= f(x+) = \lim_{t\downarrow x} f(t).
$$}

\pf{Since $\forall~  t \geq 0$, $\mathsf{P}(X_{t} \neq Y_{t}) = 0$, the two processes agree almost surely at every fixed rational time. More precisely,
    $$
        \mathsf{P}\left(\bigcup_{q\in \mathbb{Q}^{+}} X_{q} \neq Y_{q}\right) \leq \sum_{q\in \mathbb{Q}^{+}} \mathsf{P} (X_{q} \neq Y_{q}) = 0.
    $$
    Therefore,
    $$
        \mathsf{P}\left(\bigcap_{q\in \mathbb{Q}^{+}} X_{q} = Y_{q}\right) = 1.
    $$
    On the other hand, if two right-continuous functions $f$ and $g$ agree on $\mathbb{Q}^{+}$, then they must agree on all of $[0,\infty)$, since the rationals are dense.

    Applying this observation pathwise to the sample paths of $X$ and $Y$, we conclude that
    $$
        \mathsf{P}\left(X_{t} = Y_{t},~ \forall~ t \geq 0\right) = 1.
    $$
}

\ex{}{Let $T \sim \mathrm{Exp}(1)$. Define 
    $$
    \begin{aligned}
        &X_{t} \equiv 0,\quad\forall~ t \geq 0 \\
        &Y_{t} = \begin{cases}
            0,& t \neq T\\
            1,&t = T.
        \end{cases}
    \end{aligned}
    $$
    Then for every fixed $t \geq 0$,
    $$
        \mathsf{P}(X_{t} \neq Y_{t}) = \mathsf{P}(Y_{t} = 1) = \mathsf{P}(t = T) = 0.
    $$
    Hence $Y$ is a modification of $X$. However, $X$ and $Y$ are not indistinguishable, because on the event $\{T<\infty\}$ they differ at time $T$.
}

\section{Filtration}

We now introduce filtrations, which describe the information available up to each time $t$.

Let $(\mathcal{F}_{t})_{t \geq 0}$ be a filtration of $(\Omega, \mathcal{F})$, that is,
$$
    \forall~ 0 \leq s < t < \infty,\quad \mathcal{F}_{s} \subseteq \mathcal{F}_{t} \subseteq \mathcal{F}.
$$

\ex{}{(Natural filtration) Let $(X_{t})_{t \geq 0}$ be a stochastic process. Define
    $$
        \mathcal{F}_{t}^{X} = \sigma(X_{s},0 \leq s \leq t).
    $$
    Then $(\mathcal{F}_{t}^{X})$ is a filtration.
}
By construction, $X_{t} \in \mathcal{F}_{t}^{X}$ for every $t \geq 0$.

\exe{1.7}{Let $X$ be a process whose sample paths are all \underline{RCLL}. Let
    $$
        A = \{\omega\colon X(\omega)\text{ is continuous on }[0,t_{0})\}
    $$    
    for some fixed $t_{0} > 0$. Show that $A \in \mathcal{F}_{t_{0}}^{X}.$
}

\underline{RCLL}: $f(x) = f(x+)$ and $f(x-)$ exists.

\pf{
    Since every sample path is RCLL, continuity on $[0,t_0)$ is equivalent to continuity at every point in $(0,t_0)$. Hence
    $$
    \begin{aligned}
        A &= \{\omega\colon \forall~ s\in (0,t_{0}) ,~ X(s-,\omega)= X(s,\omega) \}\\
        &=\{\omega\colon \forall~q \in (0,t_{0})\cap \mathbb{Q},~X(q-,\omega) = X(q,\omega) \}
    \end{aligned}
    $$
    because for an RCLL function, possible discontinuities are jumps, and it is enough to check them on rational times.

    Let $(0, t_{0})\cap \mathbb{Q} = (q_{m})_{m \geq 1}$. Then
    $$
    \begin{aligned}
        A &= \left\{
            \omega \colon \forall~ q_{m}\text{ with } m \geq 1,~\forall~k \geq 1,~\exists~ N \geq 1,~\forall~ n \geq N, \abs{X(q_{m}-\frac{1}{n},\omega) - X(q_{m},\omega)} < \frac{1}{k}
        \right\}\\
        &= \bigcap_{m=1}^{\infty}\bigcap_{k=1}^{\infty}\bigcup_{N=1}^{\infty}\bigcap_{n=N}^{\infty}\left\{\omega\colon \abs{X(q_{m} - \frac{1}{n},\omega) - X(q_{m},\omega)} < \frac{1}{k}\right\}\in \mathcal{F}_{t_{0}}^{X}.
    \end{aligned}
    $$
    Therefore $A \in \mathcal{F}_{t_0}^X$.
}

\exe{1.8}{Let $X$ be a process whose sample paths are RCLL a.s. Then $A \notin \mathcal{F}_{t_{0}}^{X}$ is possible.}

\pf{
    [TODO]
}

\defnn{A filtration $\{\mathcal{F}_{t}\}_{t \geq 0}$ is said to satisfy the \underline{usual conditions} if
    \begin{enumerate}[label=(\alph*)]
        \item $\displaystyle \mathcal{F}_{t} = \mathcal{F}_{t+} = \bigcap_{s>t} \mathcal{F}_{s}$;
        \item $\mathcal{F}_{0}$ contains all the \underline{$\mathsf{P}-$negligible} sets in $\mathcal{F}$.
    \end{enumerate}
}

\defnn{An event $A$ is called $\mathsf{P}$-negligible if there exists $N \in \mathcal{F}$ such that $A \subseteq N$ and $\mathsf{P}(N) = 0$.}

\defnn{
    We define
    $$
        \mathcal{F}_{t+} = \bigcap_{s >t}\mathcal{F}_{s},
        \qquad
        \mathcal{F}_{t-}= \sigma\left(\bigcup_{s<t}\mathcal{F}_{s}\right).
    $$
    We say that the filtration is right-continuous if $\mathcal{F}_{t} = \mathcal{F}_{t+}$.
}

\defnn{
    The stochastic process $X$ is \underline{adapted} to the filtration $\{\mathcal{F}_{t}\}_{t \geq 0}$ if $\forall~t \geq 0$, $X_{t}$ is $\mathcal{F}_{t}$-measurable.
}

In particular, every process is adapted to its natural filtration.

\defnn{The stochastic process $X$ is called progressively measurable with respect to $(\mathcal{F}_{t})$ if $\forall~t \geq 0$, the map
    $$
        \begin{aligned}
            (s,\omega) &\mapsto X_{s}(\omega)\\
            ([0,t]\times \Omega,\mathcal{B}([0,t])\otimes \mathcal{F}_{t}) & \to (\mathbb{R}^{d}, \mathcal{B}(\mathbb{R}^{d}))
        \end{aligned}
    $$
    is measurable.
}

\rmk{An adapted process need not be progressively measurable in general. The useful point here is that, under additional path regularity, adaptedness implies progressive measurability. In the usual hierarchy one has
$$
    \text{previsible} \Longrightarrow \text{progressively measurable} \Longrightarrow \text{adapted}.
$$
}



\propp{If $X$ is adapted to the filtration ($\mathcal{F}_t$) and every sample path is right-continuous or left-continuous, then $X$ is also progressively measurable.}{
    It is enough to treat the right-continuous case. Fix $t>0$. For every $n \geq 1$ and $0 \leq k \leq 2^n -1$, define
    $$
    X_s^{(n)}(\omega) = X_{\frac{k+1}{2^n}t}(\omega),\quad \forall~  \frac{kt}{2^n} < s \leq \frac{k+1}{2^n}t
    $$
    Then, by right-continuity of the sample paths,
    $$
        \lim_{n \to \infty}X_{s}^{(n)}(\omega) = X_{s}(\omega)
    $$
    for all $(s,\omega) \in [0,t]\times \Omega$.

    Therefore, it suffices to show that
    $$
    \begin{aligned}
        (s,\omega) &\mapsto X_{s}^{(n)}(\omega)\\
        ([0,t]\times \Omega,\mathcal{B}([0,t])\otimes \mathcal{F}_{t}) & \to (\mathbb{R}^{d}, \mathcal{B}(\mathbb{R}^{d}))
    \end{aligned}
    $$
    is measurable.

    Indeed, for every $A \in \mathcal{B}(\mathbb{R}^{d})$,
    $$
    \begin{aligned}
        &\left\{(s,\omega)\colon X_{s}^{(n)}(\omega) \in A\right\}\\
        =&\bigcup_{k=0}^{2^{n} -1}\left(\frac{k}{2^{n}}t , \frac{k+1}{2^{n}}t \right]\times \left\{X_{\frac{k+1}{2^{n}}t} \in A\right\}\in \mathcal{B}([0,t])\otimes \mathcal{F}_{t}
    \end{aligned}
    $$
    because $X$ is adapted and therefore $X_{\frac{k+1}{2^n}t}$ is $\mathcal{F}_t$-measurable.

    Thus $X_{s}^{(n)}(\omega)$ is measurable for every $n$. Passing to the pointwise limit, we conclude that $X$ is progressively measurable.
}

\section{Stopping Times}
We now study random times that are compatible with the filtration $(\mathcal{F}_{t})_{t \geq 0}$ on $(\Omega, \mathcal{F},\mathsf{P})$.

\defnn{A random time $T$ is called a \underline{stopping time} with respect to $(\mathcal{F}_{t})$ if
$$
    \forall~ t \geq 0,\quad \{T \leq t\} \in \mathcal{F}_{t}.
$$
}

\defnn{A random time $T$ is called an \underline{optional time} with respect to $(\mathcal{F}_{t})$ if
$$
    \forall~ t \geq 0,\quad \{T < t\} \in \mathcal{F}_{t}.
$$
}

\propp{Every stopping time is an optional time. Moreover, if the filtration is right-continuous, then the two notions are equivalent.}{
    ``$\Rightarrow$'' If $T$ is a stopping time, then
    $$
        \{T<t\} = \bigcup_{n=1}^{\infty}\left\{T \leq t - \frac{1}{n}\right\} \in \mathcal{F}_t.
    $$
    
    ``$\Leftarrow$'' Now assume that $T$ is an optional time and that the filtration is right-continuous. Then for every $m \geq 1$,
    $$
        \{T \leq t\} = \bigcap_{n=m}^{\infty}\left\{T < t+\frac{1}{n}\right\}.
    $$
    Since $\{T < t + \frac{1}{n}\} \in \mathcal{F}_{t+\frac{1}{n}} \subseteq \mathcal{F}_{t+\frac{1}{m}}$, we obtain
    $$
        \{T \leq t\} \in \mathcal{F}_{t+\frac{1}{m}}.
    $$
    
    Therefore
    $$
        \{T \leq t\} \in \bigcap_{m=1}^{\infty}\mathcal{F}_{t+\frac{1}{m}}= \mathcal{F}_{t^+} = \mathcal{F}_t,
    $$
    so $T$ is a stopping time.
}

\lemnn{If $T$ and $S$ are stopping times, then $T \wedge S$, $T \vee S $, $T+S$ are stopping times.}

\pf{$\forall~ t \geq 0$, 
    \begin{enumerate}
        \item $\{T \wedge S \leq t\} = \{T \leq t\} \cup \{S \leq t\} \in \mathcal{F}_{t}.$
        \item $\{ T \vee S \leq t\} = \{T \leq t\}\cap \{S \leq t\} \in \mathcal{F}_{t}.$
        \item To prove $\{T+S \leq t\} \in \mathcal{F}_{t}$, it is enough to show that $\{T+S > t\} \in \mathcal{F}_{t}$.
        $$
        \begin{aligned}
            &\{T + S > t\} \\
            = &\{T = 0,~S >t\}\cup \{S=0,~ T> t\}\cup\{0<T<t,~T+S > t\}\cup \{S>0,~ T \geq t\}\\
             \in &\mathcal{F}_{t} \cup \mathcal{F}_{t} \cup \{0<T<t,~T+S > t\} \cup  \mathcal{F}_{t}.
        \end{aligned}
        $$
        Thus it remains to show that
        $$\{0< T < t ,~ T+S > t\} = \bigcup_{\stackrel{r\in \mathbb{Q}^{+}}{0 < r <t}}\{r < T < t,~ S> t-r\}.
        $$
        The right-hand side belongs to $\mathcal{F}_{t}$.
        
        ``$\Leftarrow~$'' $S > t-r > t - T$, so $S+T > t$.

        ``$\Rightarrow~$'' Since $S+T > t$, we can choose $r\in \mathbb{Q}^{+}$ with $0<r<T$ and $S+r>t$.
    \end{enumerate}
}

\lemnn{If $\{T_n\}_{n \geq 1}$ are stopping times, then $\sup_{n \geq 1} T_n$ is a stopping time.}

\pf{
    $\forall~ t \geq 0,$ 
    $
        \{\sup_{n \geq 1} T_n \leq t\} = \bigcap_{n=1}^{\infty}\{T_n \leq t\} \in \mathcal{F}_t.
    $
}

\defnn{The stopping $\sigma$-algebra associated with a stopping time $T$ is
$$
    \mathcal{F}_T = \{A \in \mathcal{F} \mid \forall~ t \geq 0,\; A \cap \{T \leq t\} \in \mathcal{F}_t\}.
$$}

\prop{If $T(\omega) = t_0$ for some fixed $t_0 \geq 0$, $\forall~ \omega \in \Omega$, then $\mathcal{F}_T = \mathcal{F}_{t_0}$.}

\pf{Note that $\forall~  t \geq 0$,
    $$
        \{T \leq t\} = \begin{cases}
            \varnothing,& t_0 > t\\
            \Omega, &t_0 \leq t
        \end{cases}
    $$

    \noindent``$\mathcal{F}_{t_{0}} \subseteq \mathcal{F}_{T}$''

    $\forall~  A \in \mathcal{F}_{t_{0}}$, $\forall~  t \geq 0$,
    \begin{enumerate}[label=\textcircled{\arabic*}]
        \item $t \geq t_{0}$, $A \cap \{T \leq t\} = A \in \mathcal{F}_{t_{0}} \subseteq \mathcal{F}_{t};$
        \item $t < t_{0}$, $A \cap \{T \leq t\} = \varnothing \in \mathcal{F}_{t}$.
    \end{enumerate}
    Hence $\mathcal{F}_{t_{0}}\subseteq \mathcal{F}_{T}$.

     \noindent``$\mathcal{F}_{T} \subseteq \mathcal{F}_{t_{0}}$''

    $\forall~  A \in \mathcal{F}_{T}$, we have $A \cap \{T \leq t\} \in \mathcal{F}_{t}$ by definition. Let $t = t_{0}$, then $\{T \leq t_{0}\}= \Omega$, so

    $$
        A \cap \{T \leq t_{0}\} \in \mathcal{F}_{t_{0}} \Rightarrow~ A \cap \Omega \in \mathcal{F}_{t_{0}} \Rightarrow~ A \in \mathcal{F}_{t_{0}}.
    $$
    Hence $\mathcal{F}_{T}\subseteq \mathcal{F}_{t_{0}}$.
}
    
\lemnn{$\forall~$ stopping time $T$ and $S$, $\forall~A \in \mathcal{F}_S$, we have $A\cap \{S \leq T\} \in \mathcal{F}_T.$ In particular, if $S \leq T$ for all $\omega \in \Omega$, then $\mathcal{F}_S \subseteq \mathcal{F}_T.$}

\pf{
    First note that $T \wedge t$ is also a stopping time. Hence $T \wedge t$ is $\mathcal{F}_t$-measurable.

    To prove $A\cap \{S \leq T \} \in \mathcal{F}_T$, it is enough to show that for every $t \geq 0$,
    $$
        A \cap \{S \leq T\} \cap \{T \leq t\}\in \mathcal{F}_t.
    $$

    We have 
    $$
    A \cap \{S \leq T\} \cap \{T \leq t\} = \underbrace{A \cap \{S \leq t\}}_{\in \mathcal{F}_t} \cap \underbrace{\{T \leq t\}}_{\in \mathcal{F}_t} \cap \underbrace{\{S \wedge t \leq T \wedge t\}}_{\mathcal{F}_t}.
    $$
    Indeed, on $\{S \leq t\}\cap \{T \leq t\}$, the event $\{S \leq T\}$ is exactly the same as $\{S \wedge t \leq T \wedge t\}$. Therefore the whole intersection belongs to $\mathcal{F}_t$, and hence $A\cap \{S \leq T \} \in \mathcal{F}_T$.
}

\lemnn{
    Let $S,T$ be two stopping times. Then $\mathcal{F}_{T \wedge S} = \mathcal{F}_T \cap \mathcal{F}_S$ and $\{T < S\}$, $\{S < T\}$, $\{T \leq S\}$, $\{S \leq T\}$, $\{T = S\}$ belong to $\mathcal{F}_S \cap \mathcal{F}_T = \mathcal{F}_{S \wedge T}.$
}

\pf{
    $\mathcal{F}_{T \wedge S} \subseteq \mathcal{F}_S$ and $\mathcal{F}_{T \wedge S} \subseteq \mathcal{F}_T \Rightarrow~\mathcal{F}_{T \wedge S} \subseteq \mathcal{F}_T \cap \mathcal{F}_S$.
    
    To prove $\mathcal{F}_T \cap \mathcal{F}_S \subseteq \mathcal{F}_{T \wedge S}$, $\forall~ A \in \mathcal{F}_T \cap \mathcal{F}_S$, we will prove $\forall~ t \geq 0$, $A \cap \{T \wedge S \leq t\} \in \mathcal{F}_t.$

    $$
    \begin{aligned}
        A \cap \{T \wedge S \leq t\} &= A \cap \left(\{T \leq t\} \cup \{S \leq t\}\right)\\
        &= \left(A \cap \{T \leq t\}\right)\cup \left(A \cap \{S \leq t\}\right)\in \mathcal{F}_t.
    \end{aligned}
    $$

    Next, $\{S \leq T\}\in \mathcal{F}_T$. It remains to prove $\{S <T\} \in \mathcal{F}_T.$

    Assuming this, by interchanging $T$ and $S$, we get $\{T \leq S\}$, $\{T > S\}$, $\{T <S\}\in \mathcal{F}_S$.

    Therefore $\{S \leq T\} \in \mathcal{F}_S \cap \mathcal{F}_T$.

    $\{T= S\} = \{ S \leq T\}\cap \{T \leq S\}\in \mathcal{F}_S \cap \mathcal{F}_T$.

    To prove $\{S < T\}\in \mathcal{F}_T$, let $R = S \wedge T$. Then $\{S < T\}= \{R < T\}$. Since $R$ is a stopping time, we have $R \in \mathcal{F}_R \subseteq \mathcal{F}_T$, and therefore $\{R < T\} \in \mathcal{F}_T$.
}

\textbf{Summary:}
\begin{enumerate}[label=(\roman*)]
    \item $T \in \mathcal{F}_T$;
    \item If $S \leq T$, $\mathcal{F}_S \subseteq \mathcal{ F}_T$;
    \item $T \wedge t \in \mathcal{F}_{T \wedge t} \subseteq \mathcal{F}_t$;
    \item $\mathcal{F}_{T \wedge S} = \mathcal{F}_T \cap \mathcal{F}_S.$
\end{enumerate}

For each $\omega \in \Omega$, if $T(\omega) \in [0,\infty)$ and $X_t(\omega) \in \mathbb{R}^d$ for all $t \geq 0$, we define
$$
    X _T(\omega) = X _{T(\omega)}(\omega).
$$

\prop{
    Let $\{X_t\}_{t \geq 0}$ 
    be a progressively measurable process, and let $T$ be a stopping time. Then $X_T$, defined on $\{T < \infty\}$, is $\mathcal{F}_T$-measurable, and the stopped process $\{X _{T \wedge t}\}_{t \geq 0}$ is progressively measurable.
}

\pf{
    Let $B \in \mathcal{B}(\mathbb{R}^d)$. By definition of $\mathcal{F}_T$,
    $$
        \{X_T \in B\} \in \mathcal{F}_T
        \quad \Longleftrightarrow \quad
        \forall~ t \geq 0,\;
        \{ X _T \in B\}\cap \{T \leq t\}\in \mathcal{F}_t.
    $$
    On the event $\{T \leq t\}$ we have $X_T = X_{T\wedge t}$, so
    $$
        \{ X _T \in B\}\cap \{T \leq t\}
        =
        \{X _{T \wedge t} \in B\} \cap \{T \leq t\}.
    $$

    Therefore, it is enough to prove that $X_{T \wedge t}$ is $\mathcal{F}_t$-measurable for every $t \geq 0$.

    Observe that the map
    $$
        (s,\omega)\mapsto (T(\omega) \wedge s, \omega)
    $$
    is $\mathcal{B}([0,t])\otimes \mathcal{F}_t$-measurable.

    By assumption, $(s,\omega)\mapsto X_s(\omega)$ is $\mathcal{B}([0,t])\otimes \mathcal{F}_t$-measurable.

    Hence the composition
    $$
        (s,\omega)\mapsto X_{T(\omega) \wedge s}(\omega)
    $$
    is also $\mathcal{B}([0,t])\otimes \mathcal{F}_t$-measurable. In particular, evaluating at $s=t$ shows that $X_{T \wedge t}$ is $\mathcal{F}_t$-measurable.

    This also proves that the stopped process $(X_{T\wedge t})_{t \geq 0}$ is progressively measurable.
}

\section{Cts-time martingale}
We now turn to continuous-time martingales and some of their basic properties.

\defnn{On $(\Omega,\mathcal{F},(\mathcal{F}_{t})_{t \geq 0},\mathbb{P})$, a martingale is an adapted stochastic process $(X_t)_{t \geq 0}$ such that
    $$
    \begin{cases}
        (i)~ \forall~ t \geq 0,~\mathbb{E}|X_t| < \infty \\
        (ii)~ \forall~ t \geq s \geq 0,~\mathbb{E}(X_t \mid \mathcal{F}_s) = X _s.
    \end{cases}
    $$

    If $\mathbb{E}(X_t \mid \mathcal{F}_s) \geq X _s$, we call $(X_t)$ a submartingale.

    If $\mathbb{E}(X_t \mid \mathcal{F}_s) \leq X _s$, we call $(X_t)$ a supermartingale.
}

\prop{Let $\varphi: \mathbb{R} \to \mathbb{R}$ be a convex (concave up) function. If $(X_t)_{t \geq 0}$ is a martingale, then $(\varphi(X _t))_{t \geq 0}$ is a submartingale.}

\textbf{Examples.} $\abs{X_t}$, $X _{t}^2$, $\abs{X _t}^p$ for all $p \geq 1$, and $(X _t)^+$ are all submartingales.

\rmk{We assume $\mathbb{E}\abs{\varphi(X_t)}<\infty$, $\forall~ t \geq 0$.}

\pf{
    Jensen's inequality.
}

Define
    $$
    \begin{aligned}
    &T_{b}^{(0)} = -\infty,\\
    & T_a^{(1)} = \inf \{t \geq (T _b^{(0)})^+\colon X _t \leq a\},\\
    &T_b^{(1)} = \inf \{t \geq T _a^{(1)}\colon X _t \geq b\}.
    \end{aligned}
    $$

    $\forall~ l \geq 2$,
    $$
    \begin{aligned}
        &T_a^{(l)} = \inf \{t \geq T _b^{(l-1)} \colon X _t \leq a\},\\
        &T_b^{l} = \inf \{ t \geq T _a^{l}\colon X _t \geq b\}
    \end{aligned}
    $$
    Define $$U_{ab}^X [0,n] = \max \{k \geq 0\colon T_{b}^{(k)} \leq n\}$$ to be the number of upcrossings of the interval $[a,b]$ on $[0,n]$.
\thm{Continuous-time Doob's upcrossing inequality}{
     Let $(X_{t})_{t \geq 0}$ be a right-continuous submartingale. For all $T > 0$,
        $$
        \mathbb{E}\left(U_{ab}^X [0,T]\right) \leq \frac{\mathbb{E}(X_{T} -a )^+}{b-a}
        $$
}

\pf{
    For each $n \geq 1$, let $D_n = \mathbb{Z}/ 2^n = \{0, \frac{\pm 1}{2^n}, \frac{\pm 2}{2^n},\cdots\}$.

    Consider the upcrossings of $X$ on $[0,T]\cap D_{n}$. The sampled process $\left\{\frac{X_{j}}{2^n}\right\}_{j=0}^\infty$ is a discrete-time submartingale, so
    $$
        U_{ab}^X\left([0,T]\cap D_{n}\right) \leq U_{ab}^X [0,T],
        \qquad
        U_{ab}^X\left([0,T]\cap D_{n}\right) \uparrow U_{ab}^X([0,T]).
    $$

    Applying the discrete-time upcrossing inequality, we obtain
    $$
        \mathbb{E}\left(U_{ab}^X\left([0,T]\cap D_{n}\right)\right) \leq \frac{1}{b-a}\mathbb{E}(X_{T_n} - a)^+,
    $$
    where $T_n = \frac{[T \cdot 2^n]}{2^n}$. Letting $n\to \infty$, we get
    $$
        \mathbb{E}\left(U_{ab}^X([0,T])\right) \leq \frac{1}{b-a}\mathbb{E}(X_{T} - a)^+.
    $$  
}

\prop{
    If $(X_{t})_{t \geq 0}$ is a right-continuous submartingale and $\sup_{t \geq 0} \mathbb{E} X_{t}^+ < \infty$, then there exists $X \in L^1$ such that $X_t \xlongrightarrow{a.s.}X$ as $t \to \infty.$ 
}

\pf{
    $$
    \begin{aligned}
        \mathbb{E}\left(U_{ab}^X([0,\infty))\right) &= \lim_{n \to \infty} \mathbb{E}\left(U_{ab}^X([0,n])\right)\\
        &\leq \limsup_{n\to \infty}\frac{1}{b-a}\mathbb{E}(X_{n} - a)^+ \\
        &\leq \limsup_{n\to \infty}\frac{1}{b-a}\left(\mathbb{E}(X_n^+)+ \abs{a}\right) < \infty.
    \end{aligned}
    $$
    Hence for every $a < b$, we have $U_{ab}^X[0,\infty)< \infty$ a.s. 
    
    Therefore, with probability one, $U_{ab}^X[0,\infty) < \infty$ for all rational $a < b$.

    Hence the event
    $$
        \left\{\liminf_{t\to \infty}X_{t} \leq a \leq b < \limsup_{t \to \infty} X_t\right\}
    $$
    cannot occur.

    Therefore,
    $$
        \liminf_{t\to \infty} X_{t} = \limsup_{t \to \infty}X_t
    $$
    almost surely.
}

\defnn{
    A family of random variables $(X_t)_{t \geq 0}$ is uniformly integrable (U.I.) if 
    $$
        \lim_{M \to \infty}\sup_{t \geq 0} \mathbb{E}\left(\abs{X_t} \mathsf{1}_{\abs{X _t} > M}\right) = 0.
    $$
}

\rmkb{If $\sup_{t \geq 0} \mathbb{E} \abs{X_t}^p < \infty$ for some $p > 1$, then U.I.}

\thmnn{
    If $X_n \xlongrightarrow{P} X$, then the following are equivalent:
    $$
    \begin{cases}
        X_n \xlongrightarrow{L^1} X, \\
        \mathbb{E}\abs{X_n} \to \mathbb{E}\abs{X},\\
        \{X_n\} \text{ is U.I.}.
    \end{cases}
    $$
}

\prop{Let $Z \in L^1$. Then
$$
    \{\mathbb{E}(Z\mid \mathcal{G}) \colon \mathcal{G}\text{ is a sub-}\sigma \text{-field of }\mathcal{F}\}
$$
is uniformly integrable.

In particular, $\{\mathbb{E}(Z \mid \mathcal{F}_t)\}_{t \geq 0}$ is uniformly integrable, and it is a martingale by the tower property.
}

\prop{
    A right-continuous martingale $(X_t)_{t \geq 0}$ is uniformly integrable if and only if there exists $X _{\infty} \in L^1$ such that
    $$
        X _t = \mathbb{E}(X _\infty \mid \mathcal{F}_t).
    $$
}

\pf{
    ``$\Rightarrow$'' Uniform integrability implies $\sup_{t \geq 0}\mathbb{E}\abs{X_t} < \infty$.

    Hence the martingale convergence theorem yields some $X_{\infty}$ such that $X_t \xlongrightarrow{a.s.} X_\infty$ and
    $$
        \mathbb{E}\abs{X _{\infty}} = \mathbb{E}\left(\lim _{t \to \infty} \abs{X _t}\right) \leq \liminf _{t \to \infty} \mathbb{E}\abs{X _t} < \infty.
    $$

    To prove that $X_t = \mathbb{E}(X _{\infty} \mid \mathcal{F}_t)$ for every $t \geq 0$, it is enough to show that
    $$
        \mathbb{E}(X _t \mathsf{1}_A)=\mathbb{E}(X _{\infty}\mathsf{1}_A),
        \qquad \forall~A \in \mathcal{F}_t.
    $$

    For any $s \geq 0$, the martingale property gives
    $$
        \mathbb{E}(X_{t+s} \mid \mathcal{F}_{t}) = X_{t},
    $$
    and hence
    $$
        \mathbb{E}(X_{t+s}\mathsf{1}_{A}) = \mathbb{E}(X_{t}\mathsf{1}_{A}).
    $$
    On the other hand,
    $$
        \abs{\mathbb{E}(X_{t+s}\mathsf{1}_A) - \mathbb{E}(X_{\infty} \mathsf{1}_A)}
        \leq \mathbb{E}\abs{X_{t+s}-X_{\infty}} \to 0.
    $$

    Letting $s\to \infty$, we obtain
    $$
        \mathbb{E}(X_{t} \mathsf{1}_{A}) = \mathbb{E}(X_{\infty} \mathsf{1}_{A}).
    $$

    ``$\Leftarrow$'' If there exists $X_\infty \in L^1$ such that $X_t = \mathbb{E}(X_\infty \mid \mathcal{F}_t)$, then $\{X_t\}_{t \geq 0}$ is uniformly integrable by the previous proposition.
}

\thm{Optional sampling Theorem}{
    Let $(X_t)_{t \geq 0}$ be a right-continuous martingale. Let $S \leq T$ be two stopping times. Suppose that either
    \begin{enumerate}[label=\textcircled{\arabic*}]
        \item $\{X_t\}$ is uniformly integrable;
        \item there exists a constant $K >0$ such that $S,T \leq K$.
    \end{enumerate}
    Then
    $$
        \mathbb{E}(X_T \mid \mathcal{F}_S) = X _S.
    $$
    In particular,
    $$
        \mathbb{E}(X _T) = \mathbb{E}(X _S) = \mathbb{E}(X _0).
    $$
}

\pf{
    In case \textcircled{1}, let $Z = X_\infty$. In case \textcircled{2}, let $Z=X_{K}$.

    In either case, we have $X_t = \mathbb{E}(Z\mid \mathcal{F}_t)$: for all $t \geq 0$ in case \textcircled{1}, and for all $t \leq K$ in case \textcircled{2}.

    It therefore suffices to prove that $X_T = \mathbb{E}(Z \mid \mathcal{F}_T)$. Then similarly $X _S = \mathbb{E}(Z \mid \mathcal{F}_S)$.

    Hence 
    $$
    \begin{aligned}
    \mathbb{E}(X_T \mid \mathcal{F}_S) &= \mathbb{E}\left(\mathbb{E}(Z \mid \mathcal{F}_T) \mid \mathcal{F}_S\right)\\
    &= \mathbb{E}(Z \mid \mathcal{F}_S)\\
    &= X_S.
    \end{aligned}    
    $$

    First suppose that $T$ takes values in some countable set $\{t_1,t_2,\dots\}$. Then for every $A \in \mathcal{F}_T$, 
    $$
    \begin{aligned}
    \mathbb{E}(X_T \mathsf{1}_A)&= \sum_{n=1}^{\infty}\mathbb{E}(X_{t_n}\mathsf{1}_{\{T=t_n\}} \mathsf{1}_A)\\
    &= \sum_{n=1}^{\infty}\mathbb{E}\left(X_{t_n}\mathsf{1}_{A \cap \{T=t_n\}}\right)\\
    &= \sum_{n=1}^{\infty}\mathbb{E}\left(\mathbb{E}(Z\mid \mathcal{F}_{t_n})\mathsf{1}_{A \cap \{T=t_n\}}\right)\\
    &= \sum_{n=1}^{\infty}\mathbb{E}\left(Z\mathsf{1}_{A \cap \{T = t_n\}}\right)\\
    &= \mathbb{E}(Z \mathsf{1}_A).
    \end{aligned}
    $$
    because $A \cap \{T = t_n\}\in \mathcal{F}_{t_n}$ and $\mathbb{E}(Z \mid \mathcal{F}_{t_n}) = X_{t_n}$.

    Now consider a general stopping time $T$. Define $T_k = \frac{[2^k T] + 1}{2^k}$.

    For every $A \in \mathcal{F}_T \subseteq \mathcal{F}_{T _k}$, the first step gives
    $$
        \mathbb{E}(Z \mathsf{1}_A) = \mathbb{E}(X_{T_k} \mathsf{1}_A).
    $$

    Letting $k \to \infty$, we conclude that
    $$
        \mathbb{E}(Z \mathsf{1}_A) = \mathbb{E}(X_T \mathsf{1}_A).
    $$

    Indeed, $X_{T_k} \xlongrightarrow{a.s.} X_T$ by right-continuity, and $\{X_{T_k}\}$ is uniformly integrable because $X_{T_k} = \mathbb{E}(Z \mid \mathcal{F}_{T_k})$. Hence $X_{T_k} \xlongrightarrow{L^1}X_T$.
}

\thm{Doob's Maximal Inequality}{
    Let $(X_t)_{t \geq 0}$ be a right-continuous submartingale. Then for every $\lambda > 0$,
    $$
        \begin{aligned}
            &\mathsf{P}(\sup_{0 \leq s \leq t} X_s > \lambda) \leq \frac{1}{\lambda}\mathbb{E}X_t^+,\\
            &\mathsf{P}(\inf_{0 \leq s \leq t} X_s < -\lambda) \leq \frac{1}{\lambda}(\mathbb{E}X_t^+ - \mathbb{E}X_0).
        \end{aligned}
    $$
}

\pf{
    For $\lambda > 0$, let
    $$
        T = \inf \{t \geq 0: X_t > \lambda\}.
    $$
    Then
    $$
        \mathsf{P}(\sup_{0 \leq s \leq t}X_s > \lambda) = \mathsf{P}(T \leq t).
    $$
    
    Consider the stopped process $\{X_{t \wedge T}, t \geq 0\}$. Then
    $$
        \begin{aligned}
        \lambda \mathsf{P}(T \leq t) = \lambda \mathbb{E}(\mathsf{1}_{T \leq t}) &\leq \mathbb{E}(X_T \mathsf{1}_{T \leq t})\\
        & = \mathbb{E}(X_{T \wedge t} \mathsf{1}_{T \leq t}) \\
        &= \mathbb{E}(X_{T \wedge t}^+ \mathsf{1}_{T \leq t})\\
        &= \mathbb{E}(X_{T \wedge t}^+) \leq \mathbb{E}(X_t^+).
        \end{aligned}
    $$

    Next let
    $$
        S = \inf\{t \geq 0 : X_t < -\lambda\}.
    $$
    Then
    $$
        \mathsf{P}(\inf_{0 \leq s \leq t}X_s < -\lambda) = \mathsf{P}(S \leq t).
    $$
    Therefore
    $$
        \lambda \mathsf{P}(S \leq t) = \lambda \mathbb{E}(\mathsf{1}_{S \leq t}).
    $$
    
    Because $X_S \leq - \lambda$, $\lambda \leq - X_S$, so
    $$
        \lambda \mathsf{P}(S \leq t) \leq \mathbb{E}(-X_S \mathsf{1}_{S \leq t}) = - \mathbb{E}(X_{S \wedge t} \mathsf{1}_{S \leq t}).
    $$

    Applying optional stopping to the submartingale, we obtain
    $$
    \begin{aligned}
    \mathbb{E}X_0 \leq \mathbb{E}X_{t \wedge S}&= \mathbb{E}(X_{t \wedge S}\mathsf{1}_{S \leq t}) + \mathbb{E}(X_{t \wedge S}\mathsf{1}_{S > t})\\
    &\leq - \lambda \mathsf{P}(S \leq t) + \mathbb{E}(X_{t \wedge S}^+) \\
    & \leq -\lambda \mathsf{P}(S \leq t)+ \mathbb{E}(X_t^+).
    \end{aligned}
    $$
}

\thmnn{$\forall~p > 1$, $\displaystyle\mathbb{E}\left(\sup_{s \leq t}\abs{X_s} \right)^p \leq \left(\frac{p}{p-1}\right)^p \mathbb{E}\abs{X_t}^p.$}

\pf{
    Let $Y= \sup_{0 \leq s \leq t}\abs{X_s}.$
    $$
    \begin{aligned}
    \lambda \mathsf{P}(Y \geq \lambda) &\leq \mathbb{E}(\abs{X_{T \wedge t}}\mathsf{1}_{T \leq t}) \\
    & = \mathbb{E}\left( \abs{X_{T \wedge t}} \right) - \mathbb{E}\left( \abs{X_{T \wedge t}} \mathsf{1}_{T > t} \right)\\
    & \leq \mathbb{E}\left( \abs{X_t} \right) - \mathbb{E}\left( \abs{X_t}\mathsf{1}_{T >t} \right)\\
    &= \mathbb{E} \abs{X_t} - \mathbb{E}\left( \abs{X_t} \mathsf{1}_{Y < \lambda} \right)\\
    &= \mathbb{E}\left( \abs{X_t}\mathsf{1}_{Y \geq \lambda} \right).
    \end{aligned}
    $$
    Hence $\mathsf{P}(Y \geq \lambda) \leq \frac{1}{\lambda}\mathbb{E}\left( \abs{X_t}\mathsf{1}_{Y \geq \lambda} \right).$
    $$
    \begin{aligned}
        \mathbb{E}Y^p &= \int_{0}^{\infty} p \lambda^{p-1}\mathsf{P}(Y \geq \lambda) ~d \lambda \\
        &\leq \int_{0}^{\infty} p \lambda^{p-2}\mathbb{E}\left( \abs{X_t}\mathsf{1}_{Y \geq \lambda} \right)~d \lambda\\
        &=\mathbb{E}\left( \abs{X_t} \int_{0}^{\infty} p \lambda^{p-2}\mathsf{1}_{Y \geq \lambda} ~d\lambda \right)\\
        &=\mathbb{E}\left( \abs{X_t}\int_{0}^{Y}p \lambda^{p-2}~d\lambda \right)\\
        &= \mathbb{E}\left(\abs{X_t}\frac{p}{p-1}Y^{p-1}\right)\\
        &\leq \frac{p}{p-1}\left( \mathbb{E}\abs{X_t}^p \right)^{\frac{1}{p}}\left( \mathbb{E}(Y^{p-1})^q \right)^{\frac{1}{q}}\\
        &=\frac{p}{p-1}\left( \mathbb{E}\abs{X_t}^p \right)^{\frac{1}{p}}\left( \mathbb{E}Y^p \right)^{\frac{1}{q}}.
    \end{aligned}
    $$
    Therefore
    $$
        \left( \mathbb{E}Y^p \right)^{\frac{1}{p}} \leq \frac{p}{p-1}\left( \mathbb{E}\abs{X_t}^p \right)^{\frac{1}{p}}.
    $$

    To justify the argument rigorously, consider $Y \wedge k$ for $k > 0$. Then
    $$
    \begin{aligned}
        \mathbb{E}\left( Y \wedge k \right)^p &= \int_{0}^{\infty} p \lambda^{p-1}\mathsf{P}(Y\wedge k \geq \lambda)~d \lambda \\
        &= \int_{0}^{k}p \lambda^{p-1}\mathsf{P}(Y \wedge k \geq \lambda)~d \lambda + \underbrace{\int_{k}^{\infty} p \lambda^{p-1}\mathsf{P}(Y\wedge k \geq \lambda)~d \lambda}_{=0} \\
        &= \int_{0}^{k}p \lambda^{p-1}\mathsf{P}(Y \wedge k \geq \lambda)~d\lambda\\
        &= \int_{0}^{k}p \lambda^{p-1}\mathsf{P}(Y \geq \lambda) ~d \lambda\\
        &\leq \mathbb{E}\left( \abs{X_t}
        \int_{0}^{k} p \lambda^{p-1} \mathsf{1}_{Y \geq \lambda}~d \lambda \right)\\
        &= \mathbb{E}\left( \abs{X_t}\int_{0}^{k \wedge Y} p \lambda^{p-2}~d \lambda \right)\\
        &= \mathbb{E}\left[ 
            \abs{X_t} \frac{p}{p-1}(Y \wedge k )^{p-1}
         \right].
    \end{aligned}
    $$

    $\implies \mathbb{E}(Y \wedge k)^p \leq \frac{p}{p-1}\left( \mathbb{E}\abs{X_t}^p \right)^{\frac{1}{p}}\left( 
        \mathbb{E}(Y \wedge k)^p
     \right)^{\frac{1}{q}}.$

    $\implies \left( \mathbb{E}(Y \wedge k)^p \right)^{\frac{1}{p}} \leq \frac{p}{p-1}\left( \mathbb{E}\abs{X_t}^p \right)^{\frac{1}{p}}$.

    Letting $k\to \infty$, we obtain
    $$
        \left( \mathbb{E} Y^p \right)^{\frac{1}{p}} \leq \frac{p}{p-1}\left( \mathbb{E}\abs{X_t}^p \right)^{\frac{1}{p}}.
    $$
}

\subsection{The Doob-Meyer Decomposition}
We now turn to the continuous-time analogue of the discrete Doob decomposition.

Recall the discrete-time Doob decomposition
$$
    X_n = M_n + A_n,
$$
where $X_n$ is a submartingale, $M_n$ is a martingale, and $A_n$ is a predictable increasing process.

$$
    \begin{aligned}
        &\mathbb{E}\left( M_{n+1} \mid \mathcal{F}_n \right) = M_n \\
        \implies& \mathbb{E}(X_{n+1}- A_{n+1}\mid \mathcal{F}_n) = X_n - A_n \\
        \implies& \mathbb{E}(X_{n+1}\mid \mathcal{F}_n) - A_{n+1} = X_n - A_n \\
        \implies& A_{n+1} - A_n = \mathbb{E}(X_{n+1}\mid \mathcal{F}_n) - X_ n \geq 0.
    \end{aligned}
$$

$$
\begin{cases}
A_n = \sum_{k=0}^{n-1}\left( \mathbb{E}(X_{k+1} \mid \mathcal{F}_k) - X_k \right)\\
A_0 = 0.
\end{cases}
$$

For a right-continuous submartingale $(X_t)_{t \geq 0}$, we seek a decomposition of the form
$$
    X_t = M_t + A_t,
$$
where $M_t$ is a martingale and $A_t$ is an increasing and ``predictable'' process.

So we first clarify what properties the increasing part should satisfy.

\defnn{
        An adapted process $A$ is called increasing if, almost surely,
    \begin{enumerate}[label=(\alph*)]
        \item $A_0 = 0$;
        \item $t \mapsto A_t$ is nondecreasing and right-continuous;
        \item $\mathbb{E}(A_t) < \infty$, $\forall~ t \geq 0 $. 
    \end{enumerate}
    We say that $A$ is integrable if
    $$
        \mathbb{E}(A_\infty) = \mathbb{E}(\lim_{t \to \infty}A_t) < \infty.
    $$
}

\defnn{
    An increasing process $A$ is called natural if for every bounded, right-continuous martingale $(M_t)_{t \geq 0}$ and every $t > 0$,
    $$
    \mathbb{E}\int_{0}^{t} M_s ~d A_s = \mathbb{E} \int_{0}^{t} M_{s-} ~d A_s.
    $$
}
\rmkb{If $s \mapsto A_s$ is a.s. continuous, then $A$ is natural.}

\defnn{
    An increasing sequence $(A_n,\mathcal{F}_n)$ is called natural if for every bounded martingale $\{M_n,\mathcal{F}_n\}$,
    $$
    \mathbb{E}(M_n A_n) = \mathbb{E}\left( \sum_{k=1}^n M_{k-1}(A_k - A_{k-1}) \right).
    $$
}

\prop{
    In discrete time, predictability is equivalent to naturality.
}

\pf{
    ``$\Rightarrow$'' Assume that $A$ is predictable, i.e. $A_n \in \mathcal{F}_{n-1}$. We show that $A$ is natural.

    Indeed,
    $$
    \begin{aligned}
        \mathbb{E}(M_n A_n)
        &= \mathbb{E}\left(\sum_{k=1}^n M_kA_k - \sum_{k=1}^n M_{k-1}A_{k-1}\right)\\
        &= \sum_{k=1}^n \mathbb{E}(M_kA_k) - \sum_{k=1}^n \mathbb{E}(M_{k-1}A_{k-1})\\
        \Longleftrightarrow& \sum_{k=2}^n \mathbb{E}(M_{k-1}A_{k-1}) + \mathbb{E}(M_n A_n) = \sum_{k=1}^n \mathbb{E}(M_{k-1}A_k)\\
        \Longleftrightarrow &
        \sum_{k=1}^n \mathbb{E}(M_k A_k) = \sum_{k=1}^n \mathbb{E}(M_{k-1}A_k)\\
        \Longleftrightarrow &
        \sum_{k=1}^n \mathbb{E}\left[ 
            (M_k - M_{k-1})A_{k}
         \right]
         = 0.
    \end{aligned}
    $$

    Since $A_n \in \mathcal{F}_{n-1}$,
    $$
    \mathbb{E}\left( 
        A_n (M_n - M_{n-1})
     \right) = \mathbb{E} \left[ 
        A_n \underbrace{\mathbb{E}(M_n - M_{n-1} \mid \mathcal{F}_{n-1} )}_{=0}
      \right] = 0.
    $$
    Therefore $A$ is natural.

    ``$\Leftarrow$'' Now assume that $A$ is natural. Then for every bounded martingale $(M_n)$,
    $$
    \mathbb{E}\left[ 
        A_n (M_n - M_{n-1})
     \right] = 0,\quad \forall~ n \geq 1.
    $$

    Our goal is to prove that
    $$
        A_n = \mathbb{E}(A_n \mid \mathcal{F}_{n-1}) \quad \text{a.s.},
    $$
    which is equivalent to $A_n \in \mathcal{F}_{n-1}$.
    $$
    \begin{aligned}
        \mathbb{E}\left[ 
            M_n\left( A_n - \mathbb{E}(A_n \mid \mathcal{F}_{n-1}) \right)
         \right]
         &= \mathbb{E}(M_nA_n) - \mathbb{E} \left[ 
            M_n \, \mathbb{E}(A_n \mid \mathcal{F}_{n-1})
          \right]\\
         &= \mathbb{E}(M_nA_n) - \mathbb{E}\left[
            \mathbb{E}(A_n \mid \mathcal{F}_{n-1}) \, \mathbb{E}(M_n \mid \mathcal{F}_{n-1})
          \right]\\
         &= \mathbb{E}(M_nA_n) - \mathbb{E}\left[
            \mathbb{E}(A_n \mid \mathcal{F}_{n-1}) M_{n-1}
          \right]\\
         &= \mathbb{E}(M_nA_n) - \mathbb{E}(A_nM_{n-1})\\
         &= \mathbb{E}\left[ A_n(M_n - M_{n-1}) \right] = 0.
    \end{aligned}
    $$

    Fix $n \geq 1$, and let
    $$
        X = \mathrm{sgn} \left( A_n - \mathbb{E}(A_n \mid \mathcal{F}_{n-1}) \right)
    $$
    Then $X\left( A_n - \mathbb{E}(A_n\mid \mathcal{F}_{n-1}) \right) = \abs{A_n - \mathbb{E}(A_n \mid \mathcal{F}_{n-1})}$.

    Define
    $$
    M_k = \begin{cases}
       \mathbb{E}(X \mid \mathcal{F}_{k}),& 0 \leq k \leq n-1\\
       X, &k=n
    \end{cases}
    $$

    Then $(M_k)$ is a bounded martingale.

    $$
    \begin{aligned}
        0 &= \mathbb{E}\left[ M_n (A_n - \mathbb{E}(A_n \mid \mathcal{F}_{n-1})) \right]\\
        &= \mathbb{E} \abs{A_n - \mathbb{E}(A_n \mid \mathcal{F}_{n-1})}
    \end{aligned}
    $$
    Therefore $A_n = \mathbb{E}(A_n \mid \mathcal{F}_{n-1})$ a.s.
    
    Hence $A_n \in \mathcal{F}_{n-1}$, so $A$ is predictable.
}

The previous proposition gives the discrete-time characterization of natural increasing processes.
We now explain why the continuous-time definition is the correct analogue.

$$
    \begin{aligned}
        &\mathbb{E}(M_n A_n) = \sum_{k=1}^n \mathbb{E}\left( M_{k-1} (A_k - A_{k-1}) \right) \\
        \Longleftrightarrow &
        \mathbb{E}(M_tA_t) = \mathbb{E}\left( \int_{0}^{t}M_{s-}~dA_s \right) = \mathbb{E}\left( \int_{0}^{t} M_s ~dA_s \right).
    \end{aligned}
$$

Indeed, the discrete sum should be viewed as a Riemann--Stieltjes approximation of the corresponding continuous-time integral.

$$
    \int_{0}^{t} M_s ~d A_s = \lim_{\norm{\Pi}\to 0} \int_{0}^{t} M_s^\Pi ~d A_s.
$$
where $\Pi = \{t_0, t_1,\cdots,t_n\}\subseteq[0,t]$ with $0 = t_0 < t_1 < \cdots < t_n = t$ and $\norm{\Pi} = \max_{1 \leq k \leq n}(t_k - t_{k-1})$. By right-continuity,
$$
    M_s^\Pi = \sum_{k=1}^{n} M_{t_k}\mathsf{1}_{(t_{k-1},t_k]}(s) \xlongrightarrow{a.s.} M_s.
$$
For each $\Pi$,
$$
\begin{aligned}
\mathbb{E}\int_{0}^{t}M_s ^\Pi ~d A_s
&= \mathbb{E} \sum_{k=1}^n M_{t_k}(A_{t_k} - A_{t_{k-1}}) \\
&= \sum_{k=1}^n \mathbb{E}(M_{t_k}A_{t_k}) - \sum_{k=1}^n \mathbb{E}(M_{t_k}A_{t_{k-1}})\\
&= \mathbb{E}(M_{t_n}A_{t_n}) + \sum_{k=1}^{n-1}\mathbb{E}(M_{t_{k}} A_{t_{k}}) - \sum_{k=2}^{n}\mathbb{E}(M_{t_k} A_{t_{k-1}})\\
&= \mathbb{E}(M_t A_t) + \sum_{k=1}^{n-1}\mathbb{E}(M_{t_k} A_{t_k}) - \sum_{k=1}^{n-1}\mathbb{E}(M_{t_{k+1}}A_{t_k}) \\
&= \mathbb{E}(M_t A_t) - \underbrace{\sum_{k=1}^{n-1}\mathbb{E}\left( A_{t_k} (M_{t_{k+1}} - M_{t_k}) \right)}_{=0}.
\end{aligned}
$$

Thus the continuous-time identity is exactly the limit of the discrete-time naturality relation.

\defnn{
    Let
    $$
        S = \{\text{stopping time }T \text{ with }T < \infty \text{ a.s.}\}.
    $$
    We say that a right-continuous process $\{X_t\}_{t \geq 0}$ is of class $D$ if $\{X_T\}_{T \in S}$ is uniformly integrable. In particular, this implies that $\{X_t\}_{t \geq 0}$ itself is uniformly integrable.

    Fix $a > 0$, and let
    $$
        S_a = \{\text{stopping time } T \text{ with }T \leq a \text{ a.s.}\}.
    $$
    We say that $\{X_t\}_{t \geq 0}$ is of class $DL$ if $\{X_T\}_{T \in S_a}$ is uniformly integrable for every $a > 0$.
}


\thm{Doob-Meyer Decomposition}{
    If $\{X_t\}$ is of class $DL$, then $X_t = M_t + A_t$, where $\{M_t\}$ is a martingale and $\{A_t\}$ is increasing.

    If $(A_t)$ is natural, then $(M_t, A_t)$ is unique.
}

\pf{
    Assume that
    $$
        X_t = M_t + A_t = M_t' + A_t',
    $$
    where $(M_t)$ and $(M_t')$ are martingales, and $(A_t)$ and $(A_t')$ are natural increasing processes.

    Let
    $$
        B_t = M_t - M_t' = A_t' - A_t.
    $$
    Then $(B_t)$ is a martingale, and also a process of bounded variation. Moreover, it is natural.
    
    To prove uniqueness, let $(\xi_t)_{t \geq 0}$ be any bounded right-continuous martingale. Then
    $$
    \begin{aligned}
        \mathbb{E}\bigl[\xi_t (A_t - A_t')\bigr]
        &= \mathbb{E}\int_{0}^{t}\xi_{s-} (d A_s - d A_s') \\
        &= \mathbb{E}\left[ \int_{0}^{t}\xi_{s-}  dB_s \right]\\
        &= \lim_{n \to \infty} \mathbb{E}\sum_{j=1}^{m_n}\xi_{t_{j-1}^{(n)}} \left[ B_{t_j^{(n)}} - B_{t_{j-1}^{(n)}} \right],
    \end{aligned}
    $$
    where $\Pi_n = \{t_0^{(n)}, \cdots, t_{m_n}^{(n)}\}$ is a partition of $[0,t]$.

    For each $j$,
    $$
    \mathbb{E}\left[ 
        \xi_{t_{j-1}^{(n)}} (B_{t_j^{(n)}} - B_{t_{j-1}^{(n)}})
     \right] = \mathbb{E} \left[ 
        \xi_{t_{j-1}^{(n)}} \underbrace{\mathbb{E}(B_{t_j^{(n)}} - B_{t_{j-1}^{(n)}} \mid \mathcal{F}_{t_{j-1}^{(n)}})}_{=0}
      \right].
    $$
    Hence $\mathbb{E}\bigl[\xi_t (A_t - A_t')\bigr] = 0.$

    Now let $\xi$ be any bounded $\mathcal{F}_t$-measurable random variable. Applying the preceding identity with the martingale
    $$
        \xi_s = \mathbb{E}(\xi \mid \mathcal{F}_s), \qquad 0 \leq s \leq t,
    $$
    we get
    $$
        \mathbb{E}\left[ \xi (A_t - A_t') \right] = 0.
    $$
    
    Taking $\xi = \mathrm{sgn}(A_t - A_t')$, we obtain
    $$
        \mathbb{E}\abs{A_t - A_t'} = 0.
    $$
    Hence $A_t = A_t'$ a.s. for every $t > 0$. In particular, this holds for all $q \in \mathbb{Q}_+$. By right-continuity, we conclude that $A_t = A_t'$ for all $t \geq 0$ a.s.

    \textbf{Existence: } Fix $a > 0$.
    
    It suffices to construct the decomposition on $[0,a]$.

    Consider
    $$
        Y_t \triangleq X_t - \mathbb{E}(X_a \mid \mathcal{F}_t), \qquad 0 \leq t \leq a.
    $$
    Then $(Y_t)_{0 \leq t \leq a}$ is a submartingale, because
    $$
    \begin{aligned}
        \mathbb{E}(Y_t \mid \mathcal{F}_s)
        &= \mathbb{E}(X_t \mid \mathcal{F}_s) - \mathbb{E}(X_a \mid \mathcal{F}_s) \\
        &\geq X_s - \mathbb{E}(X_a \mid \mathcal{F}_s)
        = Y_s, \qquad 0 \leq s \leq t \leq a.
    \end{aligned}
    $$

    Let $\Pi_n = \{t_0^{(n)},\cdots, t_{2^n}^{(n)}\}$ with $t_j^{(n)} = \frac{j}{2^n}a$, $0 \leq j \leq 2^n$. For each $n$, the discrete-time Doob decomposition gives
    $$
        Y_{t_{j}^{(n)}} = M_{t_{j}^{(n)}}^{(n)} + A_{t_j^{(n)}}^{(n)}, \qquad 0 \leq j \leq 2^n,
    $$
    where
    $$
        A_{t_j^{(n)}}^{(n)} = \sum_{k=0}^{j-1}\mathbb{E}\left( Y_{t_{k+1}^{(n)}} - Y_{t_{k}^{(n)}} \middle| \mathcal{F}_{t_{k}^{(n)}}\right).
    $$
    Since
    $$
    \begin{cases}
        Y_0 = X_0 - \mathbb{E}(X_a \mid \mathcal{F}_0) \leq 0,\\
        Y_a = X_a - \mathbb{E}(X_a \mid \mathcal{F}_a) = 0,
    \end{cases}
    $$
    we have
    $$
        0 = Y_a = M_a^{(n)} + A_a^{(n)}.
    $$
    Hence
    $$
        M_{t_{j}^{(n)}}^{(n)} = \mathbb{E}(M_a^{(n)} \mid \mathcal{F}_{t_{j}^{(n)}}) = - \mathbb{E}(A_a^{(n)} \mid \mathcal{F}_{t_{j}^{(n)}}),
    $$
    and therefore
    $$
        Y_{t_{j}^{(n)}} = A_{t_{j}^{(n)}}^{(n)} - \mathbb{E}(A_a^{(n)} \mid \mathcal{F}_{t_{j}^{(n)}}).
    $$

    We will show that $\{A_a^{(n)}\}_{n \geq 1}$ is uniformly integrable. Then, by the Dunford-Pettis compactness criterion, there exists an integrable random variable $A_a$ such that
    $$
        A_a^{(n_k)} \rightharpoonup A_a
        \qquad \text{weakly in } L^1
    $$
    along some subsequence. Equivalently, for every bounded random variable $\xi$,
    $$
        \lim_{k \to \infty} \mathbb{E}(\xi A_a^{(n_k)}) = \mathbb{E}(\xi A_a).
    $$

    Define
    $$
        A_t = Y_t + \mathbb{E}(A_a \mid \mathcal{F}_t), \qquad 0 \leq t \leq a.
    $$

    By Problem 4.11, the convergence $\mathbb{E}(\xi A_a^{(n_k)}) \to \mathbb{E}(\xi A_a)$ for every bounded random variable $\xi$ implies
    $$
        \mathbb{E}(\xi \mathbb{E}(A_a^{(n_k)} \mid \mathcal{F}_t)) \to \mathbb{E}(\xi \mathbb{E}(A_a \mid \mathcal{F}_t)),
    $$
    that is, $\mathbb{E}(A_a^{(n_k)} \mid \mathcal{F}_t)$ converges weakly in $L^1$ to $\mathbb{E}(A_a \mid \mathcal{F}_t)$.

    Fix a dyadic time $t$. Then, for all sufficiently large $k$, we have
    $$
    \begin{cases}
        A_t = Y_t + \mathbb{E}(A_a \mid \mathcal{F}_t), \\
        A_t^{(n_k)} = Y_t + \mathbb{E}(A_a^{(n_k)} \mid \mathcal{F}_t)
    \end{cases}
    $$
    because $t$ belongs to the dyadic partition $\Pi_{n_k}$ once $k$ is large enough.
    Therefore
    $$
        A_t^{(n_k)} - A_t = \mathbb{E}(A_a^{(n_k)} \mid \mathcal{F}_t) - \mathbb{E}(A_a \mid \mathcal{F}_t).
    $$
    Since $\mathbb{E}(A_a^{(n_k)} \mid \mathcal{F}_t)$ converges weakly in $L^1$ to $\mathbb{E}(A_a \mid \mathcal{F}_t)$, it follows that
    $$
        \mathbb{E}(\xi A_t^{(n_k)}) \to \mathbb{E}(\xi A_t)
    $$
    for every bounded random variable $\xi$. 

    For any dyadic times $0 \leq s < t \leq a$,
    $$
        \mathbb{E}(\xi (A_t - A_s)) = \lim_{k \to \infty}\mathbb{E}(\xi (A_t^{(n_k)} - A_s^{(n_k)})) \geq 0
    $$
    whenever $\xi \geq 0$.

    Taking $\xi = \mathsf{1}_{A_s > A_t}$, we get
    $$
        \mathbb{E}\left[ (A_t - A_s) \mathsf{1}_{A_s > A_t} \right] \geq 0.
    $$
    Therefore $A_t - A_s \geq 0$ a.s.

    Since
    $$
        A_0^{(n)} = 0 = Y_0 + \mathbb{E}(A_a^{(n)} \mid \mathcal{F}_0),
    $$
    passing to the limit gives
    $$
        A_0 = Y_0 + \mathbb{E}(A_a \mid \mathcal{F}_0) = 0.
    $$

    Moreover, for any dyadic time $t$,
    $$
        X_t - \mathbb{E}(X_a \mid \mathcal{F}_t) = Y_t = A_t^{(n_k)} - \mathbb{E}(A_a^{(n_k)} \mid \mathcal{F}_t)
    $$
    for all sufficiently large $k$.

    Passing to the limit along the subsequence gives
    $$
        Y_t = A_t - \mathbb{E}(A_a \mid \mathcal{F}_t) = X_t - \mathbb{E}(X_a \mid \mathcal{F}_t).
    $$
    Therefore
    $$
        X_t = A_t + \mathbb{E}(X_a - A_a \mid \mathcal{F}_t)
    $$
    for every dyadic time $t \in [0,a]$. By right-continuity, the same identity then holds for all $t \in [0,a]$.

    Define
    $$
        M_t = \mathbb{E}(X_a - A_a \mid \mathcal{F}_t), \qquad 0 \leq t \leq a.
    $$
    Then $(M_t)_{0 \leq t \leq a}$ is a martingale and
    $$
        X_t = A_t + M_t.
    $$

    Finally, we prove that $(A_t)$ is natural. It is enough to show that for every bounded right-continuous martingale $(\xi_t)$,
    $$
        \mathbb{E}\int_{0}^{a} \xi_s ~dA_s = \mathbb{E}\int_{0}^{a}\xi_{s-}~dA_s.
    $$
    We first note that
    $$
        \mathbb{E}\int_{0}^{a} \xi_s ~dA_s = \mathbb{E}(\xi_a A_a) = \lim_{n \to \infty}\mathbb{E}(\xi_a A_a^{(n)}).
    $$

    On the other hand,
    $$
        \mathbb{E}\int_{0}^{a}\xi_{s-}~dA_s
        = \lim_{n \to \infty} \mathbb{E} \sum_{j=1}^{2^n}\xi_{t_{j-1}^{(n)}}(A_{t_{j}^{(n)}} - A_{t_{j-1}^{(n)}}).
    $$

    Since $\{A_{t_{j}^{(n)}}^{(n)}\}_{0 \leq j \leq 2^n}$ is predictable, and $t \mapsto \mathbb{E}(A_a \mid \mathcal{F}_t)$ is a martingale,
    $$
    \begin{aligned}
        \mathbb{E}(\xi_a A_a^{(n)})
        &= \mathbb{E}\left( \sum_{j=1}^{2^n}\xi_{t_{j-1}^{(n)}} (A_{t_{j}^{(n)}}^{(n)} - A_{t_{j-1}^{(n)}}^{(n)}) \right)\\
        &= \mathbb{E}\left( \sum_{j=1}^{2^n}\xi_{t_{j-1}^{(n)}}(Y_{t_{j}^{(n)}}-Y_{t_{j-1}^{(n)}}) \right)\\
        &= \mathbb{E}\left( \sum_{j=1}^{2^n}\xi_{t_{j-1}^{(n)}}\Bigl[(A_{t_{j}^{(n)}}-A_{t_{j-1}^{(n)}}) - \bigl(\mathbb{E}(A_a \mid \mathcal{F}_{t_j^{(n)}})-\mathbb{E}(A_a \mid \mathcal{F}_{t_{j-1}^{(n)}})\bigr)\Bigr] \right)\\
        &= \mathbb{E}\left( \sum_{j=1}^{2^n}\xi_{t_{j-1}^{(n)}}(A_{t_{j}^{(n)}}-A_{t_{j-1}^{(n)}}) \right).
    \end{aligned}
    $$
    Here the last equality follows because $t \mapsto \mathbb{E}(A_a \mid \mathcal{F}_t)$ is a martingale.
    Letting $n \to \infty$, we obtain
    $$
        \mathbb{E}\int_{0}^{a} \xi_s ~dA_s = \mathbb{E}\int_{0}^{a}\xi_{s-}~dA_s.
    $$
    Therefore $(A_t)$ is natural.
}


\section{Square Integrable}

\defnn{
    Let $(X_t)_{t \geq 0}$ be a right-continuous martingale. We say that $X$ is square-integrable if
    $$
        \mathbb{E}X_t^2 < \infty, \qquad \forall~ t \geq 0.
    $$
    Throughout this section, we assume $X_0 = 0$ and write $X \in \mathcal{M}_2$.
}

Then $Y_t = X_t^2$ is a submartingale. Recall that a process is of class $DL$ if, for every $a > 0$, the family $\{Y_T : T \leq a \text{ is a stopping time}\}$ is uniformly integrable. By the Doob--Meyer decomposition,
$$
    Y_t = X_t^2 = M_t + A_t, \qquad t \geq 0.
$$

If $X \in \mathcal{M}_2^c$ is continuous, then we may choose $M_t$ and $A_t$ to be continuous as well.

\defnn{
    For $X \in \mathcal{M}_2$, the quadratic variation of $X$ is defined by $\langle X \rangle_t \triangleq A_t$, that is, $\langle X \rangle_t$ is the unique adapted natural increasing process such that $\langle X \rangle_0 = 0$ and
    $$
        X_t^2 - \langle X \rangle_t
    $$
    is a martingale.
}

\rmkb{
    The notation $\langle X \rangle_t$ is introduced here through the Doob--Meyer decomposition of the submartingale $X_t^2$.
}

If $X, Y \in \mathcal{M}_2$, then
$$
\begin{cases}
    (X_t + Y_t)^2 - \langle X+Y \rangle_t \text{ is a martingale},\\
    (X_t - Y_t)^2 - \langle X-Y \rangle_t \text{ is a martingale}.
\end{cases}
$$
Subtracting the second relation from the first, we see that
$$
    4X_tY_t - \bigl(\langle X+Y \rangle_t - \langle X-Y \rangle_t\bigr)
$$
is a martingale.

\defnn{
    For $X, Y \in \mathcal{M}_2$, we define the cross variation by
    $$
        \langle X,Y \rangle_t \triangleq \frac{1}{4}\left( \langle X+Y \rangle_t - \langle X-Y \rangle_t \right),
    $$
    so that
    $$
        X_tY_t - \langle X,Y \rangle_t
    $$
    is a martingale.
}

\rmkb{
    $\langle X,X \rangle_t = \langle X \rangle_t$.
}

\defnn{
    Let $\Pi = \{t_0,\cdots,t_m\}$ be a partition of $[0,t]$. The $p$-th variation of $X$ over $\Pi$ is defined by
    $$
        V_t^{(p)}(\Pi) = \sum_{k=1}^{m}\abs{ X_{t_k} - X_{t_{k-1}} }^p.
    $$
    When $p = 2$, we call $V_t^{(2)}(\Pi)$ the quadratic variation along the partition $\Pi$.
}

\thmnn{
    Let $X \in \mathcal{M}_2^c$. For any partition $\Pi$ of $[0,t]$, we have
    $$
        V_t^{(2)}(\Pi) \xlongrightarrow{P} \langle X \rangle_t
    $$
    as $\norm{\Pi} \to 0$.
}

\pf{
    Let $\Pi = \{t_0,\dots,t_m\}$ and write
    $$
        \Delta_k X = X_{t_k} - X_{t_{k-1}}, \qquad \Delta_k \langle X \rangle = \langle X \rangle_{t_k} - \langle X \rangle_{t_{k-1}}.
    $$
    Then
    $$
        V_t^{(2)}(\Pi) - \langle X \rangle_t
        = \sum_{k=1}^{m}\left( (\Delta_k X)^2 - \Delta_k \langle X \rangle \right).
    $$
    Hence
    $$
    \begin{aligned}
        &\mathbb{E}\left( V_t^{(2)}(\Pi) - \langle X \rangle_t \right)^2\\
        =&
        \mathbb{E}\left[
            \sum_{k=1}^{m}\left( (\Delta_k X)^2 - \Delta_k \langle X \rangle \right)
        \right]^2 \\
        =&
        \sum_{k=1}^{m}\mathbb{E}\left[ \left( (\Delta_k X)^2 - \Delta_k \langle X \rangle \right)^2 \right]
        \\
        &+ 2\sum_{1 \leq j < k \leq m}\mathbb{E}\left[
            \bigl((\Delta_k X)^2 - \Delta_k \langle X \rangle\bigr)
            \bigl((\Delta_j X)^2 - \Delta_j \langle X \rangle\bigr)
        \right].
    \end{aligned}
    $$

    The cross terms vanish. Indeed, for $j < k$, the random variable
    $$
        (\Delta_j X)^2 - \Delta_j \langle X \rangle
    $$
    is $\mathcal{F}_{t_{k-1}}$-measurable, while
    $$
        \mathbb{E}\left( (\Delta_k X)^2 - \Delta_k \langle X \rangle \middle| \mathcal{F}_{t_{k-1}} \right) = 0
    $$
    because $X_s^2 - \langle X \rangle_s$ is a martingale. Therefore
    $$
        \mathbb{E}\left[
            \bigl((\Delta_k X)^2 - \Delta_k \langle X \rangle\bigr)
            \bigl((\Delta_j X)^2 - \Delta_j \langle X \rangle\bigr)
        \right] = 0.
    $$

    Consequently,
    $$
    \begin{aligned}
        \mathbb{E}\left( V_t^{(2)}(\Pi) - \langle X \rangle_t \right)^2
        &=
        \sum_{k=1}^{m}\mathbb{E}\left[ \left( (\Delta_k X)^2 - \Delta_k \langle X \rangle \right)^2 \right] \\
        &\leq
        2\sum_{k=1}^{m}\mathbb{E}(\Delta_k X)^4
        + 2\sum_{k=1}^{m}\mathbb{E}(\Delta_k \langle X \rangle)^2 \\
        &= I + II\\
        &=2\mathbb{E}V_t^{(4)}(\Pi)+ II
    \end{aligned}
    $$

    We first consider the bounded case. Assume that
    $$
        \abs{X_s} \leq K, \qquad \langle X \rangle_s \leq K, \qquad 0 \leq s \leq t.
    $$

    \lemnn{
        Let $X \in \mathcal{M}_2$ satisfy $\abs{X_s} \leq K < \infty$ for all $0 \leq s \leq t$. Then
        $$
            \mathbb{E}\left[ V_t^{(2)}(\Pi)^2 \right] \leq 6K^4.
        $$
    }

    Assuming this lemma, we obtain
    $$
        V_t^{(4)}(\Pi)
        \leq
        V_t^{(2)}(\Pi)\, m_t^2(X;\norm{\Pi}),
    $$
    where
    $$
        m_t(X,\delta) = \sup\{\abs{X_r - X_s} : \abs{s-r} \leq \delta,\ 0 \leq r < s \leq t\}.
    $$
    Hence, by Cauchy--Schwarz,
    $$
    \begin{aligned}
        \mathbb{E}\left[ V_t^{(4)}(\Pi) \right]
        &\leq \mathbb{E}\left[ V_t^{(2)}(\Pi)\,m_t^2(X,\norm{\Pi}) \right] \\
        &\leq \left(\mathbb{E}\left[ V_t^{(2)}(\Pi)^2 \right]\right)^{1/2}
        \left(\mathbb{E}m_t^4(X,\norm{\Pi})\right)^{1/2} \\
        &\leq (6K^4)^{1/2}\left(\mathbb{E}m_t^4(X,\norm{\Pi})\right)^{1/2}.
    \end{aligned}
    $$
    Since $X$ is continuous and bounded on $[0,t]$, bounded convergence yields
    $$
        \mathbb{E}m_t^4(X,\norm{\Pi}) \to 0
    $$
    as $\norm{\Pi} \to 0$. Thus $I \to 0$.

    For $II$, note that $\langle X \rangle$ is continuous and increasing, so
    $$
        \sum_{k=1}^{m}(\Delta_k \langle X \rangle)^2
        \leq
        \sup_{1 \leq k \leq m}\abs{\Delta_k \langle X \rangle}\,\langle X \rangle_t.
    $$
    Since $\langle X \rangle$ is uniformly continuous on $[0,t]$ and bounded by $K$, it follows that
    $$
        \mathbb{E}\left[
            \sup_{1 \leq k \leq m}\abs{\Delta_k \langle X \rangle}\,\langle X \rangle_t
        \right] \to 0
    $$
    as $\norm{\Pi} \to 0$. Hence $II \to 0$ as well.

    Therefore, in the bounded case,
    $$
        \mathbb{E}\left[
            \left( V_t^{(2)}(\Pi) - \langle X \rangle_t \right)^2
        \right] \to 0
    $$
    as $\norm{\Pi} \to 0$, and hence
    $$
        V_t^{(2)}(\Pi) \xlongrightarrow{P} \langle X \rangle_t.
    $$

    For a general $X \in \mathcal{M}_2^c$, we use localization. For $n \geq 1$, define
    $$
        T_n = \inf \{ t \geq 0 : \abs{X_t} \geq n \text{ or } \langle X \rangle_t \geq n \}.
    $$
    Then $T_n \uparrow \infty$ a.s., and
    $$
        X_t^{(n)} = X_{t \wedge T_n}
    $$
    is a bounded martingale. Moreover,
    $$
        X_{t \wedge T_n}^2 - \langle X \rangle_{t \wedge T_n}
    $$
    is a martingale, so by uniqueness of the quadratic variation,
    $$
        \langle X^{(n)} \rangle_t = \langle X \rangle_{t \wedge T_n}.
    $$
    Hence
    $$
    \begin{cases}
        \abs{X_s^{(n)}} \leq n, \\
        \langle X^{(n)} \rangle_s \leq n,
    \end{cases}
    \qquad 0 \leq s \leq t,
    $$
    and the bounded case gives
    $$
        \mathbb{E}\left[
            \left(
                V_{t \wedge T_n}^{(2)}(X) - \langle X \rangle_{t \wedge T_n}
            \right)^2
        \right] \to 0
    $$
    as $\norm{\Pi} \to 0$.

    Now fix $\varepsilon > 0$ and $\eta > 0$. Choose $n$ so large that
    $$
        \mathsf{P}(T_n < t) \leq \frac{\eta}{2}.
    $$
    Then
    $$
    \begin{aligned}
        \mathsf{P}\left( \abs{V_t^{(2)}(X) - \langle X \rangle_t} > \varepsilon \right)
        &\leq
        \mathsf{P}(T_n < t)
        + \mathsf{P}\left(
            T_n \geq t,\,
            \abs{V_{t \wedge T_n}^{(2)}(X) - \langle X \rangle_{t \wedge T_n}} > \varepsilon
        \right) \\
        &\leq
        \frac{\eta}{2}
        + \frac{1}{\varepsilon^2}
        \mathbb{E}\left[
            \left(
                V_{t \wedge T_n}^{(2)}(X) - \langle X \rangle_{t \wedge T_n}
            \right)^2
        \right].
    \end{aligned}
    $$
    Taking $\norm{\Pi}$ small enough, the second term is at most $\eta/2$. Therefore
    $$
        \mathsf{P}\left( \abs{V_t^{(2)}(X) - \langle X \rangle_t} > \varepsilon \right) \leq \eta.
    $$
    This proves that
    $$
        V_t^{(2)}(\Pi) \xlongrightarrow{P} \langle X \rangle_t.
    $$
}

\cor{
    If $V_t^{(2)}(\Pi) \xlongrightarrow{P} \langle X \rangle_t$, then
    $$
    \begin{cases}
        \forall~ q > 2,\quad V_t^{(q)}(\Pi) \xlongrightarrow{P} 0,\\
        \forall~ q < 2,\quad V_t^{(q)}(\Pi) \xlongrightarrow{P} \infty
    \end{cases}
    $$
    on $\{\langle X \rangle_t > 0\}$.
}

We use the following notation:
$$
    \mathcal{M}_2 = \{X : \mathbb{E}X_t^2 < \infty,\ \forall~ t \geq 0\},
$$
$$
    \mathcal{M}^2 = \left\{X : \sup_{t \geq 0} \mathbb{E}X_t^2 < \infty\right\},
$$
and
$$
    \mathcal{M}_0^2 = \{X \in \mathcal{M}^2 : X_0 = 0\}.
$$


\section{\texorpdfstring{$L^2$}--Theory}

The space $\mathcal{M}^2$ is naturally identified with $L^2(\mathcal{F}_\infty)$ through the correspondence
$$
    X \longmapsto X_\infty.
$$
Indeed, for every $X \in \mathcal{M}^2$, we have
$$
\begin{cases}
    X_\infty = \lim_{t \to \infty} X_t \quad \text{a.s.},\\
    X_t \xlongrightarrow{L^2} X_\infty.
\end{cases}
$$
Conversely, if $X_\infty \in L^2(\mathcal{F}_\infty)$, then
$$
    X_t \triangleq \mathbb{E}\left( X_\infty \middle| \mathcal{F}_t \right)
$$
defines a martingale in $\mathcal{M}^2$. Moreover,
$$
    \mathbb{E}\left( X_t^2 \right)
    = \mathbb{E}\left[ \mathbb{E}\left( X_\infty \middle| \mathcal{F}_t \right)^2 \right]
    \leq \mathbb{E}\left[ \mathbb{E}\left( X_\infty^2 \middle| \mathcal{F}_t \right) \right]
    = \mathbb{E}[X_\infty^2].
$$

Thus, for $X,Y \in \mathcal{M}^2$, it is natural to define
$$
    \langle X,Y \rangle \triangleq \mathbb{E}\left[ X_{\infty}Y_{\infty} \right].
$$

\rmkb{
    Here $\langle X,Y \rangle$ denotes the inner product on the Hilbert space $\mathcal{M}^2$. It is different from the time-dependent bracket process $\langle X \rangle_t$ or $\langle X,Y \rangle_t$ introduced in the previous section.
}

\thm{Doob's $L^2$ inequality}{
    $$
        \mathbb{E}\left[ \sup_{0 \leq s \leq t}\abs{ X_s }^2 \right] \leq 4 \mathbb{E}\left[ X_t^2 \right].
    $$
    Letting $t \to \infty$, we obtain
    $$
        \mathbb{E}\left[ \sup_{s \geq 0}\abs{ X_s }^2 \right] \leq 4 \mathbb{E}[X_\infty^2].
    $$
}

The space $\mathcal{M}^2$ is complete. Indeed, if $\{M^{(n)}\}$ is a Cauchy sequence in $\mathcal{M}^2$, then $\{M_\infty^{(n)}\}$ is Cauchy in $L^2(\mathcal{F}_\infty)$. Hence there exists $M_\infty \in L^2(\mathcal{F}_\infty)$ such that
$$
    M_\infty^{(n)} \xlongrightarrow{L^2} M_\infty.
$$
Let
$$
    M_t = \mathbb{E}(M_\infty \mid \mathcal{F}_t).
$$
Then
$$
    \sup_{t \geq 0}\abs{ M_t^{(n)} - M_t } \xlongrightarrow{L^2} 0.
$$

\pf{
    By Doob's $L^2$ inequality,
    $$
    \begin{aligned}
        \mathbb{E}\left[ 
            \sup_{t \geq 0}\abs{ M_t^{(n)} - M_t }^2
        \right]
        &=
        \mathbb{E}\left[
            \sup_{t \geq 0}\abs{
                \mathbb{E}\left(
                    M_\infty^{(n)} - M_\infty \middle| \mathcal{F}_t
                \right)
            }^2
        \right] \\
        &\leq 4 \mathbb{E}\left[ \left( M_\infty^{(n)} - M_\infty \right)^2 \right] \to 0.
    \end{aligned}
    $$
}

\cor{
    The space $c\mathcal{M}^2$ of continuous martingales in $\mathcal{M}^2$ is closed in the topology of $\mathcal{M}^2$.
}

\pf{
    Let $M^{(n)} \to M$ in $\mathcal{M}^2$, with each $M^{(n)} \in c\mathcal{M}^2$. Then
    $$
        \sup_{t \geq 0} \abs{ M_t^{(n)} - M_t } \xlongrightarrow{L^2} 0.
    $$
    Hence, after passing to a subsequence if necessary,
    $$
        \sup_{t \geq 0}\abs{ M_t^{(n_k)} - M_t } \xlongrightarrow{a.s.} 0.
    $$
    Therefore $M^{(n_k)}$ converges uniformly almost surely to $M$, and since each $M^{(n_k)}$ has continuous sample paths, the limit process $M$ is also continuous.
}

\defn{Weakly orthogonal}{
    Two martingales $M,N \in \mathcal{M}_0^2$ are said to be weakly orthogonal if
    $$
        \mathbb{E}[M_\infty N_\infty] = 0.
    $$
}

\defn{Strongly orthogonal}{
    Two martingales $M,N \in \mathcal{M}_0^2$ are said to be strongly orthogonal if, for every stopping time $T$,
    $$
        \mathbb{E}[M_T N_T] = 0.
    $$
}

\lemnn{
    Two martingales are strongly orthogonal if and only if $MN$ is uniformly integrable.
}

\lemnn{
    Let $M$ be a progressively measurable process such that for every stopping time $T$,
    $$
    \begin{cases}
        \mathbb{E}\abs{M_T} < \infty,\\
        \mathbb{E}M_T = 0.
    \end{cases}
    $$
    Then $M$ is uniformly integrable.
}

\defn{Stable subspaces of $\mathcal{M}_0^2$}{
    A subspace $\mathcal{N}_0 \subseteq \mathcal{M}_0^2$ is called stable if
    \begin{itemize}
        \item[\textcircled{1}] $\mathcal{N}_0$ is a closed subspace;
        \item[\textcircled{2}] for every $N \in \mathcal{N}_0$ and every stopping time $T$, the stopped process $N^T$, defined by $N_t^T = N_{t \wedge T}$, also belongs to $\mathcal{N}_0$.
    \end{itemize}
}

\ex{}{$c\mathcal{M}_0^2$ is a stable subspace.}

\thmnn{
    Let $\mathcal{N}_0$ be a stable subspace of $\mathcal{M}_0^2$. Let $\mathcal{N}_0^\perp$ be the set of $Z \in \mathcal{M}_0^2$ which are weakly orthogonal to every element of $\mathcal{N}_0$, i.e.
    $$
        \mathcal{N}_0^\perp = \{Z \in \mathcal{M}_0^2 \colon \langle Z,N \rangle = 0,\ \forall~ N \in \mathcal{N}_0 \}.
    $$
    Then $\mathcal{N}_0^\perp$ is a stable subspace of $\mathcal{M}_0^2$. Moreover, if $Z \in \mathcal{N}_0^\perp$ and $N \in \mathcal{N}_0$, then $Z$ and $N$ are strongly orthogonal.

    Finally, every $M \in \mathcal{M}_0^2$ admits a decomposition
    $$
        M = N + Z
    $$
    with $N \in \mathcal{N}_0$ and $Z \in \mathcal{N}_0^\perp$.
}

\pf{
    Since $\mathcal{N}_0$ is a closed subspace of the Hilbert space $\mathcal{M}_0^2$, its orthogonal complement $\mathcal{N}_0^\perp$ is also a closed subspace, and every $M \in \mathcal{M}_0^2$ admits an orthogonal decomposition
    $$
        M = N + Z
    $$
    with $N \in \mathcal{N}_0$ and $Z \in \mathcal{N}_0^\perp$.

    It remains to prove that $\mathcal{N}_0^\perp$ is stable and that orthogonality is in fact strong. Let $Z \in \mathcal{N}_0^\perp$, $N \in \mathcal{N}_0$, and let $T$ be a stopping time. Then
    $$
    \begin{aligned}
        \mathbb{E}[Z_\infty^T N_\infty]
        &= \mathbb{E}[Z_T N_\infty] \\
        &= \mathbb{E}\left[ Z_T \mathbb{E}\left( N_\infty \middle| \mathcal{F}_T \right) \right] \\
        &= \mathbb{E}[Z_T N_T] \\
        &= \mathbb{E}\left[ \mathbb{E}\left( Z_\infty \middle| \mathcal{F}_T \right) N_T \right] \\
        &= \mathbb{E}[Z_\infty N_T] \\
        &= \mathbb{E}[Z_\infty N_\infty^T].
    \end{aligned}
    $$
    Since $\mathcal{N}_0$ is stable, we have $N^T \in \mathcal{N}_0$, and hence
    $$
        \mathbb{E}[Z_\infty N_\infty^T] = 0.
    $$
    Therefore
    $$
        \mathbb{E}[Z_\infty^T N_\infty] = 0
    $$
    for every $N \in \mathcal{N}_0$, so $Z^T \in \mathcal{N}_0^\perp$. This proves that $\mathcal{N}_0^\perp$ is stable.

    In particular, since $Z^T \in \mathcal{N}_0^\perp$, taking the same $N \in \mathcal{N}_0$ in the preceding identity yields
    $$
        \mathbb{E}[Z_T N_T] = \mathbb{E}[Z_\infty^T N_\infty] = 0.
    $$
    Thus $Z$ and $N$ are strongly orthogonal.
}

\defnn{
    Define
    $$
        d\mathcal{M}_0^2 \triangleq (c\mathcal{M}_0^2)^\perp.
    $$
    Then every $M \in \mathcal{M}_0^2$ can be written as
    $$
        M = M^c + M^d,
    $$
    where $M^c \in c\mathcal{M}_0^2$, $M^d \in d\mathcal{M}_0^2$, and $M^c$ and $M^d$ are strongly orthogonal. In particular, $M^cM^d$ is uniformly integrable.
}

For a c\`adl\`ag function $f$, define
$$
    \Delta f(t) = f(t) - f(t-).
$$

We now extend the quadratic variation from the continuous case to a general martingale in $\mathcal{M}_0^2$. In the continuous case, the bracket $\langle M \rangle$ is characterized by the fact that $M^2 - \langle M \rangle$ is a martingale. When jumps are present, however, the increasing process must also record the jump contribution. The next theorem shows that there is a unique increasing process $[M]$ which plays this role and satisfies $\Delta[M] = (\Delta M)^2$.

\thm{Meyer}{
    Let $M \in \mathcal{M}_0^2$. Then there exists a unique increasing process $[M]$, null at $0$, such that
    \begin{itemize}
        \item[\textcircled{1}] $M^2 - [M]$ is uniformly integrable;
        \item[\textcircled{2}] $\Delta [M] = (\Delta M)^2$.
    \end{itemize}
}

\pf{
    Write
    $$
        M = M^c + M^d,
    $$
    where $M^c \in c\mathcal{M}_0^2$ and $M^d \in d\mathcal{M}_0^2$. Then
    $$
        M_t^2 = (M_t^c + M_t^d)^2 = (M_t^c)^2 + 2M_t^cM_t^d + (M_t^d)^2.
    $$
    The theorem is obtained by combining the continuous quadratic variation of $M^c$, the purely discontinuous part $M^d$, and the strong orthogonality between $M^c$ and $M^d$.
}

For $M,N \in \mathcal{M}_0^2$, define
$$
    [M,N] \triangleq \frac{1}{4}\left( [M+N] - [M-N] \right).
$$

\thm{Kunita \& Watanabe}{
    The process $[M,N]$ is the unique finite-variation process, null at $0$, such that
    \begin{itemize}
        \item[\textcircled{1}] $MN - [M,N]$ is uniformly integrable;
        \item[\textcircled{2}] $\Delta [M,N] = \Delta M \cdot \Delta N$.
    \end{itemize}
}

\section{It\^o Integral}

Let $M \in \mathcal{M}_0^2$. We begin with the simplest class of integrands.

If $H_t = Z \mathsf{1}_{(S,T]}(t),$ where $Z \in \mathcal{F}_S$ is bounded and $S \leq T$ are stopping times, define
$$
    \int_0^t H_s \diff M_s \equiv (H \cdot M)_t \equiv Z\bigl(M_{T \wedge t} - M_{S \wedge t}\bigr).
$$

We now show that $(H \cdot M)$ belongs to $\mathcal{M}_0^2$. As a first step, we prove that $(H \cdot M)$ is a martingale. By Lemma 1.6.2, it suffices to verify that for every stopping time $R$,
$$
\begin{cases}
    \mathbb{E}\abs{(H \cdot M)_R} < \infty,\\
    \mathbb{E}(H \cdot M)_R = 0.
\end{cases}
$$

\pf{
    Let $R$ be any stopping time. Since $Z$ is bounded, say $\abs{Z} \leq K$, we have
    $$
        \mathbb{E}\abs{(H \cdot M)_R}
        \leq K\mathbb{E}\abs{M_{T \wedge R}} + K\mathbb{E}\abs{M_{S \wedge R}}.
    $$
    Because $M \in \mathcal{M}_0^2$, the martingale $M$ is uniformly integrable, and
    $$
        M_{T \wedge R} = \mathbb{E}(M_\infty \mid \mathcal{F}_{T \wedge R}).
    $$
    Hence
    $$
        \mathbb{E}\abs{M_{T \wedge R}}
        \leq \mathbb{E}\left[\mathbb{E}(\abs{M_\infty} \mid \mathcal{F}_{T \wedge R})\right]
        = \mathbb{E}\abs{M_\infty},
    $$
    and similarly for $M_{S \wedge R}$. Therefore
    $$
        \mathbb{E}\abs{(H \cdot M)_R} < \infty.
    $$

    It remains to show that $\mathbb{E}(H \cdot M)_R = 0$. By linearity, it is enough to treat the case $Z = \mathsf{1}_A$ with $A \in \mathcal{F}_S$. Then
    $$
        \mathbb{E}(H \cdot M)_R
        = \mathbb{E}\left[\mathsf{1}_A(M_{T \wedge R} - M_{S \wedge R})\right].
    $$
    Define the stopping times
    $$
        S_A \equiv
        \begin{cases}
            S, & \text{on } A,\\
            \infty, & \text{on } A^c.
        \end{cases}
        \qquad
        T_A \equiv
        \begin{cases}
            T, & \text{on } A,\\
            \infty, & \text{on } A^c.
        \end{cases}
    $$
    Then
    $$
        \mathsf{1}_A M_{S \wedge R} = M_{S_A \wedge R},
        \qquad
        \mathsf{1}_A M_{T \wedge R} = M_{T_A \wedge R}.
    $$
    Hence the expectation above is the difference of expectations of stopped values of the martingale $M$, and therefore equals $0$ by optional sampling. Thus $(H \cdot M)$ is a martingale.
}

\lemnn{
    If $Z \in b\mathcal{F}_S$ and $M$ is a uniformly integrable martingale, then $H \cdot M$ is a uniformly integrable martingale.
}

More generally, consider a finite linear combination of such elementary processes:
$$
    H_t = \sum_{i=1}^n Z_i \mathsf{1}_{(S_i,T_i]}(t),
$$
where each $Z_i \in b\mathcal{F}_{S_i}$.

We define
$$
    (H \cdot M)_t
    \triangleq
    \sum_{i=1}^{n} Z_i \bigl(M_{T_i \wedge t} - M_{S_i \wedge t}\bigr).
$$

\cor{
    If $Z_i \in b \mathcal{F}_{S_i}$ and $M$ is U.I. martingale, then $(H\cdot M)_{t}$ is U.I. martingale.
}


We next introduce the predictable $\sigma$-algebra, which is the natural measurability class for It\^o integrands.

\defn{Previsible process}{
    The previsible $\sigma$-algebra $\mathcal{P}$ on $(0,\infty) \times \Omega$ is defined to be the smallest $\sigma$-algebra for which every adapted $L$-process is $\mathcal{P}$-measurable.

    An $L$-process is a process $(X_t)_{t>0}$ such that $X_t = X_{t-}$ and $X_{t+}$ exists.

    We say that a process $(X_t)_{t>0}$ is previsible if it is $\mathcal{P}$-measurable.
}

\lemnn{
    Suppose $S$ and $T$ are stopping times with $S \leq T \leq \infty$. Let $Z \in b\mathcal{F}_S$ and $H = Z\mathsf{1}_{(S,T]}$. Then $H$ is previsible.
}

\pf{
    Fix $t > 0$ and let $B$ be a Borel set. We may assume that $0 \notin B$. Then
    $$
        \{H_t \in B\}
        = \{Z \in B\} \cap \{S < t\} \cap \{t \leq T\}.
    $$
    Moreover,
    $$
        \{S < t\} = \bigcup_{n=1}^{\infty}\left\{S \leq t - \frac{1}{n}\right\}.
    $$
    Since $\{Z \in B\} \in \mathcal{F}_S$, we have
    $$
        \{Z \in B\} \cap \left\{S \leq t - \frac{1}{n}\right\} \in \mathcal{F}_{t - \frac{1}{n}} \subseteq \mathcal{F}_t.
    $$
    Also, $\{t \leq T\} \in \mathcal{F}_t$. Therefore $\{H_t \in B\} \in \mathcal{F}_t$, so $H$ is adapted. Since $H$ is clearly an $L$-process, it follows from the definition of $\mathcal{P}$ that $H$ is previsible.
}

\defnn{
    Let
    $$
        b\varepsilon
        = \left\{
            \sum_{i=1}^n Z_i \mathsf{1}_{(S_i,T_i]}
            :
            n \geq 1,\,
            Z_i \in b\mathcal{F}_{S_i},\,
            S_1 \leq T_1 \leq S_2 \leq \cdots
        \right\}.
    $$
}

\noindent\textbf{Claim.} $\sigma(b\varepsilon) = \mathcal{P}$.

\pf{
    ``$\subseteq$'' The inclusion $\sigma(b\varepsilon) \subseteq \mathcal{P}$ follows from the previous lemma.

    ``$\supseteq$'' For any bounded left-continuous process,
    $$
    H_t = \lim_{k \uparrow \infty} \lim_{n \uparrow \infty} \sum_{i=2}^{n_k} H_{\frac{i-1}{n}} \mathsf{1}_{\left(\frac{i-1}{n},\frac{i}{n}\right]}(t).
    $$

    Since $H$ is adapted, we get $H_{\frac{i-1}{n}} \in b\mathcal{F}_{\frac{i-1}{n}}$.

    Hence
    $$
    \sum_{i=2}^{n_k} H_{\frac{i-1}{n}} \mathsf{1}_{\left( \frac{i-1}{n}, \frac{i}{n} \right]}(t) \in b \varepsilon.
    $$  

    $\implies H_t \in \sigma(b \varepsilon).$

}
\pf{

    
}
\lemnn{
    Let $H \in b\varepsilon$ and $M \in \mathcal{M}_0^2$. Then $H \cdot M \in \mathcal{M}_0^2$, and if
    $$
        H = \sum_{i=1}^n Z_{i-1}\mathsf{1}_{(T_{i-1},T_i]},
    $$
    then
    $$
        \mathbb{E}(H \cdot M)_\infty^2
        =
        \sum_{i=1}^{n}\mathbb{E}\left[ Z_{i-1}^2(M_{T_i} - M_{T_{i-1}})^2 \right].
    $$
}


\pf{
    $$
        (H \cdot M)_t = \sum_{i=1}^n Z_{i-1}\left( M_{T_i \wedge t} - M_{T_{i-1} \wedge t} \right).
    $$
    By the previous corollary 1.7.1, $(H \cdot M)$ is a U.I. martingale.

    Hence $(H \cdot M)_t \to (H \cdot M)_\infty$ a.s. and in $L^1$, where
    $$
        (H \cdot M)_\infty = \sum_{i=1}^n Z_{i-1}\Delta_i,
        \qquad
        \Delta_i = M_{T_i} - M_{T_{i-1}}.
    $$

    Therefore
    $$
        \mathbb{E}\left[ (H \cdot M)_\infty^2 \right]
        = \mathbb{E} \sum_{i=1}^n Z_{i-1}^2 \Delta_i^2
        + 2 \mathbb{E}\left( \sum_{i < j} Z_{i-1}\Delta_i Z_{j-1}\Delta_j \right).
    $$

    For $i<j$, the random variable $Z_{i-1}\Delta_i Z_{j-1}$ is $\mathcal{F}_{T_{j-1}}$-measurable, so
    $$
        \mathbb{E}\left[ Z_{i-1}\Delta_i Z_{j-1}\Delta_j \right]
        =
        \mathbb{E}\left[
            Z_{i-1}\Delta_i Z_{j-1}
            \mathbb{E}\left( \Delta_j \middle| \mathcal{F}_{T_{j-1}} \right)
        \right]
        = 0.
    $$
    Thus all cross terms vanish.

    By Doob's $L^2$ inequality,
    $$
        \mathbb{E}\left[
            \sup_{t \geq 0}(H \cdot M)_t^2
        \right]
        \leq 4 \mathbb{E}\left[ (H \cdot M)_\infty^2 \right]
        < \infty.
    $$
    Therefore $H \cdot M \in \mathcal{M}_0^2$.
}

Recall from Meyer's theorem that for $M \in \mathcal{M}_0^2$ there exists a quadratic variation $[M]$ such that $M^2 - [M]$ is a uniformly integrable martingale. Applying this to the previous lemma 1.7.3, we get
$$
    \mathbb{E}(H \cdot M)_\infty^2
    =
    \sum_{i=1}^n \mathbb{E}\left[
        Z_{i-1}^2\bigl(M_{T_i} - M_{T_{i-1}}\bigr)^2
    \right].
$$
Moreover,
$$
    \begin{aligned}
        \mathbb{E}\left[
            \bigl(M_{T_i} - M_{T_{i-1}}\bigr)^2 \middle| \mathcal{F}_{T_{i-1}}
        \right]
        &=
        \mathbb{E}\left[
            M_{T_i}^2 - M_{T_{i-1}}^2 \middle| \mathcal{F}_{T_{i-1}}
        \right] \\
        &=
        \mathbb{E}\left[
            [M]_{T_i} - [M]_{T_{i-1}} \middle| \mathcal{F}_{T_{i-1}}
        \right],
    \end{aligned}
$$
because $M^2 - [M]$ is a martingale and $\mathbb{E}(M_{T_i}\mid \mathcal{F}_{T_{i-1}})=M_{T_{i-1}}$.
Therefore
$$
    \mathbb{E}(H \cdot M)_\infty^2
    =
    \sum_{i=1}^{n}\mathbb{E}\left[
        Z_{i-1}^2\bigl([M]_{T_i} - [M]_{T_{i-1}}\bigr)
    \right]
    =
    \mathbb{E}\left[
        \int_0^\infty H_s^2 \diff [M]_s
    \right].
$$

This shows that, for simple predictable integrands, the map
$$
    H \mapsto (H \cdot M)_\infty
$$
is an isometry from $(L^2(M)\cap b\varepsilon,\norm{\cdot}_M)$ into $L^2(\mathcal{F}_\infty)$ equipped with the usual $L^2$ norm.

For fixed $M \in \mathcal{M}_0^2$, define
$$
    \norm{H}_M
    =
    \left(
        \mathbb{E}\int_0^\infty H_s^2 \diff [M]_s
    \right)^{1/2}.
$$
Let
$$
    L^2(M)
    =
    \left\{
        \text{previsible } H : \norm{H}_M < \infty
    \right\}.
$$

Then the map
$$
\begin{aligned}
    I : L^2(M) \cap b\varepsilon &\to L^2(\mathcal{F}_\infty),\\
    H &\mapsto (H \cdot M)_\infty
\end{aligned}
$$
is an isometry, since
$$
    \mathbb{E}\left[
        \int_0^\infty H_s^2 \diff [M]_s
    \right]
    =
    \mathbb{E}(H \cdot M)_\infty^2.
$$
Hence $I$ extends uniquely from simple predictable processes to all of $L^2(M)$.

\defn{It\'o Integral}{
    Let $M \in \mathcal{M}_{0}^2$ and $H \in L^2(M)$. Then It\'o Integral $(H\cdot M)_{t}$ is the image of $H$ under extension of $I : L^2(M)\to L^2(\mathcal{F}_{\infty})$.
}

\noindent\textbf{Localization.}

\defnn{
    Let $T$ be a stopping time and let $H$ be any process. Define
    $$
        H(0,T](t)
        =
        \begin{cases}
            H_t, & 0 < t \leq T,\\
            0, & t > T.
        \end{cases}
    $$
}

\rmkb{
    If $H$ is previsible, then so is $H(0,T]$.
}

\defnn{
    Let $T$ be a stopping time and let $X$ be any process. Define
    $$
        X^T(t)
        =
        \begin{cases}
            X(t), & 0 \leq t \leq T,\\
            X(T), & t > T.
        \end{cases}
    $$
}
For an elementary integrand, the localization identity can be checked directly. Indeed,
$$
    (H \cdot M)^T_t = (H(0,T]\cdot M^T)_t,
$$
and
$$
    \left[ Z_i \left( M_{T_i \wedge t} - M_{T_{i-1} \wedge t} \right) \right]^T
    =
    \left( Z_i \mathsf{1}_{(T_{i-1}\wedge T,\; T_i \wedge T]} \cdot M^T \right)_t
    =
    Z_i\left( M_{T_i \wedge T \wedge t} - M_{T_{i-1}\wedge T \wedge t} \right).
$$
\thmnn{
    Fix $M \in \mathcal{M}_0^2$. For any $H \in L^2(M)$ and any stopping time $T$:
    \begin{itemize}
        \item[(i)] $(H \cdot M)^T = H(0,T] \cdot M = H \cdot M^T$.
        \item[(ii)] $(H \cdot M)_t^2 - \int_0^t H_s^2 \diff [M]_s$ is a uniformly integrable martingale. In particular,
        $$
            [H \cdot M]_t = \int_0^t H_s^2 \diff [M]_s.
        $$
        \item[(iii)] If $H, K \in b\mathcal{P}$, then
        $$
            (HK)\cdot M = H \cdot (K \cdot M) = K \cdot (H \cdot M).
        $$
        \item[(iv)] $\Delta(H \cdot M) = H \Delta M$.
        \item[(v)] If $M$ is continuous, then $H \cdot M$ is continuous.
    \end{itemize}
}

\pf{
    We first prove the identities for elementary processes and then extend them by density.

    For (i), let $H = Z\mathsf{1}_{(U,V]}$. Then
    $$
        (H \cdot M)_t = Z(M_{V \wedge t} - M_{U \wedge t}),
    $$
    so
    $$
        (H \cdot M)^T_t = (H \cdot M)_{t \wedge T}
        = Z(M_{V \wedge T \wedge t} - M_{U \wedge T \wedge t}).
    $$
    On the other hand,
    $$
        H(0,T] = Z\mathsf{1}_{(U \wedge T, V \wedge T]},
    $$
    and therefore
    $$
        (H(0,T]\cdot M)_t
        = Z(M_{V \wedge T \wedge t} - M_{U \wedge T \wedge t}).
    $$
    Similarly,
    $$
        (H \cdot M^T)_t = Z(M^T_{V \wedge t} - M^T_{U \wedge t})
        = Z(M_{V \wedge T \wedge t} - M_{U \wedge T \wedge t}).
    $$
    Hence (i) holds for elementary processes, and then for all $H \in L^2(M)$ by continuity of the integral map.

    For (ii), set
    $$
        N_t = (H \cdot M)_t^2 - \int_0^t H_s^2 \diff [M]_s.
    $$
    Since
    $$
        \mathbb{E}\left[\sup_{t \geq 0}(H \cdot M)_t^2\right]
        \leq 4\mathbb{E}\left[\int_0^\infty H_s^2 \diff [M]_s\right] < \infty,
    $$
    the family $\{N_t : t \geq 0\}$ is uniformly integrable. It therefore suffices to show that $\mathbb{E}N_T = 0$ for every stopping time $T$.

    By part (i),
    $$
        (H \cdot M)_T = (H(0,T]\cdot M)_\infty.
    $$
    Applying the isometry to the integrand $H(0,T]$, we obtain
    $$
    \begin{aligned}
        \mathbb{E}(H \cdot M)_T^2
        &= \mathbb{E}(H(0,T]\cdot M)_\infty^2 \\
        &= \mathbb{E}\int_0^\infty H(0,T]_s^2 \diff [M]_s \\
        &= \mathbb{E}\int_0^T H_s^2 \diff [M]_s.
    \end{aligned}
    $$
    Hence $\mathbb{E}N_T = 0$, and (ii) follows.

    For (iii) First take $H = Z_1 \mathsf{1}_{(S_1,T_1]}$ and $K = Z_2 \mathsf{1}_{(S_2,T_2]}$. Then
    $$
        HK = Z_1Z_2\,\mathsf{1}_{(S,T]},
    $$
    where
    $$
        S = S_1 \vee S_2,
        \qquad
        T = (T_1 \wedge T_2) \vee (S_1 \vee S_2).
    $$
    Hence
    $$
        ((HK)\cdot M)_\infty = Z_1Z_2(M_T - M_S).
    $$
    On the other hand,
    $$
    \begin{aligned}
        (H \cdot (K \cdot M))_\infty
        &= Z_1\Bigl((K \cdot M)_{T_1} - (K \cdot M)_{S_1}\Bigr) \\
        &= Z_1\Bigl((K \cdot M^{T_1})_\infty - (K \cdot M^{S_1})_\infty\Bigr) \\
        &= Z_1Z_2\Bigl[(M_{T_2 \wedge T_1} - M_{S_2 \wedge T_1}) - (M_{T_2 \wedge S_1} - M_{S_2 \wedge S_1})\Bigr] \\
        &= Z_1Z_2(M_T - M_S).
    \end{aligned}
    $$
    Therefore, for every $H,K \in b\varepsilon$,
    $$
        H \cdot (K \cdot M) = (HK)\cdot M.
    $$
    The identity
    $$
        K \cdot (H \cdot M) = (HK)\cdot M
    $$
    is proved in the same way.

    Now let $K \in b\mathcal{P}$. Choose $(K^n)_{n\geq 1} \subset b\varepsilon$ such that
    $$
        K^n \xlongrightarrow{L^2(M)} K.
    $$
    Fix $H \in b\varepsilon$. By the isometry,
    $$
    \begin{aligned}
        \mathbb{E}\left[\left(H \cdot (K\cdot M)-H\cdot(K^n\cdot M)\right)_\infty^2\right]
        &= \mathbb{E}\left[\left(H\cdot\bigl((K-K^n)\cdot M\bigr)\right)_\infty^2\right] \\
        &= \mathbb{E}\left[\int_0^\infty H_s^2 \diff [ (K-K^n)\cdot M ]_s\right] \\
        &= \mathbb{E}\left[\int_0^\infty H_s^2 (K_s-K_s^n)^2 \diff [M]_s\right] \\
        &\leq \norm{H}_\infty^2\,\mathbb{E}\left[\int_0^\infty (K_s-K_s^n)^2 \diff [M]_s\right]
        \to 0.
    \end{aligned}
    $$
    Hence
    $$
        H \cdot (K^n\cdot M) \xlongrightarrow{L^2(\mathcal{F}_\infty)} H \cdot (K\cdot M).
    $$

    On the other hand, for each $n$ we already know that
    $$
        H \cdot (K^n\cdot M) = (HK^n)\cdot M.
    $$
    Since $H \in b\varepsilon$, we also have
    $$
        \mathbb{E}\left[\int_0^\infty (H_sK_s^n-H_sK_s)^2 \diff [M]_s\right]
        \leq \norm{H}_\infty^2\,\mathbb{E}\left[\int_0^\infty (K_s^n-K_s)^2 \diff [M]_s\right]
        \to 0,
    $$
    so $HK^n \xlongrightarrow{L^2(M)} HK$. Applying the isometry once again,
    $$
        (HK^n)\cdot M \xlongrightarrow{L^2(\mathcal{F}_\infty)} (HK)\cdot M.
    $$
    By uniqueness of the $L^2(\mathcal{F}_\infty)$ limit,
    $$
        H \cdot (K\cdot M) = (HK)\cdot M.
    $$
    (iv) First let $H = Z\mathsf{1}_{(S,T]}$. Then
    $$
        (H \cdot M)_t = Z(M_{t \wedge T} - M_{t \wedge S}),
    $$
    so
    $$
        \Delta(H \cdot M)_t = Z\,\mathsf{1}_{(S,T]}(t)\,\Delta M_t = H_t \Delta M_t.
    $$
    Hence the identity holds for all $H \in b\varepsilon$.

    Now let $H \in L^2(M)$, and choose $H^n \in b\varepsilon$ such that
    $$
        H^n \xlongrightarrow{L^2(M)} H.
    $$
    By the isometry,
    $$
        H^n \cdot M \xlongrightarrow{L^2(\mathcal{F}_\infty)} H \cdot M,
    $$
    and hence, after passing to a subsequence,
    $$
        \sup_{t \geq 0}\abs{(H^{n_k}\cdot M)_t - (H\cdot M)_t} \xlongrightarrow{a.s.} 0.
    $$
    Also,
    $$
        \int_0^\infty (H^{n_k}_s - H_s)^2 \diff [M]_s \xlongrightarrow{a.s.} 0
    $$
    along a further subsequence.

    Fix $t \geq 0$. Since
    $$
        \Delta(H^{n_k}\cdot M)_t = H^{n_k}_t \Delta M_t,
    $$
    we distinguish two cases.

    \textbf{Case 1.} If $\Delta M_t = 0$, then $\Delta(H^{n_k}\cdot M)_t = 0$ for every $k$, and therefore
    $$
        \Delta(H \cdot M)_t = 0 = H_t \Delta M_t.
    $$

    \textbf{Case 2.} If $\Delta M_t \neq 0$, then $\Delta[M]_t = (\Delta M_t)^2 > 0$. Hence
    $$
        \int_0^\infty (H^{n_k}_s - H_s)^2 \diff [M]_s
        \geq
        (H^{n_k}_t - H_t)^2 \Delta[M]_t,
    $$
    and so $H^{n_k}_t \to H_t$. Letting $k \to \infty$ in
    $$
        \Delta(H^{n_k}\cdot M)_t = H^{n_k}_t \Delta M_t,
    $$
    we obtain
    $$
        \Delta(H \cdot M)_t = H_t \Delta M_t.
    $$
}


\thm{Kunita-Watanabe}{
    Let $M \in \mathcal{M}_0^2$ and $H \in L^2(M)$. Then $\forall~N \in \mathcal{M}_0^2$,
    $$
    \mathbb{E}\left[ 
        (H \cdot M)_{\infty} \cdot N_\infty
     \right] = \mathbb{E}\left[ 
        (H \cdot [M,N])_{\infty}
      \right].
    $$
    Moreover, 
    $$
        [H\cdot M , N] = H \cdot [M,N] \xlongequal{def}\int H_s \diff [M,N]_s.
    $$
}

\thmnn{
    Let $M,N \in \mathcal{M}_0^2$. Then 
    $$
    \left( \int_0^\infty \abs{ H_s K_s } \diff [M,N]_s \right)^2 \leq \int_0^\infty H_s^2 \diff [M]_s \cdot \int_0^\infty K_s ^2 \diff [N]_s. \quad \text{a.s.}
    $$
}

\pf{
    $\abs{ H_s K_s } = H_s K_s \cdot \mathrm{sgn}(H_s K_s)$.

    $\abs{ \diff [M,N]_s } = \diff [M,N]_s \cdot \mathrm{sgn}([M,N]_s)$.
    $$
    \begin{aligned}
        LHS &= \int_0^\infty H_s K_s \cdot \mathrm{sgn}(H_s K_s) \cdot \mathrm{sgn} ([M,N]_s) \diff [M,N]_s\\
        &=\int_{0}^{\infty} \tilde{H_s}K_s \diff [M,N]_s
    \end{aligned}
    $$
    $$
    \begin{aligned}
        RHS &= \int_0^\infty \tilde{H_s} \diff [M]_s \cdot \int_0^\infty K_s^2 \diff [N]_s
    \end{aligned}
    $$
    $$
    \Leftrightarrow \left( \int_0^\infty \tilde{H_s} K_s \diff [M,N]_s  \right)^2 \leq \int_0^\infty \tilde{H_s} \diff [M]_s \int_0^\infty K_s ^2 \diff [N]_s.
    $$
    $$
    \Leftrightarrow \begin{aligned}
        & H_n = \sum_{i=1}^n U_i (t_i,t_{i+1}]\\
        & K_n \equiv \sum_{i=1}^n V_i(t_i,t_{i+1}].
    \end{aligned}
    $$
    $$
    \Leftrightarrow \left( \sum_{i=1}^{n} U_i V_i \left( [M,N]_{t_{i+1}} - [M,N]_{t_i} \right) \right)^2 \leq \sum_{i=1}^2 U_i^2([M]_{t_{i+1}} - [M]_{t_i} ) \cdot \sum_{i=1}^n V_i^2 ([N]_{t_{i+1}}-[N]_{t_i}).
    $$

    It suffices to prove
    $$
    \abs{ [M,N]_{t_{i+1}} - [M,N]_{t_i} } \leq \left( [M]_{t_{i+1}} - [M]_{t_i} \right)^{\frac{1}{2}} \left( 
        [N]_{t_{i+1}} - [N]_{t_i}
     \right)^{\frac{1}{2}}.
    $$
    like $(\sum_i a_i b_i)^2 \leq \sum_i a_i^2 \sum_i b_i^2.$

    To prove this, note that $\forall~ \lambda \in \mathbb{Q}$, $[M+ \lambda N]$ is increasing and so 
    $$
        [M+ \lambda N]_{t_{i+1}} - [M+\lambda N]_{t_i} \geq 0 \quad \text{a.s..}
    $$
    
    Also 
    $$
        [M + \lambda N]_t = [M]_t + 2\lambda[M,N]_t + \lambda^2 [N]_t.
    $$
    because
    $
    (M+ \lambda N)_{t}^2 - [M+ \lambda N]_t
    $ is U.I. and by the uniqueness (?)
    $$
    M_t^2  + 2 \lambda M_t N_t + \lambda^2 N_t^2 -([M]_t + 2 \lambda [M,N]_t + \lambda^2 [N]_t).
    $$
    Hence 
    $$
    \left( [M]_{t_{i+1}} - [M]_{t_i} \right) + 2 \lambda \left( [M,N]_{t_{i+1}} - [M,N]_{t_i} \right) + \lambda^2 \left( [N]_{t_{i+1}} -[N]_{t_i} \right) \geq 0
    $$
    $$ 
    \Leftrightarrow b^2 - 4ac \leq 0.
    $$
    So
    $$
        \left( [M,N]_{t_{i+1}} - [M,N]_{t_i} \right)^2 \leq \left( [M]_{t_{i+1}} - [M]_{t_i} \right)\left( [N]_{t_{i+1}} - [N]_{t_i} \right).
    $$
    $$
        \mathbb{E}\left[ 
            (H\cdot M)_{\infty} N_{\infty}
         \right] = \mathbb{E}\left[ 
            (H \cdot [M,N])_{\infty}
          \right].
    $$

    If $H = Z(S,T]$, 
    $$\begin{cases}
     LHS = \mathbb{E}\left[ Z(M_T - M_S) \cdot N_\infty \right]\\
     RHS = \mathbb{E}\left[ Z([M,N]_{T} - [M,N]_S) \right]   
    \end{cases}$$

    $$
    \begin{aligned}
        LHS &= \mathbb{E}\left[ \mathbb{E}(N_\infty \mid \mathcal{F}_T) \cdot Z(M_T - M_S) \right] \\
        &= \mathbb{E}\left[ Z(M_T N_T - M_S N_T) \right]\\
        &= \mathbb{E}[ZM_TN_T] - \mathbb{E}\left[ \mathbb{E}(N_T \mid \mathcal{F}_S)\cdot ZM_S \right]\\
        &=\mathbb{E}\left[ Z M_T N_T \right] - \mathbb{E}[ZM_S N_S].
    \end{aligned}
    $$
    
    Because $M_t N_t - [M,N]_t$ is U.I.
    $$
    \mathbb{E}[M_T N_T] = \mathbb{E}[[M,N]_T] \implies \mathbb{E}[Z(M_T N_T - [M,N]_T)] = 0.
    $$
    And
    $$
    \begin{aligned}
        LHS &= \mathbb{E}\left[ Z(M_T N_T - M_S N_S) \right] \\
        &= \mathbb{E}\left[ Z \mathbb{E}\left[ M_TN_T -M_S N_S \mid \mathcal{F}_S \right] \right]\\
        &= \mathbb{E}\left[ Z \mathbb{E}[[M,N]_T - [M,N]_S \mid \mathcal{F}_S] \right]\\
        &=\mathbb{E}[Z([M,N]_T - [M,N]_S)].
    \end{aligned}
    $$
    Hence $\forall~ H \in b \varepsilon$
    $$
    \mathbb{E}[(H\cdot M)_{\infty} \cdot N_{\infty}] = \mathbb{E}\left[ (H \cdot [M,N])_{\infty} \right].
    $$
    
    If $H_n \xlongrightarrow{L^2(M)}H$, i.e.
    $$
    \mathbb{E}\left( \int_0^\infty (H_n(s) - H_s)^2 \diff [M]_s \right) \to 0
    $$
    $\implies (H_n \cdot M)_{\infty} \xlongrightarrow{L^2} (H\cdot M)_{\infty}$ by 
    $$
    \mathbb{E}[(H_n \cdot M - H\cdot M)_{\infty}^2] = \mathbb{E}\left[ \int_{0}^{\infty}(H_n - H) \diff [M]_s \right] \to 0.
    $$
    
    Hence $\mathbb{E}[(H_n \cdot M)_{\infty} \cdot N_{\infty}] \to \mathbb{E}[(H\cdot M)_{\infty} N_\infty]$.

    Next,
    $$
    \begin{aligned}
    \mathbb{E}\left[ \abs{ (H_n \cdot [M,N])_{\infty} - (H \cdot [M,N])_{\infty} } \right] &= \mathbb{E}\abs{ ((H_n - H) \cdot [M,N])_{\infty} } \\
    &\leq \left( \mathbb{E}\left( \int_0^\infty(H_n -H) \diff [M]_s \right) \right)^{\frac{1}{2}}\left( \mathbb{E}[N]_{\infty} \right)^{\frac{1}{2}} \to 0
    \end{aligned}
    $$
    by using 
    $$
    \left( \int_0^\infty \abs{ H_s } \diff [M,N]_s \right) \leq \left( \int_0^\infty H_s^2 \diff [M]_s \right)^{\frac{1}{2}} \left( \int_0^\infty 1 \diff [N]_s \right)^{\frac{1}{2}}.
    $$

    Hence $\mathbb{E}[(H_n \cdot [M,N])_{\infty}] \to \mathbb{E}[(H\cdot [M,N])_{\infty}]$.
}

\rmkb{
    proof of $[H\cdot M,N] = H \cdot [M,N]$
    $\Leftrightarrow$ It suffices to. prove $(H\cdot M)_{t} N_t - \int_{0}^\infty H_s \diff [M,N]_s$ is a U.I. martingale.
    
    $\Leftrightarrow \forall~$stopping time $R$
    $$
    \mathbb{E}[(H\cdot M)_{R} \cdot N_R] = \mathbb{E}\left[ \int_0^R H_s \diff [M,N]_s \right].
    $$
    $$
    \begin{aligned}
        LHS &= \mathbb{E}\left[ (H\cdot M)_{\infty}^R \cdot N_\infty^R \right] \\
        &= \mathbb{E}\left[ 
        (H^R \cdot M)_{\infty} \cdot N_\infty^R
        \right] \\
        &= \mathbb{E}[(H^R \cdot [M,N^R])_{\infty}]\\
        &= \mathbb{E}\left[ \int_0^\infty H^R(s) \diff [M,N^R]_s \right]\\
        &= \mathbb{E}\left[ \int_0^R H(s) \diff [M,N]_s \right] = RHS.
    \end{aligned}
    $$
}

\defn{local martingale}{
    $\mathcal{M}_{loc}=\{\text{local martingale}\}$. 
    
    $\exists~$stopping times $T_n \uparrow \infty$ s.t. $M^{T_n}$ is a martingale.
}
\defn{continuous local martingale}{
    $c\mathcal{M}_{loc}=\{\text{continuous local martingale}\}$.

    Quadratic variation of $c \mathcal{M}_{loc}$. 
    
    If $M \in c\mathcal{M}_{loc}$, then $\exists~[M]$ increasing process such that $M^2 - [M] \in c \mathcal{M}_{loc}$.
}

\lemnn{
    If $M \in c \mathcal{M}_{loc}$. Fix $n \geq 1$, define
    $$
    T_n = \inf \{t \geq 0\colon \abs{ M_t }\geq n\}.
    $$
    Then $M^{T_n}$ is a bounded martingale.
}

\pf{
    By definition, $\exists~$stopping $(S_k)\uparrow \infty$ s.t. $M^{S_k}$ is a martingale.

    $\implies M^{S_k \wedge T_n}$ is also a martingale.

    So $\forall~ s < t$, $$\mathbb{E}(M^{S_k \wedge T_n}_t \mid \mathcal{F}_s) = M_s^{S_k \wedge T_n}.$$

    Let $k \to \infty$ to get $\mathbb{E}(M_t^{T_n} \mid \mathcal{F}_s) = M_s^{T_n}.$ 

    pf of Quadratic variation:

    \textbf{(Existence.)} $\forall~ n \geq 1$, we let
    $$
    T_n = \inf \{t \geq 0 \colon \abs{ M_t } \geq n\}.
    $$
    Then define $\forall~ t < T_n$,
    $$
    [M]_t = [M^{T_n}]_t .
    $$

    We need to check that $[M^{T_n}]_t$ is well-defined.
    $$
    \begin{aligned}
        [M^{T_2}]^{T_1} = [M^{T_2}]_{t\wedge T_1}
    \end{aligned}
    $$
    Because $[M]^T = [M^T]$,
    $$
    [M^{T_2}]^{T_1} = [M^{T_2 \wedge T_1}] =[M^{T_1}].
    $$
    By $T_n \uparrow \infty$, we have $[M]_t$ is defined for all $t \geq 0$.

    \textbf{(uniqueness.)} If $\exists~A,A'$ increasing such that $N = M^2 - A$ and $N' =M^2 - A' \in c \mathcal{M}_{loc}$.
    Then $N-N'= A' - A \in c \mathcal{M}_{loc}$ is of finite variation.
    
    \lemnn{
        If $M \in c \mathcal{M}_{loc}$ and $M$ is of finite variation, then $M_t \equiv 0.$
    }
    \pf{
        Let $V_M(t) =$ total variation of $M$ on $[0,t]$. Define $T_n = \inf \{t \geq 0 \colon V_M(t) > n\}.$ 

        Then $M^{T_n}$ is a bounded martingale and is of finite variation. Next, define 
        $$
        A_t^n \equiv \sum_{k=1}^{2^n}\left( M(t_{k}^n) - M(t_{k-1}^n) \right)^2
        $$
        with $t_k^n = \frac{k}{2^n} t$.

        So
        $$
        A_t^n \leq \sup_{1 \leq k \leq 2^n} \abs{ M(t_k^n) - M(t_{k-1}^n) }\cdot V_M(t) \xlongrightarrow{a.s.} 0
        $$
        by continuity of $M$.

        Thus $A_t^n \to 0$, $\forall~ t \geq 0$.

        We will prove $A_t^n \to A_t$ such that $M_t^2 - A_t$ is martingale.

        $\implies A_t = 0$ and $\mathbb{E}M_t^2 = \mathbb{E}M_0^2 = 0.$

        $\implies M_t = 0$ a.s. $\forall~ t \geq 0$.

        By a.s. continuity, $M_t \equiv 0$, a.s.
    }
}

\subsection{It\^o Integral of local martingale}

\thmnn{
    Let $M \in c \mathcal{M}_{loc}$ with $M_0 = 0$. Let $H \in \mathcal{P}$ such that W.P.1
    $$
    \int_0^t H_s ^2 \diff [M]_s < \infty
    $$
    for all $t > 0$. Then $H \cdot M$ is well-defined such that $H\cdot M \in c \mathcal{M}_{0,loc}$.
}
\pf{
    Let $U_n = \inf \{t \colon \abs{ M_t } > n\}$, $V_n = \inf\{t \colon \int_0^t H_s^2 \diff [M]_s > n\}.$

    (Recall $L^2(M) \implies \mathbb{E}\int_0^\infty H_s^2 \diff [M]_s < \infty$).

    $\implies T_n = U_n \wedge V_n$. In this way, $M^{T_n} \in c \mathcal{M}_0^2$ and $H^{T_n} \in L^2(M^{T_n}).$

    So $H^{T_n} \cdot M^{T_n}$ is defined by previous $L^2$ theory.

    $\implies (H\cdot M)^{T_n} = H^{T_n} \cdot M^{T_n}$ ($t < T_n$) is a $L^2$-bounded martingale.

    $\implies H \cdot M \in c \mathcal{M}_{0,loc}$.
}

\defn{semimartingale}{
    We say $X$ is a semimartingale if 
    $$
    X_t = X_0 + M_t + A_t.
    $$
    $M_t$ is local mart, $A_t$ is finite variation.
}

\defnn{
    Let $X$ be a continuous semimartingale. Then $\exists~$unique 
    $$
    \begin{cases}
        M \in c \mathcal{M}_{0,loc}\\
        A \in c F V_0
    \end{cases}
    $$
    such that $X_t = X_0 + M_t + A_t$, $\forall~ t \geq 0$.
}

\pf{
    If $\exists~M'$ and $A'$ such that 
    $$
        X= X_0 + M' + A' = X_0 + M + A.
    $$
    Then $M - M' = A' - A \in c \mathcal{M}_{0,loc} \cap c F V = \{0\}$.
}

We write $M= X^{cm}$ and 
$$
    [X,X] = [M,M] = [X^{cm},X^{cm}].
$$

Similarly, $[X,Y] = [X^{cm},Y^{cm}]$.

For any $H \in lb \mathcal{P}$ define
$$
    H \cdot X =H \cdot M + H \cdot A.
$$
Obviously $H\cdot M \in c \mathcal{M}_{0,loc}$ and $H \cdot A \in FV$.

And 
$$
\begin{aligned}
[H\cdot X, K \cdot Y] &= [H \cdot X^{cm},K\cdot Y^{cm}] \\
&= H \cdot [X^{cm},K\cdot Y^{cm}] \\
&= H\cdot K [X^{cm},Y^{cm}]\\
&= (HK)\cdot [X^{cm},Y^{cm}]=(HK)\cdot [X,Y].
\end{aligned}$$

\subsection{It\^o's Formula} 
\thm{It\^o's Formula}{
    Let $f \colon \mathbb{R}^n \to \mathbb{R}$ be $C^2$. Let 
    $$
    X_t \equiv(X_t^1,\cdots,X_t^n)
    $$
    be a continuous semimartingale.

    Then 
    $$
    f(x_t) - f(x_0) = \int_0^t \nabla f(X_s) \diff X_s + \frac{1}{2}\int_0^t \sum_{i,j} D_{ij}f(X_s)\diff [X^i,X^j]_{s}
    $$
    where $\nabla f(X_s) = (D_1f(X_s),\cdots,D_n f(X_s))$. $D_i f(x) = \frac{\partial}{\partial x_i}f(x)$, $D_{ij} = D_i D_j = D_j D_i.$
}

\thm{Integration by Parts}{
    Let $X,Y$ be continuous semimartingales. Then 
    $$
    X_tY_t - X_0Y_0 = \int_0^t X_s \diff Y_s + \int_0^t Y_s \diff X_s + [X,Y]_t.
    $$
}
\pf{
    Assume IBP holds.
    
    Let $$\mathcal{A} = \left\{f \in C^2 (\mathbb{R}^n) \colon \diff f(X_t) = \sum_i D_i f(X_t) \diff X_t + \frac{1}{2} \sum_{i,j}D_{ij}f(X_t) \diff [X^i,X^j]_t\right\}.$$

    (i). $f,g \in \mathcal{A} \implies f+ g \in \mathcal{A}$.
    (ii). $f,g \in \mathcal{A} \implies fg \in \mathcal{ A}$.

    $$
    \begin{aligned}
    \diff (fg(X_t)) =& \diff (f(X_t)\cdot g(X_t)) \\
    =& f(X_t) \cdot \diff g(X_t) + g(X_t) \diff f(X_t) + \diff [f(X_t),g(X_t)]_t\\
    =& f(X_t)\left( \sum_i D_i g(X_t) \diff X_t + \frac{1}{2} \sum_{i,j} D_{ij}g(X_t) \diff [X^i,X^j]_t \right) \\
    &+ g(X_t)\left( \sum_i D_i f(X_t)\diff X_t + \frac{1}{2}\sum_{i,j}D_{ij}f(X_t) \diff [X^i,X^j]_t \right) \\
    &+ \diff \left[ \int_0^t \sum_{i} D_i f(X_t) \diff X_t, \int_0^t \sum_j D_j g(X_t) \diff X_t^j \right]\\
    =&\diff X_t \left( \sum_i f D_i g(X_t) + \sum_i g D_i f(X_t) \right) \\
    &+ \frac{1}{2} \sum_{i,j} \left( g D_{ij}f (X_t) + f D_{ij}g(X_t) + 2D_i fD_j g(X_t) \right)\diff [X^i,X^j]_t\\
    =& \sum_i D_i(fg)(X_t) \diff X_t + \frac{1}{2}\sum_{i,j}D_{ij}(fg)(X_t)\diff [X^i,X^j]_t\\
    \implies \quad& fg \in \mathcal{A}.
    \end{aligned}
    $$

    (iii). $f(x) = x \in \mathcal{A} \implies \forall~$polynomial $\in \mathcal{A}$.
}

\lemnn{$\forall~ f \in C^2(\mathbb{R}^d)$, $\exists~$polynomials $\{P_n(x)\}_{n \geq 1}$ such that 
    $$
    \begin{cases}
        P_n(x) \to f(x)\\
        D_i P_n(x) \to D_i f(x)\\
        D_{ij}P_n(x) \to D_{ij}f(x)
    \end{cases}
    $$
    uniformly on compacts.
}

Define $U_N = \inf\{t \geq 0\colon \abs{ X_t }+ [M]_{t} > N\}$ for each $N >0$.

Then $U_N \uparrow \infty$. So it suffices to prove It\^o formula holds for $\forall~ t \leq U_N$.

Note that $$P_n(X_t) = P_n(X_0) + \int_{0}^{t}\sum_{i} D_i P_n(X_s) \diff  X_s + \frac{1}{2} \int_{0}^{t} \sum_{i,j} D_{ij}P_n(X_s) \diff [X^i,X^j]_s$$ for $0 \leq t \leq U_N$.

By letting $n \to \infty$, we will prove the limits are 
$$
f(X_t) = f(X_0) + \int_{0}^{t}\sum_{i} D_i f(X_s) \diff X_s + \frac{1}{2} \int_{0}^{t} \sum_{ij}D_{ij}f(X_s) \diff [X^i,X^j]_s.
$$

$$
\begin{aligned}
    I_{1,n}=&\int_{0}^{t} \sum_i D_i P_n(X_s) \diff X_s \\
    =& \int_0^t \sum_i D_i P_n(X_s) \diff M_s + \int_0^t \sum_i D_i P_n(X_s) \diff A_s.
\end{aligned}
$$

Notice that 
$$
    \abs{ \int_0^t \sum_i D_i P_n (X_s) \diff A_s - \int_0^t \sum_i D_i f(X_s) \diff  A_s } \leq d\cdot \sup_{i} \underbrace{\norm{ D_iP_n - D_i f }}_{\text{max on compacts}}\cdot A_t \xlongrightarrow{n \to \infty} 0 \text{ a.s.}
$$

Next,
$$
    \begin{aligned}
    &\mathbb{E} \sup_{s \leq t} \abs{ \int_0^t \sum_i \left( D_i P_n(X_s) - D_i f(X_s) \right) \diff  M_s }^2 \\
    \leq & \mathbb{E}\int_0^t \left( \sum_i D_i P_n(X_s) - D_i f(X_s) \right)^2 \diff [M]_s \\
    \leq & d \cdot \sup_{1 \leq i \leq d}\norm{ D_iP_n - D_i f }^2 \cdot \underbrace{\mathbb{E}[M]_{t \wedge U_N}}_{N} \xlongrightarrow{n \to \infty} 0.
    \end{aligned}
$$

By taking subsequence, $\int_0^t \sum_i D_i P_{n_k}(X_s) \diff  M_s \xlongleftarrow{\text{a.s.}}\int_0^t \sum_i D_i f(X_s) \diff  M_s.$

$$
    \begin{aligned}
        I_{2,n} = \frac{1}{2} \int_0^t \sum_{i,j} D_{ij}P_n(X_s) \diff [X^i,X^j]_s.
    \end{aligned}
$$

Notice that 

$$
    \abs{ I_{2,n} - \frac{1}{2}\int_0^t \sum_{i,j}D_{ij}f(X_s) \diff [X^i,X^j]_s } \leq C \sum_{i,j} \norm{ D_{ij}P_n - D_{ij}f }\cdot \underbrace{\sum_{i,j}[X^i,X^j]_t}_{\leq C_N} \xlongrightarrow{n \to \infty} \to 0.
$$

\lemnn{
    Let $M$ be a bounded continuous martingale with $M_0 = 0$. Let $V$ be continuous $FV_0$
    $$
    \begin{aligned}
        &M_t^2 = \int_0^t 2M_s \diff M_s + [M]_t ,\\
        &M_t V_t = \int_0^t M_s \diff V_s + \int_0^t V_s \diff  M_s.
    \end{aligned}
    $$
}
\pf{
    $$
    \begin{aligned}
        M_t^2 &= \left[ \sum_{k=1}^{2^n}M(t_{k}^{n}) - M(t_{k-1}^{n}) \right]^2\\
        &= \sum_{k=1}^{2^n}\left( M(t_{k}^{n}) - M(t_{k-1}^{n}) \right)^2 + 2\sum_{k=1}^{2^n}\left[ M(t_{k}^{n})M(t_{k-1}^{n}) - M(t_{k-1}^{n})^2 \right]\\
        &=M(t_{2^n}^{n})^2 - M(t_{0}^{n})^2 = M(t)^2
    \end{aligned}
    $$
    where $t_{k}^{n} = \frac{k}{2^n} t$.

    Note that $\sum_{k=1}^{2^n}\left( M(t_{k}^{n}) - M(t_{k-1}^{n}) \right)^2 \to [M]_t$ and 
    $$
    2\sum_{k=1}^{2^n}\left[ M(t_{k}^{n})M(t_{k-1}^{n}) - M(t_{k-1}^{n})^2 \right] = 2 \sum_{k=1}^{2^n}M(t_{k-1}^{n})(M(t_{k}^{n}) - M(t_{k-1}^{n}))\to 2\int_0^t M_s \diff M_s.
    $$

    $$
    \begin{aligned}
        M_t V_t &= M_{t_{2^n}^n}V_{t_{2^n}^n} - M_{t_{0}^n}V_{t_0^n} \\
        &= \sum_{k=1}^{2^n}\left[ M(t_{k}^{n})V(t_{k}^{n}) - M(t_{k-1}^{n})V(t_{k-1}^{n}) \right] \\
        &= \sum_{k=1}^{2^n} M(t_{k}^{n})(V(t_{k}^{n}) - V(t_{k-1}^{n})) + \sum_{k=1}^n V(t_{k-1}^{n})\left( M(t_{k}^{n}) - M(t_{k-1}^{n}) \right) \\
        &\to \int_{0}^{t}M_s \diff V_s + \int_0^t V_s \diff M_s.
    \end{aligned}
    $$
}

\thm{L\'evy Thm}{
    Let $\{B_t^i\}_{t \geq 0} \in c \mathcal{M}_{0,loc}$ for $i=1,2,\cdots,n$. If $\forall~i,j$, we have 
    $$ 
    [B^i,B^j]_t = \begin{cases}
        t, & i = j,\\
        0, & i \neq j.
    \end{cases}
    $$
    
    Then $B_t = (B_t^1,\cdots,B_t^n)$ is a Brownian motion on $\mathbb{R}^n$ with $\mathcal{N}(0,I_n).$
}

\pf{
    $\forall~ s < t,$ we will prove 
    $$
    \mathbb{E}\left[ e^{i \theta (B_t - B_s)} \middle| \mathcal{F}_s \right] = e^{-\frac{1}{2}(t-s) \abs{ \theta }^2}
    $$
    for all $\theta \in \mathbb{R}^n$.

    Define $f(x,t) = e^{i\theta x + \frac{1}{2}\abs{ \theta }^2 t }$, $\mathbb{R}^{n+1} \to \mathbb{C}$.

    Apply It\^o's formula to $f(B_t,t) \equiv f(X_t).$

    $X_t = (B_t^1,B_t^2,\cdots,B_t^n,t).$

    $$
    \begin{aligned}
    f(X_t) =& f(x_0) + \int_0^t \sum_{i=1}^n (D_i f(X_s \diff B_s^i)) + \int_0^1 D_{n+1}f(X_s)ds \\
    &+ \frac{1}{2}\int_0^t \sum_{1 \leq i,j \leq n}D_{ij}f(X_s) \diff [B^i,B^j]_s \\
    &+ \frac{1}{2}\int_0^t \sum_{1 \leq k \leq n}D_{k,n+1}f(X_s) \diff [B^k,s]_s
    \end{aligned}
    $$
    $$
    \begin{aligned}
        \implies f(B_t,t)= &f(B_0,0) + \int_0^t \sum_i D_i f(B_s,s)\diff B_s^i \\
        &+ \int_0^t \left( e^{i \theta B_s + \frac{1}{2}\abs{ \theta }^2 s} \cdot \frac{1}{2} \abs{ \theta }^2 \right) \diff s \\
        &+ \int_0^t \frac{1}{2}\left( \sum_{j=1}^n e^{i \theta B_s + \frac{1}{2}\abs{ \theta }^2 s }(i \theta_j)^2 \right) \diff s.
    \end{aligned}
    $$

    We conclude that $f(B_t,t) = f(B_0,0) + c \mathcal{M}_{0,loc}$
    $$
    \implies e\underbrace{^{i\theta B_t + \frac{1}{2} \abs{ \theta }^2 t}}_{M_t^\theta \in c \mathcal{M}_{0,loc}} = 1 + c \mathcal{M}_{0,loc}
    $$
    Since $\abs{ M_t^\theta } = e^{\frac{1}{2} \abs{ \theta }^2 t} < \infty$ and 
    $$
    \mathbb{E}(M_t^{T_n} \mid \mathcal{F}_s) = M_s^{T_n}
    $$
    where $T_n$ reduces $M$.

    We may take $n \to \infty$ to get $\mathbb{E}(M_t \mid \mathcal{F}_s) = M_s$.

    $$
    \implies \mathbb{E}(e^{i \theta B_t + \frac{1}{2}\abs{ \theta }^2 t} \mid \mathcal{F}_s) = e^{i\theta B_s + \frac{1}{2} \abs{ \theta }^2 s}.
    $$

    $$
    \implies \mathbb{E}\left( e^{i \theta(B_t - B_s)} \mid \mathcal{F}_s \right) = e^{\frac{1}{2}\abs{ \theta }^2 (t-s)}.
    $$

    $$
    \implies B_t - B_s \sim \mathcal{N}(0,t-s)
    $$

    Q independent of $\mathcal{F}_s$.
}

\thm{Dubins and Schwarz}{
    let $M \in c \mathcal{M}_{0,loc}$ such that $[M]_t \uparrow \infty$ as $t \to \infty$. Then $M_t = B([M]_t)$ where $B$ is some Brownian motion.

    Let $\tau_t = \inf\{u \geq 0 \colon [M]_u > t\}$. Set $\mathcal{G}_t \equiv \mathcal{F}(\tau_t)$. Then $B_t \equiv M(\tau_t)$ is a B.M. with respect to $\mathcal{G}_t.$
}

\pf{
    (i) $M(\tau_t) \in \mathcal{F}_{\tau_t} \implies B_t \in \mathcal{G}_t$.
    
    (ii) $t\mapsto B_t$ is continuous.

    $$
    \begin{cases}
        \tau_{a-} = b, \\
        \tau_{a+} = c.
    \end{cases}
    $$

    We only need to prove $M(b) = M(c)$.

    Assume $M$ is bounded $\forall~ q \in \mathbb{Q}_+$, define 
    $$
    S_q \equiv \inf\{t > q \colon [M]_t > [M]_q\}.
    $$

    We will show $M_t \equiv M_q$ for all $t \in [q,S_q)$
    
    Note $S_q \geq q$, by OST,
    $$\mathbb{E}\left( M_{S_q}^2 - [M]_{S_q} \mid \mathcal{F}_q \right) = M_q^2 - [M]_q.
    $$

    Because $[M]_{S_q} = [M]_q$,
    $$
    \implies \mathbb{E}[M_{S_q}^2 - M_q^2 \mid \mathcal{F}_q] = 0.
    $$
    $$
    \implies \mathbb{E}((M_{S_q} - M_q)^2 \mid \mathcal{F}_q) = 0.
    $$
    $$
    \implies \mathbb{E}(M_{S_q} - M_q)^2 = 0.
    $$
    $$
    \implies M_{S_q} = M_q.
    $$

    (iii) $B_t$ and $B_t^2 - t \in c \mathcal{M}_{0,loc} \implies [B]_t = t \implies$ B.M.

    Recall $B_t = M(\tau_t)$. Define 
    $$
    \begin{aligned}
    &T_n = \inf \{ t \geq 0 \colon \abs{ M_t } > n\},\\
    &U_n = [M]_{T_n}.    
    \end{aligned}
    $$

    We will check $\tau_{t \wedge U_n} = T_n \wedge \tau_t.$
    $$
    t \wedge U_n \leq t \implies \tau_{t \wedge U_n } \leq \tau_t.
    $$
    $$
    \tau_{t\wedge U_n} =\inf \{
        u\colon [M]_u > t \wedge U_n = t \wedge[M]_{T_n}
    \}.
    $$

    \textbf{Claim:} $\tau_t \wedge T_n = \tau_{t \wedge U_n}$
    
    \textbf{case 1.} $t \leq U_{n}$ we have
    $$
        \tau_{t \wedge U_{n}} = \inf \{u\colon [M]_{u}> t\} = \tau_{t},
    $$
    and
    $$
    t \leq U_{n} = [M]_{T_{n}} \implies T_{n} \geq \tau_{t} \implies \tau_{t} \wedge T_{n} = \tau_{t} = \tau_{t \wedge U_{n}}.
    $$

    \textbf{case 2.} $t > U_{n}$ we have
    $$
    \tau_{t \wedge U_{n}} = \tau_{U_{n}} = \inf \{u\colon [M]_{u} > U_{n}\}
    $$
    and
    $$
    t > U_{n} > [M]_{T_{n}} \implies T_{n} < \tau_{t} \implies \tau_{t} \wedge T_{n} = T_{n}.
    $$
    It suffices to prove $T_{n} = \tau_{[M]_{T_{n}}}= \tau_{U_{n}}$.

    One hand, $\tau_{U_{n}} \geq T_{n}$ can be proved by def of $\tau_{U_{n}}$.

    On the other hand, $T_{n} = \inf \{t \colon \abs{ M_{t} } > n\}$. Take $\forall~\varepsilon>0$, $u_{\varepsilon} \in (T_{n} , T_{n} + \varepsilon)$ s.t. $\abs{ M_{u_{\varepsilon}} }>n .$

    Hence $M_{u_{\varepsilon}} \neq M_{T_{n}}$. Suppose $[M]_{u_{\varepsilon}} = [M]_{T_{n}}$, then $M$ is constant by the property of continuous local martingale. Contradiction. So $[M]_{u_{\varepsilon}} > [M]_{T_{n}} = U_{n} \implies \tau_{U_{n}} \leq u_{\varepsilon}$. Let $\varepsilon \downarrow 0$ to get $\tau_{U_{n}} \leq T_{n}$.

    Then $B_{t \wedge U_n} = M(\tau_{t \wedge U_n}) = M(T_n \wedge \tau _t)$, $B_t^{U_n} = M_{\tau_t}^{T_n}.$

    By OST 
    $$
    \mathbb{E}\left( B_t^{U_n} \middle| \mathcal{F}_{\tau_s} \right) = \mathbb{E}\left( M_{\tau_t}^{T_n} \middle| \mathcal{F}_{\tau_s} \right) = M_{\tau_s}^{T_n} = B_s^{U_n}.
    $$

    Then $B^{U_n}$ is a martingale w.r.t. $\mathcal{G}_t = \mathcal{F}_{\tau_t}$ + (2) $U_n$ is $\mathcal{G}_t$-stopping time $\implies B \in c \mathcal{M}_{0,loc}$

    \textbf{pf of (2)}

    $$
    \{[M]_{T_n} \leq t\}=\{U_n \leq t\} \in \mathcal{F}_{\tau_t}.
    $$

    $$
    \{T_n \leq \tau_t\} \in \mathcal{F}_{\tau_t}.
    $$

    $\left( M^2 - [M] \right)^{T_n}$ is martingale. Then 
    $$
    \mathbb{E}\left[ (M_{\tau_t}^{T_n})^2 - [M]_{\tau_t}^{T_n} \middle|  \mathcal{F}_{\tau_s}\right]= \left( M_{\tau_s}^{T_n} \right)^2 - [M]_{\tau_s}^{T_n}
    $$  
    $$
    \implies \mathbb{E}\left[ 
        (B_t^{U_n})^2 - t \wedge U_n \middle| \mathcal{F}_s 
     \right] = (B_s^{U_n})^2 - s \wedge U_n.
    $$  
    $
    \implies (B_t^2 - t)^{U_n}
    $ is a martingale $\implies B_t^2 - t \in \mathcal{M}_{loc} \implies[B]_t = t.$
}   

If $[M]_{\infty} < \infty$, we get $\tau_t = \infty$, $\forall~ t \geq [M]_{\infty}$.
$$
B_t = M(\infty), \quad \forall~ t \geq [M]_{\infty}.
$$

So this def does not give a B.M.

Define 
$$B_t = 
\begin{cases}
    M(\tau_t), & t < [M]_{\infty},\\
    M(\infty)+W(t) - W([M]_{\infty}), & t \geq [M]_{\infty},
\end{cases}
$$
where $W$ is a standard B.M. independent of $M$.

Or 
$$B_t
\begin{cases}
    M(\tau_t),& t <[M]_{\infty}\\
    M(\infty)+W(t-[M]_{\infty}), & t \geq[M]_{\infty}.
\end{cases}
$$
Then $B([M]_{t})= M_t.$

\lemnn{
    On $\{[M]_{\infty
} < \infty\}$, the limit $M_\infty \equiv \lim_{t \uparrow \infty} M_t$ exists a.s.
}
\pf{
    Define $S_K = \inf \{t \geq 0 \colon [M]_t > K\}$. Then on $\{[M]_{\infty} < \infty\}$, $S_K = \infty$, $\forall~ K >[M]_\infty$

    $M^{S_K}$ is a $L^2$-bounded martingale.

    $\implies \lim_{t \uparrow \infty} M_t^{S_K}$ exists a.s.

    Let $K \uparrow \infty$ to get $\lim_{t\uparrow \infty}M_t$ exists a.s.

    By an enlargement of the prob. space $(\Omega. \mathcal{F},\mathsf{P})$
    $$
    \begin{aligned}
    \implies &\tilde{\Omega} = \Omega \times \Omega^*\\
        & \tilde{\mathcal{F}}_t = \mathcal{F}_t \times \mathcal{F}_{t}^* \\
        & \tilde{\mathsf{P}} = \mathsf{P} \times \mathsf{P}^*
    \end{aligned}
    $$
    where on $(\Omega^*, \mathcal{F}^*,\mathcal{F}_{t}^*, \mathsf{P}^*)$ we have $W^*$ is B.M. independent of $M$.
}

\thm{Dubins-Schwarz (general version)}{
    Let $M \in c \mathcal{M}_{0,loc}$. Then $\exists~$B.M. B s.t. $M_t = B([M]_t)$ for all $t \geq 0$.
}

\thm{Dol\'eans' Thm}{
    Let $X$ be a semimartingale with $X_0 = 0$. Suppose $Z_0 \in \mathcal{F}_0$. Then $\exists~$a unique continuous semimartingale $Z$ s.t. 
    $$
    Z_t = Z_0 + \int_{0}^{t}Z_s \diff X_s.
    $$
    The solution $Z_t = Z_0 e^{X_t - \frac{1}{2}[X]_t}$. 

    \textbf{Notation:} $\mathcal{E}(X_t) = e^{X_t - \frac{1}{2}[X]_t}$ is called the exponential semimartingale of $X$.
}

\pf{
    $\mathcal{E}(X_t) = e^{X_t - \frac{1}{2}[X]_t} = f(X_t, [X]_t)$ with $f(x,y) = e^{x - \frac{1}{2}y}.$

    By apply It\^o's formula, 
    $$
    \diff f(X_t,[X]_t) = \partial_x f(X_t,[X]_t) \diff X_t + \partial_y f(X_t, [X]_t) \diff [X]_t + \frac{1}{2}\left( \partial_{xx}f(X_t,[X]_t) \diff [X]_t \right).
    $$
    Because 
    $$
    \begin{cases}
        \diff X_t \diff X_t = \diff [X]_t\\
        \diff X_t \diff [X]_t = 0\\
        \diff [X]_t \diff [X]_t = 0.
    \end{cases}
    $$
    Hence 
    $$
    \diff \mathcal{E}(X_t) = \mathcal{E}(X_t) \diff X_t.
    $$
    So $\mathcal{E}(X)$ is a solution to 
    $$
    \diff Z_t = Z_t \diff X_t.
    $$

    \textbf{(Uniqueness)} Assume $Z$ is another solution. We will prove 
    $$
        Z_t = Z_0 \mathcal{E}(X_t).
    $$

    We now have $\diff Z_t = Z_t \diff X_t$. 
    
    Define 
    $$
    Y_t = e^{- X_t + \frac{1}{2}[X]_t} = (\mathcal{E}(X)_t)^{-1}.
    $$

    It suffices to prove $\diff (Z_t Y_t) = 0$ which implies 
    $$
    Z_t Y_t = Z_0 Y_0 = Z_0, \quad \forall~t.
    $$
    
    We have
    $$
    \diff (Z_t Y_t) = Z_t \diff Y_t + Y_t \diff Z_t + \diff [Y,Z]_t.
    $$

    Notice that by It\^o's formula,
    $$
    \diff Y_t = (-1)Y_t \diff X_t + \frac{1}{2} Y_t \diff [X]_t + \frac{1}{2}(-1)^2 Y_t \diff [X]_t = -Y_t \diff X_t + Y_t \diff [X]_t.
    $$

    $$
    \implies \diff (Z_t Y_t) = Z_t(-Y_t \diff X_t + Y_t \diff [X]_t) + Y_t(Z_t \diff X_t) + Y_t(Z_t \diff X_t) + (-Y_t Z_t) \diff [X]_t = 0.
    $$
}

Let $M \in c \mathcal{M}_{0,loc}$. Then $\mathcal{E}(\theta M)_t = e^{\theta M_t - \frac{1}{2}\theta^2[M]_t}$ is in $c \mathcal{M}_{loc}$.

By Fatou's lemma, we get 
$$
\begin{aligned}
    \mathbb{E}(\mathcal{E}(\theta M)_t \mid \mathcal{F}_s) &= \mathbb{E}(\lim_{n\to \infty} \mathcal{E}(\theta M)_{t \wedge T_n} \mid \mathcal{F}_s) \\
    &\leq \liminf_{n \to \infty} \mathbb{E}(\mathcal{E}(\theta M)_{t}^{T_n} \mid \mathcal{F}_s)\\
    & = \liminf_{n \to \infty} \mathcal{E}(\theta M)_{s}^{T_n} \\
    &= \mathcal{E}(\theta M)_s.
\end{aligned}
$$

Hence $\mathcal{E}(X_t)$ is a supermartingale and $\mathbb{E}(\mathcal{E}(M)_t) \leq 1$.

\thmnn{
    Suppose $\forall~ t > 0$, $\exists~ K_t > 0$ s.t.
    $$
    [M]_t < K_t \quad \text{a.s.}
    $$

    If for some $t > 0$, $\exists~K_t >0$ s.t.
    $$
    [M]_t < K_t \quad \text{a.s.}
    $$
    
    Then for this $t > 0$,
    $$
    \mathsf{P}\left( \max_{s \leq t}M_s > y \right) \leq e^{-\frac{y^2}{2K_t}}, \quad \forall~ y > 0.
    $$
    and $\mathcal{E}(\theta M)$ is a true martingale.
}

\pf{
    By using $\mathcal{E}(\theta M)$ is a supermartingale, by Doob's inequality,
    $$
    \begin{aligned}
    \mathsf{P}\left( \max_{s \leq t} M_s > y \right) &= \mathsf{P}\left( \max_{s \leq t} \mathcal{E}(\theta M)_s > e^{\theta y - \frac{1}{2} \theta^2 K_t} \right)\\
    & \leq \frac{\mathbb{E}(\mathcal{E}(\theta M)_t)}{e^{\theta y -\frac{1}{2} \theta^2 K_t}}\\
    & \leq e^{-\theta y + \frac{1}{2} \theta^2 K_t}.
    \end{aligned}
    $$
    Pick $\theta = \frac{y}{K_t}$ to get $e^{-\frac{y^2}{2K_t}}$.

    Next, define $Z_t = \max_{s \leq t} M_s$. Then $\forall~ \theta > 0$, 
    $$
        \mathbb{E}e^{\theta Z_t} = \int_{\mathbb{R}}e^{\theta y} \mathsf{P}(Z_t \in \diff y).
    $$

    Notice that $\forall~ z \in \mathbb{R}$, 
    $$
        e^{\theta Z} = 1 + \int_{0}^{Z} \theta e^{\theta y} \diff y.
    $$
    $$
    \begin{aligned}
    \implies 
        \mathbb{E}(e^{\theta Z_t}) &= 1 + \mathbb{E}\left( \int_0^{Z_t} \theta e^{\theta y} \diff y \right) \\
        &= 1 + \mathbb{E}\left( \int_0^\infty \mathsf{1}_{y < Z_t} \theta e^{\theta y} \diff  y \right)\\
        & = 1+ \int_0^\infty \theta e^{\theta y} \diff  y  \mathsf{P}(Z_t > y)\\
        & \leq 1 + \int_0^\infty \theta e^{\theta y} e^{-\frac{y^2}{2K_t}} \diff y \\
        &< \infty.
    \end{aligned}
    $$

    By $\mathcal{E}(\theta M)_t$ is a $\mathcal{M}_{loc}$, we get $\exists~T_n \uparrow \infty$ s.t. $\mathcal{E}(\theta M)_{t}^{T_n}$ is a martingale.

    $$
        \implies \forall~ s < t,\quad \mathbb{E}[\mathcal{E}(\theta M)_{t}^{T_n} \mid \mathcal{F}_s] = \mathcal{E}(\theta M)_{s}^{T_n}.
    $$

    By DCT for $\mathcal{E}(\theta M)_{t}^{T_n} \leq Z_t$ and $\mathbb{E} Z_t < \infty$, we get 
    $$
    \lim_{n \to \infty} \mathbb{E}(\mathcal{E}(\theta M)_{t}^{T_n} \mid \mathcal{F}_s) = \mathbb{E}(\mathcal{E}(\theta M)_t \mid \mathcal{F}_s).
    $$
}

\cor{
    Let $B$ be a B.M. in $\mathbb{R}^n$ and let $H$ be a bounded previsible process. Then 
    $$
    e^{\int H \diff B - \frac{1}{2}\int \abs{ H_s }^2 \diff s}
    $$
    is a true martingale.
}

\rmkb{$[M]_t = \int_0^t \abs{ H_s }^2 \diff  s < K_t$.}

\cor{
    Let $M \in c \mathcal{M}_{0,loc}$. Then $\forall~ t > 0$, $y> 0$, $K >0$,
    $$
    \mathsf{P}\left( \max_{s \leq t} M_s \geq y, [M]_{t} \leq K \right) \leq e^{-\frac{y^2}{2K}}.
    $$
    $$
    \implies \mathsf{P}\left( \max_{s \leq t}\abs{ M_s } \geq y, [M]_t \leq K \right) \leq 2 e^{-\frac{y^2}{2K}}.
    $$
}

\pf{
    Let $T =\inf \{u \geq 0 \colon [M]_{u} \geq K\}$. Then 
    $$
    \mathsf{P}\left( \max_{s \leq t}M_s \geq y, [M]_{t} \leq K \right) \leq \mathsf{P}\left( \max_{s \leq t} M_{s \wedge T} \geq y \right) \leq e^{- \frac{y^2}{2K}}
    $$  
    since $[M^T]_t \leq K$.
}

\thm{Novikov's Thm}{
    If $\forall~t > 0$, $\mathbb{E}(e^{\frac{1}{2} [M]_t} ) < \infty$ then $\mathcal{E}(X)_t$ is a true martingale.
}

\pf{
    See P198 of [KS].
}

\subsection{Burkholder-Davis-Gundy inequalities}

\defnn{
    We say a function $F: \mathbb{R}^+ \to \mathbb{R}^+$ is \underline{moderate} if continuous and increasing, $F(x) = 0 \Leftrightarrow x=0$, and $\forall~\alpha > 1$,
    $$
        \sup_{x > 0} \frac{F(\alpha x)}{F(x)} < \infty.
    $$
}
\ex{}{
    $F(x) = x^{p}$, $\forall~ p > 0$.
    $F(x) = x^p \left( \log x \vee 1 \right)$, $\forall~ p \geq 0$.
}

\thm{BDG inequality}{
    Let $F$ be moderate. Then $\exists~$universal constants $c_F \leq C_F < \infty$ s.t. $\forall~ M \in c \mathcal{M}_{0,loc}$,
    $$
    c_{F} \mathbb{E}\left( F([M]_{\infty}^{\frac{1}{2}}) \right) \leq \mathbb{E} F(M_\infty^*) \leq C_F \mathbb{E}\left( F([M]_{\infty}^{\frac{1}{2}}) \right).
    $$

    $M_t^* = \sup_{s \leq t}\abs{ M_s }.$ Fix $t > 0$. Let $\tilde{M}_{s}^{t} = M_{s \wedge t}$. Then 
    $$
    c_{F} \mathbb{E}\left( F([M]_{t}^{\frac{1}{2}}) \right) \leq \mathbb{E}\left( F(M_t^*) \right) \leq C_F \mathbb{E} F([M]_{t}^{\frac{1}{2}}).
    $$
}

\ex{}{
    If $F(x) = x^p$, then 
    $$
    c_p \mathbb{E}[M]_{t}^{\frac{p}{2}} \leq \mathbb{E} \abs{ M_t^* }^p \leq C_p \mathbb{E}[M]_t^{\frac{p}{2}}.
    $$
}

\lemp{``good'' $\lambda$ inequality}{
    Let $X$ and $Y$ be nonegative R.V.s. Suppose $\exists~\beta > 1$ s.t. $\forall~ \lambda > 0$, $\delta >0$,
    $$
    \mathsf{P}(X > \beta \lambda, Y \leq \delta \lambda) \leq \psi(\delta) \mathsf{P}(X > \lambda)
    $$
    where $\psi(\delta) \downarrow 0$ as $\delta \downarrow 0$.

    Then for each moderate $F$, $\exists~C=C(\beta,\psi, F) > 0$ s.t. 
    $$
    \mathbb{E}[F(X)] \leq c \mathbb{E}[F(Y)].
    $$
}{
    Assume good $\lambda$ inequality. Then it holds by replacing $X$ by $X \wedge a$ for any $a>0$.

    If $a > \beta \lambda > \lambda$, then
    $$
    \begin{aligned}
    &X \wedge a > \beta \lambda \Leftrightarrow X > \beta \lambda\\
    & X \wedge a > \lambda \Leftrightarrow X > \lambda.
    \end{aligned}
    $$

    Hence we may assume $\mathbb{E}F(X) < \infty$.

    [ $\mathbb{E}F(X\wedge a) \leq \mathbb{E} F(Y)$ for all $a >0$. Then let $a \uparrow \infty$.]

    Notice that $\sup_{x>0} \frac{F(x)}{F(\frac{x}{\beta})}$.

    Let $\gamma > 0$ s.t. $F(x) \leq \frac{1}{\gamma} F(\frac{x}{\beta})$ for all $x > 0$.

    Now 
    $$
    \int F(\diff \lambda) \psi(\delta) \mathsf{P}(X > \lambda) \geq \int F(\diff \lambda) \mathsf{P}(X > \beta \lambda, Y \leq \delta \lambda).
    $$
    $$
    LHS = \psi(\delta) \mathbb{E}\left( \int_0^X F(\diff \lambda) \right) = \psi(\delta)\mathbb{E}(F(X)).
    $$
    $$
    \begin{aligned}
    RHS &= \int F(\diff \lambda) \mathbb{E}(\mathsf{1}_{\frac{Y}{\delta} \leq \lambda < \frac{X}{\beta}}) \\
    &= \mathbb{E}\left( \int_{\frac{Y}{\delta}}^{\frac{X}{\beta}} F(\diff \lambda) \right) \\
    &= \mathbb{E}\left( F\left(\frac{X}{\beta}\right) - F\left(\frac{Y}{\delta}\right) \right)\\
    & \geq \mathbb{E}\left( \gamma F(X) \right) - \mathbb{E}\left( F\left( \frac{Y}{\delta} \right) \right).
    \end{aligned}
    $$

    Let $\delta > 0$ be small sucht that 
    $$
    \psi(\delta) < \frac{\gamma}{2} \implies \gamma -\psi(\delta) \geq \frac{\gamma}{2}.
    $$

    Hence 
    $$
    \mathbb{E}\left( F(X) \right)\left( \psi(\delta) - \gamma \right) \geq - \mathbb{E}\left( F\left( \frac{Y}{\delta} \right) \right).
    $$

    $$ \implies
    \mathbb{E}\left( F\left( \frac{Y}{\delta} \right) \right) \geq \left( \gamma - \psi(\delta) \right) \mathbb{E}F(X) \geq \frac{\lambda}{2} \mathbb{E}F(X).
    $$

    Using $\displaystyle\sup_{y > 0}\frac{F(y/\delta)}{F(y)} < \infty$, $\implies \exists~\mu > 0$ s.t. $F(y/\delta) \leq \mu F(y)$ for all $y > 0$ 
    $$\implies \mathbb{E}(F(Y/\delta)) \leq \mu \mathbb{E} F(Y)$$
    $$
    \implies \frac{\gamma}{2}\mathbb{E}F(X) \leq \mu \mathbb{E}F(Y).
    $$
}

\pf{
    of BDF inequality.

    $$\mathsf{P}\left(M_\infty^* > \beta \lambda , [M]_{\infty}^{\frac{1}{2}} \leq \delta \lambda\right) \leq e^{-\frac{(\beta \lambda)^2}{2 (\delta \lambda)^2}} = e^{-\frac{\beta^2}{2 \delta^2}}.$$

    Define $\tau = \inf \{u \geq 0 \colon \abs{ M_u }> \lambda\}$. 

    Then 
    $$\begin{aligned}
        \mathsf{P}\left( M_\infty^* > \beta \lambda, [M]_{\infty}^{\frac{1}{2}} \leq \delta \lambda \right) &= \mathsf{P}\left( M_\infty^* > \beta \lambda, [M]_{\infty}^{\frac{1}{2}} \leq \delta \lambda , \tau < \infty\right)\\
        &= \mathbb{E}\left[ \mathsf{1}_{\tau < \infty} \mathsf{P}\left( M_\infty^* > \beta \lambda, [M]_{\infty}^{\frac{1}{2}} \leq \delta \lambda \,\middle|\, \mathcal{F}_{\tau} \right) \right]
    \end{aligned}
    $$

    Define $N_t \equiv (M_{t + \tau} - M_{\tau})^2 - ([M]_{t + \tau} - [M]_{\tau})$.

    Then $N_t \in \mathcal{F}_{t + \tau}$ is a continuous local martingale
    $$
    \mathbb{E}\left[ M_{t+\tau}^2 - [M]_{t + \tau} \mid \mathcal{F}_{s + \tau} \right] = M_{s+\tau}^2 - [M]_{s+\tau}.
    $$

    $ \implies N_0 = 0$ and on $\{M_{\infty}^* > \beta \lambda, [M]_{\infty}^{\frac{1}{2}} \leq \delta \lambda\}$ we have $\exists~\tilde{t} > 0$ s.t.
    $$
    M_{\tilde{t} + \tau} \geq \beta \lambda \implies N_{\tilde{t}} = \left( M_{\tilde{t}+\tau} \pm \lambda \right)^2 - \left( [M]_{\tilde{t} + \tau} - [M]_{\tau} \right) \geq \left[ (\beta - 1) \lambda \right]^2 - (\delta \lambda)^2.
    $$

    Moreover $N_t \geq - \delta^2 \lambda^2$, $\forall~ t \geq 0$. 

    Hence $N_t$ reaches $[(\beta-1)\lambda]^2 - (\delta \lambda)^2$ before $(-\delta^2 \lambda^2)$.

    By (Gambler's Ruin $\mathsf{P}_{k}(\tau_0 < \tau_N) = \frac{N-k}{N}$)
    $$ 
    \mathsf{P}(\tau_b < \tau_a) \leq \frac{-a}{b-a}
    $$
    $$
    \implies \mathsf{P}\left( M_\infty^* > \beta \lambda, [M]_{\infty}^{\frac{1}{2}} \leq \delta \lambda \mid \mathcal{F}_\tau \right) \leq \mathsf{P}\left( \tau_{[(\beta - 1)\lambda]^2 - (\delta \lambda)^2} < \tau_{-(\delta \lambda)^2} \right) \leq \frac{(\delta \lambda)^2}{[(\beta -1)\lambda]^2} = \frac{\delta^2}{(\beta -1)^2}.
    $$

    Here we take $0 < \delta < \beta -1$. Therefore
    $$
    \mathsf{P}\left( M_\infty^* > \beta \lambda, [M]_{\infty}^{\frac{1}{2}} \leq \delta \lambda  \right) \leq \delta^2 \frac{1}{(\beta - 1)^2} \mathsf{P}(\tau < \infty) = \delta^2 \frac{1}{(\beta-1)^2} \mathsf{P}\left( M_\infty^* > \lambda \right).
    $$

    Similarly, 
    $$
    \mathsf{P}\left( [M]_{\infty}^{\frac{1}{2}} > \beta \lambda, M_{\infty}^*  \leq \delta \lambda\right).
    $$

    $$
    \tilde{N}_t = [M]_{t+\tilde{\tau}} - [M]_{\tilde{\tau}} - (M_{\tilde{\tau} + t}-M_{\tilde{\tau}})^2
    $$
    with $\tilde{\tau} = \inf\{u \geq 0 \colon [M]_{u}^{\frac{1}{2}} > \lambda\}$.

    $$
    \implies \mathsf{P}\left( [M]_{\infty}^{\frac{1}{2}} > \beta \lambda, M_{\infty}^*  \leq \delta \lambda, \tilde{\tau} < \infty \right) = \mathbb{E}\left[ \mathsf{P}(\dots \mid \mathcal{F}_{\tilde{\tau}})\mathsf{1}_{\tilde{\tau} < \infty} \right].
    $$

    Then $\tilde{N}_t$ reaches $(\beta \lambda)^2 - \lambda^2 - (\delta \lambda)^2$ before $-(\delta \lambda)^2$.

}

\section{Semimartingale local time}
For a Brownian motion $(B_t)_{t \geq 0}$. We consider
$$
l_t^0 = \int_0^t \delta_0 (B_s) \diff s,
$$
where $\delta_0$ is the Dirac delta function such that $\forall~$function $f$,
$$
\int f(x) \delta_0(x) \diff x = f(0).
$$

$$
\mu([a,b]) = l_b^0 - l_a^0,
$$
$\mathrm{supp}~(\mu)=\{t \geq 0 \colon B_t = 0\}$, level set $=\{t \geq 0 \colon B_t = x\}$, using local time 
$$
l_t^x = \int_0^t \delta_x (B_s) \diff s.
$$

\thm{Tanaka's formula}{
    Let $X$ be a continuous semimartingale. Then $\exists~$continuous increasing adapted process $\{l_t, t \geq0\}$
    $$
    \abs{ X_t } - \abs{ X_0 } = \int_0^t \mathrm{sgn}(X_s) \diff X_s + l_t,
    $$
    where 
    $$
    \mathrm{sgn}(x) = \begin{cases}
        1, & x >0\\
        0, & x = 0\\
        -1,&x< 0.
    \end{cases}
    $$
    Then $\{l_t , t \geq 0\}$ is called the local time of $X$ at 0, and 
    $$\int_0^t \mathsf{1}_{\{X_s \neq 0\}} \diff l_s = 0.$$
}

\rmkb{
    \begin{itemize}
        \item[(1)] $f(x) = \abs{ x } \implies f'(x) = \mathrm{sgn}(x)$, $\forall~ x \neq 0 \implies f''(x) = 2 \delta(x).$ (distributional derivative)
        \item[(2)] $\abs{ X_t - x} - \abs{ X_0 - x } = \int_0^t \mathrm{sgn}(X_0 - x) \diff X_s + l_t^x.$
        \item[(3)] Note that $\abs{ X_t } = X_t^+ + X_t^-$, $X_t = X_t^+ - X_t^- $
        $$\implies 2X_t^+ = \abs{ X_t } + X_t = \int_0^t \mathsf{1}_{\{X_s > 0\}} \diff X_s - \int_0^t \mathsf{1}_{\{X_s \leq 0\}} \diff X_s + l_t + \int_0^t 1 \diff X_s + X_0 + \abs{ X_0 }$$
        $$
        \implies 2X_t^+ = 2 \int_0^t \mathsf{1}_{\{X_s > 0\}} \diff X_s + l_t + 2X_0^+.
        $$
        $$
        \implies X_t^+ - X_0^+ = \int_0^t \mathsf{1}_{\{X_s > 0\}} \diff X_s + \frac{1}{2} l_t.
        $$

        Similarly, $2X_t^- = \abs{ X_t } - X_t$
        $$
        \implies X_t^- - X_0^- = -\int_0^t \mathsf{1}_{\{X_s \leq 0\}} \diff X_s + \frac{1}{2} l_t.
        $$
    \end{itemize}
}

\pf{
    Apply It\^o's formula to the function $f(x) = \abs{ x }.$
    $$
    \begin{aligned}
    \implies \abs{ X_t }&= \abs{ X_0 } + \int_0^t f'(X_s) \diff X_s + \frac{1}{2}\int_0^t f''(X_s) \diff [X]_s\\
    &= \abs{ X_0 } + \int_0^t \mathrm{sgn}(X_s) \diff X_s + \frac{1}{2} \underbrace{\int_0^t 2 \delta(X_s) \diff [X]_s}_{=l_t}.
    \end{aligned}
    $$

    $\exists~ \varphi(x) \in C^\infty(\mathbb{R})$, define $\varphi_n(x) = \begin{cases}
        -1, &x \leq 0\\
        \varphi(nx), & 0 < x < \frac{1}{n}\\
        1,& x \geq \frac{1}{n},
    \end{cases} $
    then $\varphi_n(x) \to \mathrm{sgn}(x)$, $\forall~ x \neq 0$.

    Let $$f_n(x) = \begin{cases}
        -x, & x \leq0\\
        \int_0^x \varphi_n(t) \diff t, & x >0.
    \end{cases}$$
    $$
    \implies f_n'(x) = \varphi_n(x.)
    $$
    Because $f(x) = \abs{ x } = \int_0^{\abs{ x }} 1 \diff t.$
    $$
    \implies \abs{ f(x) - f_n(x) } \leq \int_0^{\abs{ x }}\abs{ \varphi_n(t) - 1 }\diff  t \leq \int_0^{\frac{1}{n}}\abs{ \varphi_n(t)-1 }\diff  t \leq 2 \frac{1}{n}.
    $$

    Apply It\^o formula to $f_n$ to get that
    $$
    f_n(X_t) = f_n(X_0) + \int_0^t f_n'(X_s) \diff X_s + \frac{1}{2}\int_0^t f_n''(X_s) \diff [X]_s.
    $$
    Note
    $$
    \int_0^t f_n'(X_s) \diff X_s = \int_0^t f_n'(X_s) \diff M_s + \int_0^t f_n'(X_s) \diff A_s,
    $$
    and
    $$
    \int_0^t f_n'(X_s) \diff A_s = \int_0^t \varphi_n(X_s) \diff A_s.
    $$

    We have $\varphi_n(x) \uparrow \mathrm{sgn}(x)$, $\forall~x$. By Monotone convergence,
    $$
    \abs{ \int_0^t \varphi_n(X_s) - \mathrm{sgn}(X_s) \diff A_s } \leq \int_0^t \abs{ \varphi_n(X_s) - \mathrm{sgn}(X_s) } \abs{ \diff A_s } \to 0,
    $$
    which implies $\int_0^t \varphi_n(X_s)\diff A_s \to \int_0^t \mathrm{sgn}(X_s) \diff A_s$ uniform in $t$.

    Next, $T_k = \inf\{t \geq 0 \colon \abs{ M_t }\geq k\}$. Then $M^{T_k}$ is bounded. 
    
    We may assume $M$ is bounded. Then
    $$
    \norm{\int_0^\infty \left( \mathrm{sgn}(X_s) - f_n'(X_s) \right) \diff M_s}_2^2 = \mathbb{E}\int_0^\infty \left( \mathrm{sgn}(X_s) - f_n'(X_s) \right)^2 \diff [M]_s \to 0.
    $$

    So $$
    \mathbb{E}\left[ \sup_{t \leq T} \abs{ \int_0^t \left( \mathrm{sgn}(X_s) - f_n'(X_s) \right) \diff M_s}^2  \right] \xlongrightarrow{n \to \infty} 0.
    $$

    Hence 
    $$
    \sup_{t \geq 0} \int_0^t \left( \mathrm{sgn}(X_s) - f_n'(X_s) \right) \diff M_s \xlongrightarrow{a.s.} 0.
    $$
    $\implies \int_0^t f_n'(X_s) \diff X_s \xlongrightarrow{a.s.} \int_0^t \mathrm{sgn}(X_s) \diff  X_s$, $f_n(X_t) \xlongrightarrow{a.s.} \abs{ X_t }$, $f_n(X_0) \xlongrightarrow{a.s.} \abs{ X_0 }$.
    
    $\implies C_t^n \triangleq \int_0^t \frac{1}{2} f_n''(X_s) \diff [X_s]$, $C_t^n \xlongrightarrow{a.s.} l_t \triangleq \abs{ X_t } - \abs{ X_0 } - \int_0^t \mathrm{sgn}(X_s) \diff X_s.$

    $\diff C_s^n = \frac{1}{2}f_n''(X_s) \diff [X]_s \implies \forall~ \delta > 0$, if $\frac{1}{n} < \delta$, then
    $$
        \abs{ X_s } > \delta > \frac{1}{n} \implies f_n''(X_s) = 0.
    $$
    $$
    \implies\int_0^t \mathsf{1}_{\{\abs{ X_s } > \delta\}} \diff C_s^n = 0.
    $$

    Hence $$\int_0^t \mathsf{1}_{\{\abs{ X_s }> \delta\}} \diff  l_s \leq \lim_{n \to \infty} \int_0^t \mathsf{1}_{\{\abs{ X_s }> \frac{\delta}{2}\}} \diff C_s^n = 0.$$

    $$
    \int_0^t \varphi_\delta(X_s) \diff l_s = \lim_{n\to \infty}\int_0^t \varphi_\delta(X_s) \diff C_s^n \leq \liminf_{n \to \infty} \int_0^t \mathsf{1}_{\{\abs{ X_s } > \frac{\delta}{2}\}} \diff  C_s^n.
    $$

    We prove $\int_0^t \mathsf{1}_{\{\abs{ X_s }> \delta\}}\diff l_s = 0$, $\forall~ \delta >0$ $\implies \int_0^t \mathsf{1}_{\{\abs{ X_s } \neq 0\}} \diff  l_s =0.$
} 

Let $X,Y$ be continuous semimartingale. Then 
$$
\int_0^t Y_t \diff X_t = \text{It\^o Integral.}
$$

Stratonovich Integral
$$
\int_0^t Y_t \partial X_t \triangleq \int_0^t Y_t \diff X_t + \frac{1}{2}[X,Y]_{t}.
$$

Recall that 
$$
\begin{aligned}
X_t Y_t - X_0 Y_0 &= \int_0^t X_s \diff Y_s + \int_0^t Y_s \diff X_s + [X,Y]_t \\&= \int_0^t X_s \diff Y_s + \frac{1}{2}[X,Y]_t + \int_0^tY_s \diff X_s + \frac{1}{2}[X,Y]_t\\
&=\int_0^t X_s \partial Y_s + \int_0^t Y_s \partial X_s.
\end{aligned}
$$

\thmnn{
    $f(X_t) = f(X_0) + \int_0^t f'(X_s) \partial X_s$.
}

\pf{
    By It\^o formula,
    $$
    f(X_t) = f(X_0) + \int_0^t f'(X_s) \diff X_s + \frac{1}{2} \int_0^t f''(X_s) \diff [X_s].
    $$

    By definition, 
    $$
    \int_0^t f'(X_s) \partial X_s = \int_0^t f'(X_s) \diff X_s + \frac{1}{2}\left[ f'(X),X \right]_t.
    $$
    $$
    \diff f(X_t) = f'(X_t) \diff X_t + \frac{1}{2}f''(X_t) \diff [X]_t
    $$
    $$
    \implies \diff f'(X_t) = f''(X_t) \diff X_t + \frac{1}{2} f'''(X_t) \diff [X]_t.
    $$
    So
    $$
    f'(X_t) = f'(X_0) + \int_0^t f''(X_s) \diff X_s  + \frac{1}{2}\int_0^t f'''(X_s) \diff [X]_s.
    $$
    $$
    \implies \left[ f'(X),X \right]_t = \left[ \int_0^t f''(X_s) \diff X_s, X_t \right] = \int_0^t f''(X_s) \diff [X]_s.
    $$
    ($[H\cdot M, K \cdot N]_t = \int_0^t HK \diff [M,N]_s$).

    So 
    $$
    f(X_t) = f(X_0) + \int_0^t f'(X_s) \diff X_s + \frac{1}{2} \left[ f'(X), X \right]_t = f(x_0) + \int_0^t f'(X_s) \partial X_s.
    $$
}

\thm{Riemam-Sum Approximation}{
    Let $C$ be a locally bounded adapted $L$-process. Let $X$ be a continuous semimartingale. Then
    $$
    \sum_{i=1}^{\infty}c(t_{i-1}^n)(X_{t_i^n} - X_{t_{i-1}^n}) \to \int_0^t c(s) \diff X_s.
    $$
    with $t_i^n = t \wedge \frac{i}{2^n}$. The convergence is uniform in probability on compacts, i.e. $\forall~ K > 0$,
    $$
    \sup_{t \leq K}\abs{ I^n(t)-I(t) } \xlongrightarrow{P} 0.
    $$
}

\pf{
    Let $c^n(t) = c(t_{i-1}^n)$ for $t \in (t_{i-1}^n,t_i^n]$. Then $c^n(t) \to c(t)$ for all $t \geq 0.$

    Also, 
    $$
    \begin{aligned}
    I^n(t) &= \sum_{i=1}^\infty c(t_{i-1}^n)(X(t_{i}^n) - X(t_{i-1}^n))\\
    &= \int_0^t c^n(s) \diff X_s \\
    &= (c^n,X)_t = \int_0^t c^n(s)\diff M_s + \int_0^t c^n(s) \diff A_s.
    \end{aligned}$$

    $\forall~ t \leq k$, $\abs{ c^n(t) } \leq c_k$
    $$
    \implies\int_0^t c^n(s) \diff A_s \to \int_0^t c(s)\diff A_s.
    $$

    By localization, we may assume $M$ is bounded. If $c$ is bounded and $M \in \mathcal{M}_0^2$, then 
    $$
    \int_0^t c^n(s) \diff M_s \to \int_0^t c(s) \diff M_s.
    $$
    $$
    \mathbb{E}\left[ \sup_{t \leq k}\abs{ \int_0^t (c^n(s) - c(s)) \diff M_s }^2 \right] \leq c_k \mathbb{E}\int_0^k (c^n(s) - c(s))^2 \diff [M]_s \to 0.
    $$
}
\par\noindent\hrulefill\par


\subsection{Stratonovich Calculus}
Let $X,Y$ be continuous semimartingale. Define
$$\begin{aligned}
&S^n(t) \equiv \sum_{i=1}^n \frac{1}{2} \left( Y(t_{i}^n) +Y(t_{i-1}^n) \right)\left( X(t_{i}^n) - X(t_{i-1}^n) \right)\\
&S(t) = \int_0^t Y_s \partial X_s.
\end{aligned}
$$
Then 
$$
\sup_{t \leq k}\abs{ S^n(t) - S(t) } \xlongrightarrow{P} 0.
$$

\pf{
    $$
    \begin{aligned}
    S^n(t) &= \sum_{i=1}^{\infty}\left[ Y(t_{i-1}^n) + \frac{1}{2}(Y(t_{i}^n) - Y(t_{i-1}^n) ) \right](X(t_{i}^n) - X(t_{i-1}^n)) \\
    &= \sum_{i=1}^{\infty} Y(t_{i-1}^n)(X(t_{i}^n) - X(t_{i-1}^n)) + \frac{1}{2} \sum_{i=1}^{\infty} \Delta Y(t_i^n)\Delta X(t_i^n).
    \end{aligned}
    $$

    It remains to prove 
    $$
    \begin{aligned}
    &\sum_{i=1}^{\infty}\Delta Y(t_{i}^n)\Delta X(t_{i}^n) \to [X,Y]_t = \frac{1}{4}([X+Y]_t - [X-Y]_t) \\
    &\Leftrightarrow \sum_{i=1}^{\infty}\Delta X(t_i^n)^2 \to [X]_t \\
    & \Leftrightarrow \sum_{i=1}^{\infty}\left( \Delta M(t_i^n) + \Delta A(t_i^n) \right)^2 \to [X]_t = [M]_t.
    \end{aligned}
    $$

    We have $\sum_{i=1}^\infty\left( \Delta M(t_i^n) \right)^2 \to [M]_t$,
    $$
    \abs{ \sum_{i=1}^{\infty}\Delta M(t_i^n) \Delta A(t_i^n) } \leq \sup_{i \geq 1}\abs{ \Delta M(t_i^n) } \sum_{i=1}^{\infty}\abs{ \Delta A(t_i^n) } \leq \sup_{\substack{t \leq k\\ i \geq 1}} \abs{ \Delta M(t_i^n) } \mathrm{Var}_A(t) \xlongrightarrow{a.s.} 0.
    $$
}

\section{Stochastic Differenctial Equation}
$$
\begin{aligned}
    &\diff X_t = b(t,X_t) \diff t + \sigma(t,X_t) \diff W_t,\\
    & X_t = X_0 + \int_0^t b(s,X_s) \diff s + \int_0^t \sigma(s,X_s) \diff W_s.
\end{aligned}
$$

d-dimension SDE ...

\defn{Strong solution}{
    A strong solution of the SDE, on a given probability space ($\Omega,\mathcal{F},\mathsf{P}$) and a given B.M. $(W_t,\mathcal{F}_t^W)_{t \geq 0}$, with initial condition $\xi$ is a process $\{X_t, t \geq 0\}$ with continuous sample path and 
    \begin{itemize}
        \item[(i)] $X \in \mathcal{F}_t$ with $(\mathcal{F}_t)_{t \geq 0}$ being the augmented filtration of the B.M. $(\mathcal{F}_t^W)_{t \geq 0}$.
        \item[(ii)] $\mathsf{P}(X_0 = \xi) = 1$.
        \item[(iii)] $$
        \mathsf{P}\left( \int_0^t \abs{ b_i(s,X_s) } + \sigma_{ij}^2(s,X_s) \diff s < \infty \right) = 1 
        $$   
        for all $i,j,t$.
        \item[(iv)] $X_t = X_0 + \int_0^t b(s,X_s) \diff s + \int_0^t \sigma(s,X_s) \diff W_s$.
    \end{itemize}
}

\defn{Strong Uniqueness}{
    Given $b(t,x)$ and $\sigma(t,x)$. Suppose $\forall~$B.M. $\{W_t\}_{t \geq 0}$ on some $(\Omega,\mathcal{F},\mathsf{P})$ and $\xi$ is the initial condition, and $\{\mathcal{F}_t\}$ is the augmented filtration, if $X$ and $\tilde{X}$ are two strong solutions with respect to the sam $W$, then 
    $$
    \mathsf{P}\left(X_t = \tilde{X}_t, \forall~ t \geq 0\right)  = 1.
    $$
}

\ex{}{
    \begin{itemize}
        \item[(i)] $\diff X_t = \diff W_t \implies X_t = W_t + \xi_0.$
        \item[(ii)] $\diff X_t = b(t,X_t)\diff t + \diff W_t$ where $b:[0,\infty) \times \mathbb{R} \to \mathbb{R}$ is bounded Borel-mble and $\forall~ x \leq y$, we have $b(t,x) \geq b(t,y)$, $\forall~ t \geq 0$. Fot this SDE, strong uniqueness holds.  
    \end{itemize}
}

\pf{
    Eg.(ii) $\forall~$B.M. on $(\Omega,\mathcal{F},\mathsf{P})$, $\forall~\xi$. If $X^{(1)}$ \& $X^{(2)}$ are strong solutions, then 
    $$
    X_t^{(j)} = \xi + \int_0^t b(s,X_s^{(j)}) \diff s + W_t, \quad \forall~ j=1,2.
    $$  
    $$
    X_t^{(1)} - X_t^{(2)} = \int_0^t \left[ b(s,X_s^{(1)}) - b(s,X_s^{(2)}) \right] \diff s.
    $$
    $$
    \implies \diff  \Delta_t = \left[ b(t,X_t^{(1)}) - b(t,X_t^{(2)}) \right] \diff t.
    $$
    By It\^o's formula,
    $$
    \Delta_t^2 = \int_0^t 2 \Delta_s \diff \Delta_s + \frac{1}{2}\int_0^t 2 \diff [\Delta]_t = 2\int_0^t \left( X_s^{(1)} - X_s^{(2)} \right)\left[ b(s,X_s^{(1)}) - b(s,X_s^{(2)}) \right] \diff s.
    $$

    If $x \leq y$, then $b(s,x) - b(s,y) \geq 0.$

    If $x > y$, then $b(s,x) - b(s,y) \leq 0$.

    $\implies (x-y)(b(s,x) - b(s,y)) \leq 0$. So 
    $$
    \Delta_t^2 = 2\int_0^t \left( X_s^{(1)} - X_s^{(2)} \right)\left[ b(s,X_s^{(1)}) - b(s,X_s^{(2)}) \right] \diff s \leq 0, \quad \text{a.s.}
    $$

    $\implies \mathsf{P}\left( \Delta_t = 0,\forall~ t \geq 0 \right) = 1$.

    $\implies$ strong uniqueness.
}

\thm{It\^o Uniqueness Theorem}{
    Suppose $b(t,x)$ and $\sigma(t,x)$ are locally Lipschitz in the space variable: $\forall~ n \geq 1$, $\exists~ K_n >0$ s.t. $\forall~ t \geq 0$, $\norm{ x },\norm{ y } \leq n$, we have
    $$
    \norm{ b(t,x)-b(t,y) } + \norm{ \sigma(t,x) - \sigma(t,y) } \leq K_n \norm{ x-y }.
    $$
    Then strong uniqueness holds.
}

\lem{Gronwall's Inequality}{
    Suppose a continuous function $g(t)$ satisfies $0 \leq g(t) \leq \alpha(t) + \beta\int_{0}^{t} g(s) \diff s$ for all $t \in [0,T]$ with $\beta \geq 0$ and $\alpha:[0,T] \to \mathbb{R}$ integrable. Then 
    $$
    g(t) \leq \alpha(t) + \beta \int_0^t \alpha(s) e^{\beta(t-s)} \diff s, \quad \forall~ t \leq T.
    $$
}

\pf{
    $I(t):= \int_{0}^{t} g(s)\diff s$.

    $\implies I'(t) =  g(t) \leq  \alpha(t) + \beta I(t)$.

    $\implies I'(t) - \beta I(t) -  \alpha(t) \leq 0$.

    Notice that
    $$
    \left[ e^{-\beta t} I(t) \right]' = \left( I'(t) - \beta I(t) \right) e^{-\beta t} \leq \alpha(t) e^{-\beta t}.
    $$

    $e^{-\beta t}I(t) - 0 = \int_0^t \left[ e^{-\beta s}I(s) \right]' \diff s \leq \int_0^t \alpha(s) e^{-\beta s} \diff s$.

    $\implies I(t) \leq \int_0^t \alpha(s) e^{\beta(t-s)} \diff s$.
    
    $\implies g(t) \leq \alpha(t) + \beta I(t) = \alpha(t) + \beta \int_0^t \alpha(s) e^{\beta(t-s)} \diff s$.
}

\pf{
    of It\^o Uniqueness Theorem.

    Given $W$ and $(\Omega,\mathcal{F} , \mathsf{P})$ and $\xi$. If $X$ and $\tilde{X}$ are two strong solutions with respect to the same $W$. Define
    $$
    \begin{aligned}
        &\tau_n = \inf \{t \geq 0\colon \norm{ X_t } \geq n\}\\
        &\tilde{\tau}_n = \inf \{t \geq 0\colon \norm{ \tilde{X}_t } \geq n\}.
    \end{aligned}
    $$

    Set $\tau_n \wedge \tilde{\tau}_n$. Then $\lim_{n \to \infty} S_n = \infty$, $\mathsf{P}$-a.s.

    Note that 
    $$
        X_{t \wedge S_n} - \tilde{X}_{t \wedge S_n} = \int_0^{t \wedge S_n} \left[b(s,X_s) - b(s,\tilde{X}_s)\right] \diff s + \int_0^{t \wedge S_n}\left[ \sigma(s,X_s) - \sigma(s,\tilde{X}_s) \right]
    $$
    $$\implies \mathbb{E} \norm{ X_{t \wedge S_n} - \tilde{X}_{t \wedge S_n} }^2 \leq 4 \mathbb{E}\left[ \int_0^{t \wedge S_n}(b-b)\diff s \right]^2 + 4 \mathbb{E}\left[ \int_0^{t \wedge S_n} (\sigma - \sigma)\diff W_s \right]^2$$.

    $(X_t^1,X_t^2,\dots,X_t^d) = \left( \int_0^t b^1 \diff s + \int_0^t \sigma^1 \diff W_s,\dots,\int_0^t b^d \diff s + \int_0^t \sigma^d \diff W_s \right)$.

    $$ \begin{aligned}
    \implies \norm{ X }^2 &= \sum_{i=1}^d (X_t^j)^2 \\&= \sum_{i=1}^d \left( \int_0^t b^j \diff s + \int_0^t \sigma^j \diff W_s \right)^2 \\&\leq 2 \sum_{i=1}^d\left(\int_0^t b^j \diff s\right)^2 + \left(\int_0^t \sigma^j \diff W_s\right)^2
    \end{aligned} $$

    $$
    \begin{aligned}
        \mathbb{E} \norm{ X_{t \wedge S_n} - \tilde{X}_{t \wedge S_n} }^2  \leq &2 \sum_{i=1}^{d} \mathbb{E}\left( \int_0^{t \wedge S_n} \left( b_i(s,X_s) - b(s,\tilde{X}_s) \right) \diff s\right)^2 \\
        &+ 2 \sum_{i=1}^{d} \mathbb{E}\left( \sum_{j=1}^{r}\int_0^{t \wedge S_n} \left( \sigma_{ij}(s,X_s) - \sigma_{ij}(s,\tilde{X}_s) \right) \diff W_s^{(j)}\right)^2 \\
        \leq& 2 \sum_{i=1}^{d} \mathbb{E} \left[ t \cdot \int_0^{t \wedge S_n}\left( b_{i}(s,X_s) - b_i(s,\tilde{X}_s) \right)^2 \diff s \right]\\
        &+ 2\sum_{i=1}^{d}\sum_{j=1}^{r} \mathbb{E} \left[ \int_0^{t \wedge S_n}(\sigma_{ij} - \sigma_{ij}) \diff W_s^{(j)} \right]^2
        \end{aligned}
    $$
    And $\mathbb{E} \left[ \int_0^{t\wedge S_n} \left(\sigma_{ij} - \sigma_{ij}\right) \diff W_s^{(j)}\right]^2 = \mathbb{E}\left[ \int_0^{t \wedge S_n}\left( \sigma_{ij} - \sigma_{ij} \right)^2 \diff s \right]$.

    We conclude that
    $$
    \begin{aligned}
    \mathbb{E} \norm{ X_{t \wedge S_n} - \tilde{X}_{t\wedge S_n} }^2 \leq & 2 t \mathbb{E}\int_0^{t \wedge S_n}\norm{ b(s,X_s) - b(s,\tilde{X}_s) }^2 \diff s + 2 \mathbb{E}\int_0^{t \wedge S_n} \norm{ \sigma(s,X_s) - \sigma(s,\tilde{X}_s) }^2 \diff s \\
    \leq & 2T \mathbb{E}\int_0^{t \wedge S_n} K_n^2 \norm{ X_s - \tilde{X}_s }^2 \diff s + 2 \mathbb{E}\int_0^{t \wedge S_n} K_n^2 \norm{ X_s - \tilde{X}_s }^2 \diff s \\
    \leq& 2(T+1)K_n^2 \int_0^t \mathbb{E}\norm{ X_{s\wedge S_n} - \tilde{X}_{s \wedge S_n} }^2 \diff s,
    \end{aligned}
    $$
    since
    $$
    \norm{ \sigma }^2 = \sum_{i=1}^{d}\sum_{j=1}^{r} \sigma_{ij}^2.
    $$

    If we let $g(t ) = \triangleq \mathbb{E} \norm{ X_{t \wedge S_n} - \tilde{X}_{t \wedge S_n} }^2$, then $g(t) \leq 2(T+1)K_n^2 \int_0^t g(s) \diff s.$

    By Gronwall's inequality, we get $g(t) \leq 0$.
    
    $\implies X_{t \wedge S_n} = \tilde{X}_{t \wedge S_n}$ a,s, $\forall~ t \geq 0$.

    $\implies X_{t \wedge S_n} = \tilde{X}_{t \wedge S_n}$, $\forall~t \geq 0$, a.s.

    Let $S_n \uparrow \infty$ to conclude $X_t = \tilde{X}_t$, $\forall
    ~ t \geq 0$, a.s.
}

\thm{Existence}{
    Suppose $\exists~K > 0$ s.t. $\forall~ t \geq 0$, $\forall~ x,y \in \mathbb{R}^d$,
    $$
    \begin{cases}
        \norm{ b(t,x) - b(t,y) } + \norm{ \sigma(t,x) - \sigma(t,y) } \leq K \norm{ x-y }\\
        \norm{ b(t,x) } + \norm{ \sigma(t,x) } \leq K(1 + \norm{ x }).
    \end{cases}
    $$
    On $(\Omega,\mathcal{F},\mathsf{P})$, let $\xi$ be R.V. and r-dimensional $B.M.$ $W=\{W_t\}_{t\geq 0}$ independent of $\xi$, and $\mathbb{E}\norm{ \xi }^2 < \infty.$ Let $(\mathcal{F}_t)_{t \geq 0}$ be the augmented filtration. Then $\exists~$continuous adapted process $X=\{X_t\}_{t \geq 0}$ which is a strong solution with respect to $W$ and $X_0 = \xi$ a.s.

    Moreover, $\forall~ T > 0$, $\exists~ C= C(K,T) > 0$ s.t. $\forall~ 0 \leq t \leq T$,
    $$
    \mathbb{E} \norm{ X_t }^2 \leq C(1 + \mathbb{E}\norm{ \xi }^2) e^{Ct}.
    $$
}

\pf{
    Define $X_t^{(0)}= \xi$, $\forall~ t \geq 0$.

    For any $k \geq 1$, define 
    $$
    X_t^{(k)} \triangleq \xi + \int_0^t b(s,X_s^{(k-1)}) \diff s + \int_0^t \sigma(s,X_s^{(k-1)})\diff  W_s.
    $$

    Then $\forall~ T > 0$, $\exists~C=C(K,T) > 0$ s.t.
    $$
    \mathbb{E}\norm{ X_t^{(k)} }^2 \leq C(1+ \mathbb{E}\norm{ \xi }^2) e^{Ct}, \quad\forall~ t \leq T, \forall~k \geq 0.
    $$  

    (Exercise)

    For any $k \geq 0$,
    $$
    X_t^{(k+1)} - X_t^{(k)} = \int_0^t\left( b(s,X_s^{(k)}) - b(s,X_s^{(k-1)}) \right) \diff s + \int_0^t\left( \sigma(s,X_s^{(k)}) - \sigma(s,X_s^{(k-1)}) \right) \diff  W_s.
    $$

    $$
    \begin{aligned}
    \mathbb{E} \left[ \sup_{s \leq t} \norm{ B_s }^2 \right] &= \mathbb{E} \norm{ \sup_{s \leq t}\int_0^s \left( b(u,X_u^{(k)}) - b(u,X_u^{(k-1)}) \right) \diff  u}^2 \\
    &\leq \mathbb{E}\left[ t \cdot \int_0^t \norm{ b(u,X_u^{(k)}) - b(u,X_u^{(k-1)}) } ^2 \diff u\right]\\
    & \leq K^2 + \int_0^t \mathbb{E} \norm{ X_u^{(k)} - X_u^{(k-1)} }^2 \diff u.
    \end{aligned}
    $$

    $M_t = \int_0^t \sigma(s,X_s^{(k)}) - \sigma(s,X_s^{(k-1)}) \diff W_s$.

    $[M]_t = \int_0^t \norm{ \sigma(s,X_s^{(k)}) - \sigma(s,X_s^{(k-1)}) }^2 \diff s$.

    $\implies \mathbb{E}[M]_t \leq \int_0^t K^2\left(1+\norm{ X_s^{(k)} } + \norm{ X_s^{(k-1)} }\right)^2 \diff s < \infty$ by $\mathbb{E}\norm{ X_t^{(k)} }^2 \leq Ce^{Ct}.$

    $$
    \begin{aligned}
        \mathbb{E}\left[ \sup_{s \leq t} \norm{ M_s }^2 \right] &\leq C \mathbb{E}[M_t]\\
        &= C \mathbb{E}\int_0^t \norm{ \sigma(s,X_s^{(k)}) - \sigma(s,X_s^{k-1}) }^2 \diff s\\
        &\leq C \mathbb{E}\int_0^t K \norm{ X_s^{(k)} - X_s^{(k-1)} }^2 \diff s.
    \end{aligned}
    $$

    Combine the two inequalities to get 
    $$
    \begin{aligned}
        \mathbb{E}\left[ \max_{s \leq t} \norm{ X_s^{(k+1)} - X_s^{(k)} }^2 \right] &\leq 2 \mathbb{E}\left[ \max_{s \leq t} \norm{ B_s }^2 \right] + 2 \mathbb{E} \left[ \max_{s \leq t} \norm{ M_s }^2 \right] \\
        & \leq C\cdot 4 K^2(1+T) \int_0^t \mathbb{E}\norm{ X_s^{(k)}- X_s^{(k-1)} }^2 \diff s \\
        & \leq L \int_0^t \mathbb{E}\left( \norm{ X_u^{(k)} - X_u^{(k-1)}}^2  \right)\diff s.
    \end{aligned}
    $$

    When $k=1$, we have 
    $$
        \mathbb{E}\left( \norm{ X_s^{(1)} - X_s^{(0)} }^2 \right) \leq 2 \mathbb{E}\norm{ X_t^{(1)} }^2 + 2 \mathbb{E}\norm{ X_t^{(0)} }^2 \leq C^*
    $$
    for all $0 \leq t \leq T.$

    Hence 
    $$
    \mathbb{E}\left( \max_{0 \leq s \leq t} \norm{ X_s^{(2)} - X_s^{(1)} }^2 \right) \leq L\int_0^t C^* \diff s = C^*(Lt).
    $$

    $$\implies
        \mathbb{E}\left( \max_{0 \leq s \leq t} \norm{ X_s^{(3)} - X_s^{(2)} }^2 \right) \leq L\int_0^t C^* (Ls) \diff s = C^* \frac{1}{2}(Lt)^2.
    $$

    By induction, we can prove $\forall~ k \geq 1$,
    $$
    \mathbb{E}\left( \max_{s \leq t} \norm{ X_s^{(k+1)} - X_s^{(k)} }^2 \right) \leq C^* \frac{1}{k!} (Lt)^k.
    $$

    By chebyshev's Inequality,
    $$
    \mathsf{P}\left(\max_{0 \leq t \leq T} \norm{ X_t^{(k+1)} - X_t^{(k)} } > \frac{1}{2^{k+1}}\right) \leq 4^{k+1} C^* \frac{1}{k!} (LT)^k.
    $$

    $\sum_{k=1}^\infty \mathsf{P}(\max \norm{ \dots } > \frac{1}{2^{k+1}}) < \infty$.

    $\implies$ B.C. Lemma gives
    $$
    \mathsf{P}\left(\max_{t \leq T} \norm{ X_t^{(k+1)} - X_t^{(k)} } > \frac{1}{2^{k+1}}~\text{i.o.}\right) = 0.
    $$  

    $\implies \mathsf{P}$-a.s. $\exists~N(\omega)< \infty$ s.t. $\forall~ k \geq N$, 
    $$
        \max_{t \leq T} \norm{ X_t^{(k+1)} - X_t^{(k)} }\leq \frac{1}{2^{k+1}}.
    $$

    $\implies \forall~ k \geq N$, $\forall~ m \geq 1$, we have
    $$
    \max_{t \leq T} \norm{ X_t^{(k+m)} - X_t^{(k)} } \leq\sum_{j=k}^{k+m-1} \frac{1}{2^{j+1}} \leq \frac{1}{2^k}.
    $$
    
    Hence $\left\{ X_t^{(k)}(\omega),0 \leq t \leq T \right\}$ is a Cauchy sequence in $\norm{ \dots }_{sup}$.

    Hence $\exists~$limit $\left\{ X_t(\omega), 0 \leq t \leq T \right\}$ s.t. 
    $$
    \lim_{k\to \infty} \max_{t \leq T} \norm{ X_t^{(k)} - X_t } = 0.
    $$

    Since $T > 0$ is arbitrary, $\exists~X=\left\{ X_t,0 \leq t < \infty \right\}$ s.t. $\forall~ T > 0$, $\sup_{t \leq T}\norm{ X_t^{(k)} - X_t } \to 0 $ as $k \to \infty$. 

    $$
    \begin{aligned}
        \mathbb{E}\norm{ X_t }^2 &= \mathbb{E}\lim_{k \to \infty} \norm{ X_t^{(k)} }^2 \\
        &\leq \lim_{k \to \infty} \mathbb{E} \norm{ X_t^{(k)} }^2 \leq C(1+\mathbb{E}\norm{ \xi }^2) e^{Ct}.
    \end{aligned}
    $$

    $\implies \mathsf{P}\left( \int_{0}^{t}\abs{ b(s,X_s) }\diff s + \int_0^t \sigma_{ij}^2(s,X_s)  \diff  s < \infty\right) = 1.$

    $$
    X_t^{(k+1)} = \xi + \int_0^t b(s,X_s^{(k)}) \diff s + \int_0^t \sigma(s,X_s^{(k)})\diff W_s.
    $$

    $$
    X_t = \xi + \int_0^t b(s,X_s) \diff s + \int_0^t \sigma(s,X_s) \diff W_s.
    $$

    Need to prove $X_t^{(k+1)} \xlongrightarrow{k \to \infty} X_t$ a.s. 
}

\thm{Yamada-Watanabe}{
    Let $\diff  X_t = \sigma(X_t) \diff W_t + b(X_t) \diff t$,
    where
    $$
    \sigma(x) =
    \begin{pmatrix}
    \sigma_{1}(x_1) & \cdots & 0 \\
    \vdots & \ddots & \vdots \\
    0 & \cdots & \sigma_{n}(x_ )
    \end{pmatrix}
    $$
    $$
    b(x) = (b^1(x),b^2(x),\dots,b^n(x))
    $$

    such that 
    \begin{itemize}
        \item[(i)] $\exists~$positive increasing function $\rho(u)$, $u\in (0,\infty)$ s.t. 
        $$
        \abs{ \sigma_i(\xi) - \sigma_i(\eta) } \leq \rho\left( \abs{ \xi - \eta } \right)
        $$
        for all $\xi,\eta \in \mathbb{R}^1$, $i=1,\dots,n$ and 
        $$
        \int_{0^+} \rho^{-2}(u) \diff u = \infty.
        $$
        i.e. $\forall~ \varepsilon >0$, $\int_0^\varepsilon \rho^{-2}(u) \diff u = \infty$.

        \item[(ii)] $\exists~$positive increasing concave function $k(u)$, $u\in(0,\infty)$ s.t. 
        $$
        \abs{ b_i(x) - b_j(y) } \leq k \left( \abs{ x-y } \right)
        $$
        for all $x,y \in \mathbb{R}^n$, $i=1,\dots,n$ and $\int_{0^+} k^{-1}(u) \diff u = \infty.$
    \end{itemize}
    Then strong uniqueness holds.
}

\rmkb{
    If $\rho(u) = u^{\alpha}$, then strong uniqueness holds if $\alpha \geq \frac{1}{2}$; fails if $0 < \alpha < \frac{1}{2}$.
}


\prop{(Yamada \& Watanabe 1971)

    ``$d=1$'' Consider
    $$
    \diff X_t = b(t,X_t) \diff t + \sigma(t,X_t) \diff W_t.
    $$

    Assume that $\exists~$constant $K > 0$ and $h:[0,\infty) \to [0,\infty)$ is a strictly increasing function with $h(0) = 0$ such that
    $$
    \begin{cases}
        \abs{ b(t,x) - b(t,y) } \leq K \abs{ x-y }\\
        \abs{ \sigma(t,x) - \sigma(t,y) } \leq h(\abs{ x-y })
    \end{cases}
    $$
    hold for all $t \geq 0$ and $x,y \in \mathbb{R}$ with 
    $$
    \int_0^\varepsilon h(u)^{-2}\diff u = \infty, \quad \forall~ \varepsilon >0.
    $$ 
    Then strong uniqueness holds. 
}

\pf{
    Idea: Constant a $C^2$ function $\psi_n(x)$ s.t. $\psi_n(x) \uparrow \abs{ x }$. Let $a_0 = 1 > a_1 >\dots > a_n \downarrow 0$ such that
    $$
    \begin{aligned}
        &\int_{a_1}^{a_0} h(u)^{-2}\diff  u = 1,\\
        &\int_{a_2}^{a_1} h(u)^{-2} \diff u = 2,\\
        &\dots \\
        & \int_{a_n}^{a_{n-1}} h(u)^{-2} \diff  u = n,
    \end{aligned}
    $$
    So 
    $$
    \int_{a_n}^{a_{n-1} } \frac{2}{n }h(x)^{-2}\diff  x = 2.
    $$

    Then $\exists~\rho_n(x)$ s.t.
    $$
    \int_{a_n}^{a_{n-1}} \rho_n(x)\diff x = 1
    $$
    and $\rho_n(x) = 0$, $\forall~x  \notin (a_n,a_{n-1})$, and $x \mapsto \rho_n(x)$ is continuous.

    $$
    \implies \int_0^\infty \rho_n(x) \diff x = 1.
    $$

    Let $\psi_n'(x) = \int_0^x \rho_n(y) \diff y$, $\forall~ x \geq 0$.

    Let $\psi_n(x) = \int_0^x \psi_n'(y) \diff y$ for $x \geq 0$ and 
    $$
    \psi_n(x) = \psi_n(-x), \quad \forall~ x < 0.
    $$

    One can check that $\psi_n(x) \uparrow \abs{ x }$, $\forall~ x \in \mathbb{R}$.

    \textbf{Summary:} \begin{itemize}
        \item[(i)] $\psi_n(x) \to \abs{ x }$, $\forall~ x \in \mathbb{R}$,
        \item[(ii)] $\abs{ \psi_n'(x) } \leq 1$, $\forall~ x \in \mathbb{R}$,
        \item[(iii)] $\abs{ \psi_n''(x) } = \abs{ \rho_n(x) } \leq \frac{2}{n} h(x)^{-2}$.
    \end{itemize}

    Suppose $X^{(1)}$ \& $X^{(2)}$ are two strong solutions with $X_0^{(1)} = X_0^{(2)}$ a.s.

    Define 
    $$
    \tau_N = \inf \left\{ t \geq 0 \colon \int_0^t \sigma(s,X_s^{(1)})^2 \diff s \geq N \text{ or } \int_{0}^{t}\sigma(s,X_s^{(2)})^2 \diff s \geq N\right\}.
    $$
    $$
    \implies \mathbb{E}\int_0^{t \wedge \tau_N} \sigma(s,X_s^{(i)})^2 \diff s \leq N
    $$
    for $i=1,2$.

    We may assume that 
    $$
    \mathbb{E}\int_0^t \sigma(s,X_s^{(i)})^2 \diff s < \infty,
    $$
    for $\forall~ t \geq 0$, $i=1,2$, due to $\tau_N \uparrow \infty$ a.s.

    Otherwise we may prove
    $$
    X_s^{(1)} = X_s^{(2)}, \quad \forall~0 \leq s \leq \tau_N \text{ a.s. }\implies X_s^{(1)} = X_s^{(2)}, \quad \forall~ s \geq 0.
    $$

    Define $\Delta_t = X_t^{(1)} - X_t^{(2)}$.
    
    $$
    \begin{aligned}
    \implies \diff \Delta_t &= \diff X_t^{(1)} - \diff X_t^{(2)}\\
    &= \left[ \sigma(t,X_t^{(1)}) - \sigma(t,X_t^{(2)}) \right] \diff W_t + \left[ b(t,X_t^{(1)}) - b(t,X_t^{(2)}) \right] \diff t.
    \end{aligned}
    $$

    Apply It\^o's formula to get 
    $$
    \begin{aligned}
    \psi_n(\Delta_t) &= \psi_n(\Delta_0) + \int_0^t \psi_n'(\Delta_s) \diff \Delta_s + \frac{1}{2} \int_0^t \psi_n''(\Delta_s) \diff [\Delta]_s \\
    &=\psi_n(0) + \int_0^t\psi_n'(\Delta_s)[\sigma - \sigma]\diff W_s + \int_0^t \psi_n'(\Delta_s)[b-b]\diff s + \frac{1}{2}\int_0^t \psi_n''(\Delta_s)[\sigma - \sigma]^2 \diff s.
    \end{aligned}
    $$

    Take expectation to see that
    $$
    \begin{aligned}
        \mathbb{E}\psi_n(\Delta_t) &= \mathbb{E}\int_0^t \psi_n'(\Delta_s)\left[ b(t,X_t^{(1)}) - b(t,X_t^{(2)})  \right]\diff s + \mathbb{E}\int_0^t \frac{1}{2}\psi_n''(\Delta_s) [\sigma - \sigma]^2 \diff s \\
    &\leq \mathbb{E}\int_0^t K \abs{ X_s^{(1)} - X_s^{(2)} } \diff s + \mathbb{E}\int_0^t \frac{1}{2} \frac{2}{n} h(\abs{ \Delta_s })^{-2}\cdot h(\abs{ X_s^{(1)} - X_s^{(2)} })^2 \diff s\\
    &= K\int_0^t \mathbb{E}\abs{ X_s^{(1)} - X_s^{(2)} }\diff s + \frac{1}{n}t.
    \end{aligned}
    $$

    Let $n\to \infty$ to get 
    $$
    \mathbb{E}\abs{ \Delta_t } \leq K \int_0^t \mathbb{E} \abs{ \Delta_s }\diff s.
    $$

    By Gronwall's Inequality,
    $$
    \mathbb{E}\abs{ \Delta_t } = 0 \implies X_t^{(1)} = X_t^{(2)} \text{ a.s.}.
    $$

    By continuity,
    $$
    X_t^{(1)} = X_t^{(2)}, \quad \forall~ t \geq 0 \text{ a.s.}.
    $$
}

\prop{(R \& W)

    If $Y = X- X'$, and $h \colon [0,\infty) \to [0,\infty)$ is increasing s.t.
    $$
    \int_0^\varepsilon h(x)^{-2} \diff x = \infty, \quad \forall~\varepsilon > 0,
    $$
    then $l_{\cdot}^0(Y) = 0$ if 
    $$
    \int_0^t h(Y_s)^{-2} \mathsf{1}_{\{Y_s > 0 \}}\diff [Y]_s < \infty
    $$
    a.s. $\forall~ t > 0$.

    Here $l_t^{0} (Y)$ is the local time of $Y$ at 0 
    
    (Recall $\abs{ Y_t } = \abs{ Y_0 } + \int_0^t \mathrm{sgn}(Y_s) \diff Y_s + l_t^{0}(Y)$)
}

\pf{
    $$
    l_t^a ``=" \int_0^t \delta_a (Y_s) \diff [Y]_s.
    $$  
    Then $\forall~ $function $f \colon : \mathbb{R} \to \mathbb{R}$,
    $$
    \begin{aligned}
    \int_0^t f(Y_s) \diff [Y]_s & = \int_0^t \diff [Y]_s \int_{-\infty}^{\infty}f(a)\delta_a (Y_s)\diff a \\
    &= \int_{-\infty}^{\infty}f(a)\diff a \int_0^t \delta_a(Y_s)\diff [Y]_s \\
    &= \int_{-\infty}^{\infty}f(a)l_t^a(Y)\diff a.
    \end{aligned}
    $$
    
    So 
    $$
    \int_0^t h(Y_s)^{-2}\mathsf{1}_{\{Y_s > 0\}} \diff [Y]_s = \int_{-\infty}^{\infty}h(a)^{-2}\mathsf{1}_{\{a > 0\}} l_t^a \diff a = \int_0^\infty h(a)^{-2}l_t^a \diff a.
    $$

    If $l_t^a(Y) > 0$, then by the right-continuity of $a \mapsto l_t^a(Y)$, we get $\int_t^a (Y) \geq \varepsilon_0$ for some $\varepsilon_0 > 0$ and $0 < a < \varepsilon_0$.
    $$
    \implies \int_0^\infty h(a)^{-2}l_t^a \diff a \geq \int_0^{\varepsilon_0} h(a)^{-2} \varepsilon_0 \diff a = \infty.
    $$
    Contradiction.

    Hence we conclude that $l_t^0(Y) = 0$, $\forall~ t \geq 0$, a.s.
}

\textbf{Assumption:} $\begin{cases}
    \abs{ b(t,x)-b(t,y) } \leq K \abs{ x-y }\\
    \abs{ \sigma(t,x) - \sigma(t,y) } \leq h (\abs{ x-y })
\end{cases}$ and $\int_0^{\varepsilon}h(x)^{-2} \diff x = \infty$, $\forall~\varepsilon > 0$.

\pf{Second proof of (Y-W 1971) (Le Gall) 

    Since (Recall $\diff Y_s = \diff X_s^{(1)} - \diff X_s^{(2)}$)
    $$
    \begin{aligned}
    \int_0^t h(Y_s)^{-2}\mathsf{1}_{\{Y_s > 0\}} \diff [Y]_s &= \int_0^t h(Y_s)^{-2} \mathsf{1}_{\{Y_s > 0\}}(\sigma - \sigma)^2 \diff s\\
    &\leq \int_0^t h(Y_s)^{-2}\mathsf{1}_{\{Y_s > 0\}}h(\abs{ X_s^{(1)} - X_s^{(2)} })^2 \diff s \leq t < \infty \text{ a.s..}
    \end{aligned}
    $$

    $\implies l_t^0 (Y) = 0$ a.s.
    
    Thus 
    $$
        \abs{ Y_t } = \abs{ Y_0 } + \int_0^t \mathrm{sgn}(Y_s) \diff Y_s + l_t^0(Y).
    $$
    $$
    \begin{aligned}
        \implies \abs{ X_t^{(1)} - X_t^{(2)} } &= \int_0^t \mathrm{sgn}(X_s^{(1)} - X_s^{(2)})(b-b)\diff s + \int_0^t \mathrm{sgn}(X_s^{(1)} - X_s^{(2)}) (\sigma - \sigma)\diff W_s.
    \end{aligned}
    $$
    $$
    \implies \mathbb{E}\abs{ X_t^{(1)} - X_t^{(2)} } \leq \int_0^t \mathbb{E} K \abs{ X_s^{(1)} - X_s^{(2)} }\diff s.
    $$

    By Gronwall's, the proof is done.
}


\thm{Comparision Principle}{
    Given $(\Omega,\mathcal{F},\mathsf{P})$ and a B.M. $\{W_t,\mathcal{F}_t,t \geq 0\}$. If two continuous adapted process $X^{(1)},X^{(2)}$ satisfy that 
    $$
    \begin{cases}
        X_t^{(1)} = X_0^{(1)} + \int_0^t b_1(s,X_s^{(1)}) \diff s + \int_0^t \sigma(s,X_s^{(1)}) \diff  W_s,\\
        X_t^{(2)} = X_0^{(2)} + \int_0^t b_2(s,X_s^{(2)}) \diff s + \int_0^t \sigma(s,X_s^{(2)}) \diff  W_s.
    \end{cases}
    $$

    Assume \begin{itemize}
        \item[(1)] $\sigma(t,x)$, $b_j(t,x)$ are continuous, real-values functions on $[0,\infty) \times \mathbb{R}$.
        \item[(2)] $\abs{ \sigma(t,x) -\sigma(t,y) } \leq h(\abs{ x-y })$ with 
        $$\int_0^{\varepsilon}h(x)^{-2} \diff x = \infty.$$
        \item[(3)] $X_0^{(1)} \leq X_0^{(2)}$.
        \item[(4)] $b_1(t,x) \leq b_2(t,x)$.
        \item[(5)] either $b_1(t,x)$ or $b_2(t,x)$ satisfy $\abs{ b(t,x)-b(t,y) } \leq K \abs{ x-y }$.
    \end{itemize}

    Then $\mathsf{P}\left( X_t^{(1)} \leq X_t^{(2)}, \forall~ t \geq 0 \right) = 1$.
}

\pf{
    Assume $\abs{ b_1(t,x) - b_1(t,y) } \leq K \abs{ x-y }$. Again we assume 
    $$
    \mathbb{E}\int_0^t \sigma(s ,X_s^{(j)})^2 \diff s < \infty, \quad \forall~t \geq 0.
    $$

    Let $\psi_n(x)$ be as before. Define
    $$
    \varphi_n(x) = \psi_n(x) \mathsf{1}_{(0,\infty)}(x).
    $$
    $$
    \mathbb{E}[\varphi_n(\Delta_t)] \leq \mathbb{E}\int_0^t \varphi_n'(\Delta_s)\left[ b_1(s,X_s^{(1)}) - b_2(s,X_s^{(2)}) \right]\diff s + \frac{t}{n}
    $$
    where $\Delta_t = X_t^{(1)} - X_t^{(2)}$.

    So 
    $$
    \begin{aligned}
    \mathbb{E}\varphi_n(\Delta_t)- \frac{t}{n} &\leq \mathbb{E}\int_0^t \varphi_n'(\Delta_s)\left[ b_1(s,X_s^{(1)}) - b_2(s,X_s^{(2)}) \right]\diff s \\
    &\leq \mathbb{E}\int_0^t K \abs{ X_s^{(1)} - X_s^{(2)} }\varphi_n'(\Delta_s)\diff s\\
    & = \mathbb{E}\int_0^t K \Delta_s^+ \varphi_n'(\Delta_s^+) \diff s \\
    &\leq \mathbb{E}\int_0^t K \Delta_s^+ \diff s,
    \end{aligned}
    $$
    since $\varphi_n'(\Delta_s) = 0$ if $\Delta_s \leq 0$. $\varphi_n'(\Delta_s) > 0$ if $\Delta_s > 0$.

    By $\psi_n(x) \to \abs{ x }$, $\varphi_n(x) = \psi_n(x) \mathsf{1}_{(0,\infty)}(x) \to x^+$.

    Let $n \to \infty$ to get 
    $$
    \mathbb{E}\Delta_t^+ \leq K \int_0^t \mathbb{E} \Delta_s^+ \diff s.
    $$

    By Gronwall, $\mathbb{E} \Delta_t^+ = 0 \implies \mathsf{P}\left( X_t^{(1)} \leq X_t^{(2)} \right) = 1$, $\forall~ t > 0$.
}

\defn{Weak Solution}{
    A weak solution of SDE is a triple $(X,W)$, $(\Omega, \mathcal{F}, \mathsf{P})$, $\{\mathcal{F}_t\}$ where
    \begin{itemize}
        \item[(1)] ($\Omega, \mathcal{F}, \mathsf{P}$) probability space, $\{\mathcal{F}_t\}$ is filtration.
        \item[(2)] $(X_t, \mathcal{F}_t)$ is continuous and adapted $(W_t, \mathcal{F}_t)$ is an r-dim B.M.
        \item[(3)] $\displaystyle \mathsf{P}\left( \int_0^t b_j(s,X_s) + \sigma_{ij}(s,X_s)^2 \diff s < \infty, \quad \forall~t \geq 0  \right) = 1$.
        \item[(4)] $X_t = X_0 + \int_0^tb(s,X_s)\diff s + \int_0^t \sigma(s,X_s) \diff W_s$.
    \end{itemize}
}

\defn{Pathwise Uniqueness}{
    Suppose whenever $(X,W)$, $(\Omega, \mathcal{F}, \mathsf{P})$, $\{\mathcal{F}_t\}$ and $(\tilde{X},W)$, $(\Omega, \mathcal{F}, \mathsf{P})$, $\{\tilde{\mathcal{F}}_t\}$ are two weak solutions with $X_0 = \tilde{X}_0$ a.s., if $$
    \mathsf{P}\left( X_t = \tilde{X}_t,\quad \forall~ t \geq 0 \right) = 1,
    $$
    then pathwise uniqueness holds.
}

\defn{Weak Uniqueness (Uniqueness in law)}{
    If for any two weak solutions $(X,W)$, $(\Omega,\mathcal{F}, \mathsf{P})$, $\{\mathcal{F}_t\}$ and $(\tilde{X},\tilde{W})$, $(\tilde{\Omega}, \tilde{\mathcal{F}}_t, \tilde{\mathsf{P}})$, $\{\tilde{\mathcal{F}}_t\}$ with same initial distribution $\forall~ B \in B(\mathbb{R}^d)$, $$\mathsf{P}(X_0 \in B)= \tilde{\mathsf{P}}(\tilde{X}_0 \in B).$$ 

    We get $X$ and $\tilde{X}$ have the same law, i.e. $\forall~ t_0 \leq t_1 \leq \dots \leq t_n$,
    $$
    \mathsf{P}\left( X_{t_0} \in B_0,\dots, X_{t_n}\in B_n \right) = \tilde{\mathsf{P}}\left( \tilde{X}_{t_0}\in B_0, \dots, X_{t_n} \in B_n \right)
    $$
    for all $B_i \in B(\mathbb{R}^d)$, then weak uniqueness holds.
}

\rmkb{
    (Tanaka SDE)
    $$
    X_t = \int_0^t \mathrm{sgn}(X_s) \diff W_s 
    $$
    with $\mathrm{sgn}(x) = \begin{cases}
        1,& x> 0\\
        -1,& x \leq 0
    \end{cases}$.

    \textbf{Weak Uniqueness:} If $(X,W)$, $(\Omega, \mathcal{F}, \mathsf{P})$, $\{\mathcal{F}_t\}$ is a weak solution, then the process $(X_t)$ is a continuous, square-integrable martingale with
    $$
    [X]_t = \int_0^t \mathrm{sgn}(X_s)^2 \diff  s = t .
    $$
    
    So $X$ is a B.M. $\implies$ Weak Uniqueness holds.

    However, $(-X,W)$, $(\Omega, \mathcal{F}, \mathsf{P})$, $\{\mathcal{F}_t\}$ is also a weak solution.

    Because $$X_t = \int_0^t \mathrm{sgn}(X_t) \diff W_s \implies -X_t = \int_0^t -\mathrm{sgn}(X_s) \diff W_s = \int_0^t \mathrm{sgn}(-X_s) \diff W_s.
    $$ 
    $$
    \int_0^t\left( -\mathrm{sgn}(X_t) - \mathrm{sgn}(-X_s) \right)\diff  W_s \sim \int_0^t \left( -\mathrm{sgn}(X_s) - \mathrm{sgn}(-X_s) \right)^2 \diff s.
    $$

    Notice that $-\mathrm{sgn}(x) - \mathrm{sgn}(-x) = 0$ for all $x \neq 0$ and 
    $$
    \mathsf{P}\left( X_s \neq 0,\text{ for a.e. } s \geq 0 \right) = 1.
    $$

    W.P.1 
    $$
    \int_0^t (-\mathrm{sgn}(X_s) - \mathrm{sgn}(-X_s))^2 \diff s = 0 \implies \int_0^t -\mathrm{sgn}(X_s) \diff W_s = \int_0^t \mathrm{sgn}(-X_s) \diff W_s \text{ a.s. }
    $$

    $\implies - X_t = \int_0^t \mathrm{sgn}(-X_s)\diff W_s$.

    $\implies(-X,W)$ is also a weak solution.
    $$
    \implies \mathsf{P}(X_t \neq - X_t, \text{ for a.e. } t \geq 0) = 1.
    $$

    Hence pathwise uniqueness fails.

    \textbf{Weak Existence:} $(X,W)$, $(\Omega, \mathcal{F}, \mathsf{P})$, $\{\mathcal{F}_t\}$.

    \textbf{Construction:} Given $(\Omega, \mathcal{F}, \mathsf{P})$ and a B.M. $X = \{X_t, \mathcal{F}_t^{X}, t \geq 0\}$. Let $\mathring{\mathcal{F}}_t^X = $augmention of $\mathcal{F}_t^X$.

    Assume $\mathsf{P}(X_0 = 0 ) = 1$.

    Define $W_t \triangleq \int_0^t \mathrm{sgn}(X_s) \diff X_s$.

    Then $W_t$ is a B.M. with respect to $\mathring{\mathcal{F}}_t^X$. 

    Also 
    $$
        \int_0^t \mathrm{sgn}(X_s) \diff W_s = \int_0^t \mathrm{sgn}(X_s) \mathrm{sgn}(X_s)\diff X_s = \int_0^t 1 \diff X_s = X_t.
    $$

    Hence $(X,W)$, $(\Omega,\mathcal{F},\mathsf{P})$, $\mathring{\mathcal{F}}_t^X$ is a weak solution.

    \textbf{No strong solution.}

    Given $(\Omega, \mathcal{F}, \mathsf{P})$ and a B.M. $W$, if $X$ is a strong solution, then 
    $$
    X_t \in \mathcal{F}_t^{W} \subseteq \mathring{\mathcal{F}}_t^W \implies \mathcal{F}_t^X \subseteq \mathring{\mathcal{F}}_t^W.
    $$

    By $X$ is a B.M., by Tanaka's formula,
    $$
    \abs{ X_t }= \abs{ X_0 } + \int_0^t \mathrm{sgn}(X_s) \diff X_s + L_t^X(0).
    $$

    By SDE, we get
    $$
    \int_0^t \mathrm{sgn}(X_s) \diff X_s = \int_0^t \mathrm{sgn}(X_s)^2 \diff W_s = W_t.
    $$  
    $$
    \implies W_t = \abs{ X_t }  - L_t^X(0).
    $$
    $$
    L_t^X(0) = \int_0^t \delta_0 (X_s) \diff s = \lim_{\varepsilon \downarrow 0} \frac{1}{2\varepsilon} \int_0^t \mathsf{1}_{\abs{ X_s } \leq \varepsilon} \diff s = \lim_{\varepsilon \downarrow 0} \frac{1}{2\varepsilon} \abs{ \{0 \leq s \leq y \colon \abs{ X_s } \leq \varepsilon\} }.
    $$

    Hence $W_t = \abs{ X_t } - \lim_{\varepsilon \downarrow 0} \frac{1}{2 \varepsilon} \abs{ \{0 \leq s \leq t \colon \abs{ X_s } \leq \varepsilon\} } \in \mathcal{F}_t^{\abs{ X }}$.

    $\implies \mathring{\mathcal{F}}_t^W \subseteq \mathring{\mathcal{F}}_t^{\abs{ X }}$.

    $\implies \mathring{\mathcal{F}}_t^X \subseteq \mathring{\mathcal{F}}_t^{\abs{ X }}$. Contradiction.
}

$$
\begin{cases}
    \text{strong uniqueness} \implies \text{weak uniqueness}\\
    \text{weak existence}+\text{pathwise uniqueness} \implies \text{strong existence}.
\end{cases}
$$

\subsection{Girsanov Theorem}
    Given $(\Omega, \mathcal{F},\mathsf{P})$ and a d-dim B.M. $W=\{W_t,\mathcal{F}_t, t \geq 0\}$ with $\mathsf{P}(W_0 = 0) = 1$. Let $X=\{X_t,\mathcal{F}_t,t \geq 0\}$ be a d-dim measurable, adapted process such that $\forall~ T \geq 0$, $\forall~ 1 \leq i \leq d$, $\mathsf{P}\left( \int_0^T (X_t^{(i)})^2 \diff t < \infty \right) = 1$. Then 
    $$
    \int_0^t X_s^{(i)} \diff W_s^{(i)}
    $$
    is well-defined. Set 
    $$
    Z_t(X) = \exp \left( \int_0^t X_s \diff W_s  - \frac{1}{2}\int_0^t \norm{ X_s }^2 \diff s \right).
    $$

    Assume $Z(X)$ is a martingale i.e. $\mathbb{E}Z_t(X) = 1$, $\forall~ t \geq 0$.

    Define a probability measure $\tilde{\mathsf{P}}_T$ for any $T > 0$ by 
    $$
    \tilde{\mathsf{P}}_T (A) = \mathbb{E}\left( Z_T(X) \cdot \mathsf{1}_A \right)
    $$
    for all $A \in \mathcal{F}_T$, $\frac{\diff \tilde{\mathsf{P}}_T}{\diff \mathsf{P}} = Z_T(X)$.


\rmkb{
    If $A \in \mathcal{F}_t \subseteq \mathcal{F}_T$ for $0 \leq t \leq T$, we must
    $$
    \begin{aligned}
        \tilde{\mathsf{P}}_T(A) = \mathbb{E}\left( \mathsf{1}_A Z_T(X) \right),\\
        \tilde{\mathsf{P}}_t(A) = \mathbb{E}\left( \mathsf{1}_A Z_t(X) \right).
    \end{aligned}
    $$
    Note that 
    $$
    \mathbb{E}(\mathsf{1}_A Z_T(X)) = \mathbb{E}\left[ \mathbb{E}\left(\mathsf{1}_A Z_T(X) \mid \mathcal{F}_t\right) \right] = \mathbb{E}\left[ \mathsf{1}_A \mathbb{E}\left( Z_T(X) \mid \mathcal{F}_t \right) \right] = \mathbb{E}\left[ \mathsf{1}_A Z_t(X) \right].
    $$
    Hence $\tilde{\mathsf{P}}_T|_{\mathcal{F}_t} = \tilde{\mathsf{P}}_t$.
}

\thm{Girsanov Theorem}{
    Assume $Z(X)$ is a martingale. Define $\tilde{W} = \{\tilde{W}_t, t \geq 0\}$ by $\tilde{W}_t^{(i)} = W_t^{(i)} - \int_0^t X_s^{(i)} \diff s$ for $1 \leq i \leq d$, and $t \geq 0$. Then for each fixed $T > 0$, $\{\tilde{W}_t, \mathcal{F}_t, 0 \leq t \leq T\}$ is a d-dim B.M. on $(\Omega,\mathcal{F}_T,\tilde{\mathsf{P}}_T)$.
}

\lemnn{Fix $T > 0$. If $0 \leq s \leq t \leq T$ and $Y \in \mathcal{F}_t$ with $\tilde{E}_T \abs{ Y } < \infty$. Then
    $$
    \tilde{E}_T \left( Y \mid \mathcal{F}_s \right) = \frac{1}{Z_s}\mathbb{E}\left( YZ_t \mid \mathcal{F}_s \right)
    $$
    a.s. $\tilde{\mathsf{P}}_T$ and $\mathsf{P}$.

    $\tilde{\mathsf{P}}_T << \mathsf{P}$ and $\mathsf{P} << \tilde{\mathsf{P}}_T$. $\tilde{\mathsf{P}}_T(A) = \mathbb{E}\left( \mathsf{1}_A Z_T(X) \right)$.

    If $\mathsf{P}(A) = 0 \implies \tilde{\mathsf{P}}_T(A) = 0$.

    If $\tilde{\mathsf{P}}_T(A) = 0 \implies \mathsf{P}(A) = 0$.$(Z_T(X)>0 \text{ a.s.})$.}

    \pf{
        It suffices to prove $\forall~ A \in \mathcal{F}_s$, 
        $$
        \tilde{E}_T(\mathsf{1}_A Y) = \tilde{E}_T(\mathsf{1}_A \frac{1}{Z_s} \mathbb{E}(YZ_t \mid \mathcal{F}_s)).
        $$

        $$
        \begin{aligned}
            RHS &= \tilde{E}_s (\mathsf{1}_A \frac{1}{Z_s} \mathbb{E}(Y Z_t \mid \mathcal{F}_s)) \\
            &= \mathbb{E}(\mathsf{1}_A \mathbb{E}(YZ_t \mid \mathcal{F}_s))\\
            &= \mathbb{E}(\mathsf{1}_A YZ_t)\\
            &=\tilde{E}_t(\mathsf{1}_A Y)\\
            &= \tilde{E}_T(\mathsf{1}_A Y).
        \end{aligned}
        $$
}

\prop{
    Define $M_T^{c,loc}=\{M\colon M\text{ is a continuous local martingale on }(\Omega,\mathcal{F}_T, \mathsf{P})\}$.

    $\tilde{M}_T^{c,loc}=\{\text{continuous local martingale on }(\Omega,\mathcal{F}_T,\tilde{\mathsf{P}}_T)\}$.

    \begin{itemize}
        \item[(1)] If $M \in M_T^{c,loc}$, then $$
        \tilde{M}_t \triangleq M_t - \sum_{i=1}^d \int_0^t X_s^{(i)} \diff [M,W^{(i)}]_s
        $$
        for $0 \leq t \leq T$ is in $\tilde{M}_T^{c,loc}$.
        \item[(2)] If $N \in M_T^{c,loc}$ and 
        $$
        \tilde{N}_t = N_t - \sum_{i=1}^{d}\int_0^t X_s^{(i)} \diff [N,W^{(i)}]_s,
        $$
        then 
        $$
        [\tilde{M},\tilde{N}]_t = [M,N]_t
        $$
        for all $0 \leq t \leq T$, a.s. $\mathsf{P}$ and $\tilde{\mathsf{P}}_T$.
    \end{itemize}
}


\pf{
    of Girsanov Theorem:

    If $M = W^{(j)}$ with $1 \leq j \leq d$, then 
    $$
    \tilde{M}_t = W_t^{(j)} - \int_0^t X_s^{(j)}\diff W_s^{(j)} \in \tilde{M}_T^{c,loc}
    $$
    $$
    \implies \tilde{W}_t^{(j)} \in \tilde{M}_T^{c,loc}.
    $$

    Set $N = W^{(k)} \implies \tilde{N}_t = \tilde{W}_t^{(k)} \in \tilde{M}_T^{c,loc}$ and $[\tilde{M},\tilde{N}]_t = [M,N]_t = \delta_{jk}t$.
    
    $\implies [\tilde{W}^{(j)},\tilde{W}^{(k)}]_t = \delta_{jk}t$ 

    $\implies \tilde{W}$ is a d-dim B.M..
}

\pf{
    of Prop. Assume $M,N$ are bounded martingale with bounded quadrastic variation.
    $$
    \left[ \tau_n^{M} =  \inf\left\{ t \geq 0 \colon \abs{ M_t } \geq n \text{ or } [M]_{t} \geq n \right\}\right]
    $$

    Consdier $t \leq \tau_n^M \wedge \tau_n^N$.

    Assume also that $Z_t(X)$ and $\sum_{i=1}^d\int_0^t(X_s^{(j)})^2 \diff s$ are bounded for all $ t \geq 0$ a.s.

    By Kunita \& Watanabe inequality
    $$
    \begin{aligned}
    \abs{ \int_0^t X_s^{(j)}\diff [M,W^{(j)}]_s }^2
        &\leq \int_0^t(X_s^{(j)})^2 \diff 
        [W^{(j)}]_s \cdot \int_0^t 1^2 \diff [M]_s \\
        &=\int_0^t(X_s^{(j)})^2 \diff s \cdot [M]_t \\
        & \leq C \quad \text{a.s.}.
    \end{aligned}
    $$

    Hence $\tilde{M}_t$ is also bounded for all $t \geq 0$. To prove $\tilde{M} \in \tilde{M}_T^{c,loc}$, we only need to show that 
    $$
    \tilde{E}_T(\tilde{M}_t \mid \mathcal{F}_s) = \tilde{M}_s \quad \text{ a.s.}
    $$
    $$
    LHS = (\text{by Lemma}) = \frac{1}{Z_s(X)}\mathbb{E}(\tilde{M}_t Z_t(X) \mid \mathcal{F}_s).
    $$

    By IBP, we get 
    $$
    \tilde{M}_t Z_t(X) = \tilde{M}_0Z_0(X) + \int_0^t \tilde{M}_s \diff Z_s(X) + \int_0^t Z_s(X) \diff \tilde{M}_s + [\tilde{M},Z]_t.
    $$
    Recall $Z_t(X) = e^{\int_0^t X_s \diff W_s - \frac{1}{2} \int_0^t \norm{ X_s }^2 \diff s}$,

    we have 
    $$
    \diff Z_t(X) = Z_t(X) X_t \cdot \diff W_t = \sum_{i=1}^d Z_t(X) X_t^{(i)}\diff W_t^{(i)}.
    $$

    Recall $\diff \tilde{M}_t =\diff M_t - \sum_{i=1}^d X_t^{((i))} \diff [M,W^{(i)}]_t$.

    So $$
    \begin{aligned}
        &Z_t(X) \tilde{M}_t \\=& \int_0^t \tilde{M}_s \sum_{i=1}^d Z_s(X) X_s^{(i)} \diff  W_s^{(i)} \\
        &+ \int_0^t Z_s(X)\left[ \diff M_s - \sum_{i=1}^d X_s^{(i)} \diff [M,W^{(i)}]_s \right] + \int_0^t Z_s(X) X_s^{(i)}\diff [M,W^{(i)}]_s \\
        =&\int_0^t \tilde{M}_s \sum_{i=1}^d Z_s(X) X_s^{(i)} \diff W_s^{(i)} + \int_0^t Z_s(X) \diff M_s.
    \end{aligned}
    $$

    Hence $\tilde{M}_t Z_t(X)$ is a martingale.
    $$
    \implies LHS = \frac{1}{Z_s(X)}\tilde{M}_s Z_s(X) = RHS.
    $$
    
    (2) $\diff \tilde{N}_t = \diff N_t = \sum_{i=1}^d X_t^{(i)} \diff [N,W^{(i)}]_t$.
    
    $\diff \tilde{M}_t = \diff M_t - \sum_{i=1}^d X_t^{(i)} \diff [M,W^{(i)}]_t$.
    $$
    \begin{aligned}
    \tilde{M}_t \tilde{N}_t =& \tilde{M}_0 \tilde{N}_0 + \int_0^t \tilde{M}_s \diff \tilde{N}_s + \int_0^t \tilde{N}_s \diff \tilde{M}_s + [\tilde{M},\tilde{N}]_t\\
    =&0 + \int_0^t \tilde{M}_s \left( \diff W_s -\sum_{i=1}^d X_s^{(i)} \diff [N,W^{(i)}]_s \right)\\
    &+\int_0^t \tilde{N}_s \left( \diff M_s - \sum_{i=1}^d X_s^{(i)}\diff [M,W^{(i)}]_s \right) + [M,N]_t.    \end{aligned}
    $$

    Then 
    $$
    \begin{aligned}
        \tilde{M}_t \tilde{N}_t - [M,N]_t =& \int_0^t \tilde{M}_s \diff N_s + \int_0^t \tilde{N}_s \diff M_s - \sum_{i=1}^d \int_0^t X_s^{(i)}  \tilde{M}_s \diff [N,W^{(i)}]_s \\
        &- \sum_{i=1}^d \int_0^t X_s^{(i)}\tilde{N}_s \diff [M,W^{(i)}]_s.
    \end{aligned}
    $$

    Next 
    $$
    \begin{aligned}
        &Z_t(X) \left[ \tilde{M}_t - \tilde{N}_t - [M,N]_t \right]\\
        =&\int_0^t Z_s(X)\left[ \tilde{M}_s \diff N_s + \tilde{N}_s \diff M_s - \sum_{i=1}^d X_s^{(i)}\left( \tilde{M}_s \diff [N,W^{(i)}]_s +\tilde{N}_s \diff [M,W^{(i)}]_s \right) \right]\\
        &+ \int_0^t (\tilde{M}_s \tilde{N}_s - [M,N]_s)\sum_{i=1}^d Z_s(X) X_s^{(i)}\diff W_s^{(i)}\\
        &+ \int_0^t \sum_{i=1}^d Z_s(X) X_s^{(i)}\left( \tilde{M}_s \diff [N,W^{(i)}]_s + \tilde{N}_s \diff [M,W^{(i)}]_s \right).
    \end{aligned}
    $$

    Therefore 
    $$
    \tilde{E}_T(\tilde{M}_t \tilde{N}_t - [M,N]_t \mid \mathcal{F}_s) = \frac{1}{Z_s(X)}\mathbb{E}(Z_t(X)(\tilde{M}_t \tilde{N}_t - [M,N]_t) \mid \mathcal{F}_s) = (\tilde{M}_s \tilde{N}_s - [M,N]_s).
    $$

    This proves $[\tilde{M},\tilde{N}]_t = [M,N]_t$.
}

Weak solutions constructed by means of the Girsanov Theorem.

\prop{
    Consider the SDE 
    $$
    \diff X_t = \diff W_t + b(t,X_t)\diff t
    $$
    for $0 \leq t \leq T$, W is B.M. and $b(t,x)$ is Borel-measurable with
    $$
    \norm{ b(t,x) } \leq K(1 + \norm{x })
    $$
    for all $0 \leq t \leq T$, $x\in \mathbb{R}^d$.

    Then for any probability measure $\mu$ on $(\mathbb{R}^d, \mathcal{B}(\mathbb{R}^d))$ the SDE has a weak solution ($X,W$), $(\Omega,\mathcal{F},\mathsf{P})$, $\{\mathcal{F}_t\}$ such that 
    $$
        \mathsf{P}(X_0 \in \cdot) = \mu.
    $$
}

\pf{
    Let $X=\{X_t, \mathcal{F}_t, t \geq 0\}$ be a B.M. on $(\Omega,\mathcal{F})$ with law $(\mathsf{P}^X)_{X\in \mathbb{R}^d}$ s.t.
    $$
    \mathsf{P}^X(X_0 = x) = 1.
    $$

    By corollary 3.5.16, by assuming 
    $$
    \norm{ b(t,x) } \leq K(1 + \norm{ x })
    $$
    we get 
    $$
    Z_t \triangleq \exp\left( \sum_{i=1}^d \int_0^t b_j(s,X_s) \diff X_s^{(j)} - \frac{1}{2}\int_0^t \norm{ b(s,X_s) }^2 \diff s \right)
    $$
    is a martingale.

    By Girsanov Theorem, if we define 
    $$
    Q^X(A) = \mathsf{P}^X(A \cdot Z_T)
    $$
    for all $A \in \mathcal{F}_T$, then
    $$
    W_t \triangleq X_t - X_0 - \int_0^t b(s,X_s) \diff s
    $$
    is a B.M. with $Q^X(W_0 = 0) = 1$.

    Then $X_t = X_0 + W_t + \int_0^t b(s,X_s) \diff s$ for all $0 \leq s \leq T$.

    Note that $\mathsf{P}^X(X_0 = x) = 1$.

    Then $Q^X(X_0 = x) = 1$.

    Define $Q^\mu (A) = \int_{\mathbb{R}^d} Q^X(A)\mu(\diff x)$.

    Then $Q^{\mu}(W_0 = 0) = 1$ and 
    $$
    \begin{aligned}
    Q^\mu(X_0 \in A) &= \int_{\mathbb{R}^d}Q^X(X_0\in A) \mu (\diff x)\\
    &=\int_{\mathbb{R}^d} \mathsf{1_{x\in A}}\mu(\diff x) = \mu(A) .
    \end{aligned}
    $$

    So under $Q^\mu$ we get $X_0 \sim \mu$ and $(X,W)$, $(\Omega,\mathcal{F},Q^\mu)$, $\{\mathcal{F}_t\}$.

}

(1) pathwise uniqueness $\implies$ weak uniqueness.

\pf{
    If $$
    \begin{cases}
        (X^{(1)},W^{(1)})\\
        (\Omega_1,\mathcal{F}_1,v_1)\\
        (\mathcal{F}_t^{(1)})_{t \geq 0},
    \end{cases}
    \begin{cases}
        (X^{(2)},W^{(2)})\\
        (\Omega_2,\mathcal{F}_2,v_2)\\
        (\mathcal{F}_t^{(2)})_{t \geq 0},
    \end{cases}
    $$
    are two weak solutions.

    If $v_1(X_0^{(1)} \in \cdot) = v_2 (X_0^{(2)} \in \cdot)$, then we will show that
    $$
    v_1(X^{(1)} \in \cdot) = v_2 (X^{(2)} \in \cdot).
    $$

    In fact, we will prove that if 
    $$
    Y^{(j)} = X^{(j)}-X_0^{(j)},
    $$  
    then 
    $$
    v_1((X_0^{(1)},W^{(1)},Y^{(1)}) \in A) = v_2 ((X_0^{(2)},W^{(2)},Y^{(2)}) \in A)
    $$
    for all $A\in \mathcal{B}(\Theta)$ with $\Theta = \mathbb{R}^d \times C[0,\infty)^r \times C[0,\infty)^d$ and $\mathcal{B}(\Theta) = \mathcal{B}(\mathbb{R}^d)\otimes \mathcal{B}(C[0,\infty)^r) \otimes \mathcal{B}(C[0,\infty)^d)$.

    Let $\mu(\cdot) = v_1(X_0^{(1)} \in \cdot) = v_2(X_0^{(2)} \in \cdot)$.

    Let $P_*$ denote the wiener measure on $C[0,\infty)^r$ .

    Define $(\Omega,\mathcal{F})$ by $\Omega=\Theta\times C[0,\infty)^d$, $\mathcal{F} = \mathcal{B}(\Theta)\otimes \mathcal{B}(C[0,\infty)^d)$.

    For each $\omega \in \Omega$, we write $\omega = (x,w,y_1,y_2)$.

    Define 
    $$
    \mathsf{P}(\diff \omega) = \mu(\diff x)P_*(\diff w)Q_1(x,w,\diff y_1)Q_2(x,w,\diff y_2),
    $$
    where $Q_j(x,w,y_i)$ is the regular conditional probability s.t. 
    \begin{itemize}
        \item[(1)] $\forall~ x \in \mathbb{R}^d$, $w \in C[0,\infty)^r$, $Q_j(x,w,\cdot)$ is a probability measure on $(C[0,\infty)^d,\mathcal{B}(C[0,\infty)^d)).$
        \item[(2)] $\forall~ F \in  \mathcal{B}(C[0,\infty)^d)$, the map $(x,w)\to Q_j(x,w;F)$ is measurable w.r.t $\mathcal{B}(\mathbb{R}^d)\otimes \mathcal{B}(C[0,\infty)^r)$.
        \item[(3)] $\forall~G \in \mathcal{B}(\mathbb{R}^d)\otimes \mathcal{B}(C[0,\infty)^r)$, $\forall~F\in \mathcal{B}(C[0,\infty)^d)$, 
        $$
        v_j(X_0^{(j)}, W^{(j)} \in G, Y^{(j)} \in F) = \int_G \mu(\diff x) P_*(\diff w) Q_j(x,w;F)
        $$
        i.e. $Q_j(x,w;F) = v_j(Y^{(j)}\in F \mid X_0^{(j)}= x , W^{(j)}= w )$.
    \end{itemize}

    Let $G_t \triangleq \sigma((x,w(s),y_1(s),y_2(s)) \colon x\in \mathbb{R}^d, 0 \leq s \leq t )$.

    Let $F_t \triangleq $ augmention of $\sigma_t$.

    Now we conclude that 
    $$
    \begin{aligned}
        \mathsf{P}\left( \{\omega \in \Omega \colon (x,w,y_1)\in A\} \right) &= 
        \int \mathsf{P}(\diff \omega) \mathsf{1}_{(x,w,y_1)\in A} \\
        &= \int_{A} \mu(\diff x) P_*(\diff w)Q_1(x,w;\diff y_1)\\
        &= v_1((X_0^{(1)} ,W^{(1)}, Y^{(1)})\in A).
    \end{aligned}
    $$

    So 
    $$
    (x+y_1,w) \xlongequal{d} (X^{(1)},X^{(1)})
    $$
    and 
    $$
    (x+y_2,w) \xlongequal{d} (X^{(2)},W^{(2)}).
    $$

    pathwise uniqueness implies
    $$
    \mathsf{P}(x+y_1(t) = x+y_2(t),\forall~ t \geq 0) = 1.
    $$
    $$
    \begin{aligned}
        v_1((X_0^{(1)},W^{(1)},Y^{(1)})\in A) &=
        \mathsf{P}(\{\omega \in \Omega\colon (x,w,y_1)\in A\}) \\
        &= \mathsf{P}(\{\omega \in \Omega \colon (x,w,y_2)\in A\}) \\
        &=v_2((X_0^{(2)},W^{(2)},Y^{(2)})\in A).
    \end{aligned}
    $$
}

\section{Harmonic Function}

$\forall~ $smooth function $f:\mathbb{R}^d \to \mathbb{R}$,
$$
    \frac{\Delta}{2} f(x) = \lim_{t \downarrow 0} \frac{E^x[f(W_t) - f(x)]}{t}.
$$
We call $\frac{\Delta}{2}$ the generator of Brownian mothion.
$$
\Delta u \triangleq \sum_{i=1}^{d} \frac{\partial^2 u}{\partial x_i^2}.
$$

\pf{
    $$f(W_t) = f(W_0)+\int_0^t \sum_{i=1}^d \frac{\partial f}{\partial x_i}(W_s) \diff W_s^{(i)} + \frac{1}{2}\int_0^t \sum_{i=1}^d \frac{\partial ^2 f}{\partial x_i^2}(W_s) \diff s.$$

    So 
    $$
        \frac{1}{t}E^x[f(W_t) - f(W_0)] = \frac{1}{2} \frac{1}{t} \int_0^t E[\Delta f(W_s)] \diff s.
    $$
    Let $t \downarrow 0$, we get the limit is $\frac{1}{2}\Delta f(x)$.
}

A function $u$ defined on some open subset $D \subseteq \mathbb{R}^d$ is called harmonic if $u \in C^2(D)$ and $\Delta u = 0$. 

Let $\{W_t,\mathcal{F}_t, 0 \leq t < \infty\}$, $(\Omega,\mathcal{F})$ and $\{\mathsf{P}^x\}_{x\in \mathbb{R}^d}$ is a d-dim B.M. such that 
$$
    \mathsf{P}^x(W_0 = x) = 1.
$$

For any $D \subseteq \mathbb{R}^d$ open, we set $\tau _D = \inf\{t \geq 0: W_t \notin D\}$.

If $D$ is bounded, then 
$$
    \mathsf{P}^x(\tau_D < \infty) = 1, \quad \forall~ x \in D,
$$
and $B_{\tau_D} \in \partial D$ with $\partial D$ = topological boundary, $\overline{D} = D \cup \partial D$ is the closure of $D$.

Let $B_r = \{ x\in \mathbb{R}^d: \norm{ x } <r\}$, $B_r(x) = \{y\in \mathbb{R}^d: \norm{ y-x } < r\}$ be the open ball centered at $0/x$ with radius $r > 0$. Volume 
$$V_r = \frac{2 r^d \pi^{\frac{d}{2}}}{d \Gamma(\frac{d}{2})}.$$

Surface area 
$$
S_r \triangleq \frac{2 r^{d-1}\pi^{\frac{d}{2}}}{\Gamma(\frac{d}{2})}.
$$

Define $\mu_r (\diff x) = \mathsf{P}^0(W_{\tau_{B_r}} \in \diff x)$. Then $\mu_r$ is a measure on $\partial B_r$.

By rotational invariance $\mu_r$ is a uniform measure on $\partial B_r$. In particular, we have
$$
\mu_r = \frac{\sigma_r}{S_r}
$$
where $\sigma_r(\diff x)=$surface measure on $\partial B_r$.

Recall that for any function $f$,
$$
\int_{B_r} f(x) \diff x = \int_0^t \diff  \rho \int_{\partial B_{\rho}} f(x) \sigma_\rho (\diff x) = \int_0^r \diff  \rho \int_{\partial B_{\rho}}f(x) S_\rho \mu_\rho(\diff x).
$$

\defn{Mean Value Property (MVP)}{
    A function $u: D \to \mathbb{R}$ has MVP if $\forall~a \in D$, $\forall~ r > 0$ with $B_r(a) \subseteq D$, we have 
    $$
    u(a) = \int_{\partial B_r} \mu(a+x) u_r(\diff x) = \int_{\partial B_r(a)} u(x) \mu_r (\diff x).
    $$
}

Assume MVP, we have 
$$
\int_{B_r(a)} u(x)\diff x = \int_0^t \diff  \rho \int_{\partial B_{\rho}(a)} u(x)S_\rho \mu_\rho(\diff x) = \int_0^r \diff \rho \cdot u(a)S_\rho = V_r \cdot u(a).
$$

Then 
$$
u(a) = \frac{1}{V(r)}\int_{B_r(a)} u(x) \diff x.
$$

\prop{
    If $u$ is harmonic on $D$, then $u$ has MVP. And if $u$ has MVP on $D$, then $u$ is harmonic \& $u \in C^\infty.$
}

\pf{
    (1). 
    Assume $\Delta u = 0$, $\forall~ x\in D$. Then $\forall~ r > 0 $ with $\overline{B_r(a)} \subseteq D$, $\forall~ a \in D$.

    By It\^o's formula
    $$
    u(W_{t \wedge \tau_{B_r(a)}}) = u(W_0) + \int_0^{t \wedge \tau_{B_r(a)}} \nabla u(W_s) \diff W_s + \frac{1}{2} \int_0^{t\wedge\tau_{B_r(a)}} \Delta u(W_s) \diff s, \quad \forall~ t \geq 0.
    $$

    By $\nabla u(x)$ is bounded on $\overline{B_r(a)}$, we get 
    $$
    E\left[ \int_0^{t \wedge \tau_{B_r(a)}}  \nabla u(W_0) \diff W_s\right].
    $$
    
    Hence 
    $$
    E^a\left[ u(W_{t \wedge \tau_{B_r(a)}}) \right] = u(a).
    $$

    Let $t \to \infty$ to get 
    $$
    \begin{aligned}
        u(a) = \lim_{t \to \infty}E^a\left[ u(W_{t\wedge \tau_{B_r(a)}}) \right] &= E^a\left[ u(W_{\tau_{B_r(a)}}) \right]\\
        &=E^0 \left[ u(W_{\tau_{B_r}} + a) \right]\\
        &=\int_{\partial B_r}u(x+a) \mu_r(\diff x).
    \end{aligned}
    $$

    (2) If $u$ has MVP, then $\forall~ a \in D$, $\forall~ r >0$ with $\overline{B_r(a)} \subseteq D$,
    $$
    u(a) = \int_{\partial B_r(a)} u(x)\mu_r(\diff x).
    $$

    For any $\varepsilon > 0$, define $g_\varepsilon : \mathbb{R}^d \to \mathbb{R}$ by 
    $$
    g_{\varepsilon}(x) = \begin{cases}
        c(\varepsilon)e^{-\frac{1}{\varepsilon^2 - \norm{ x }^2}}, &\norm{ x } < \varepsilon\\
        0,& \norm{ x }\geq \varepsilon
        .
    \end{cases}
    $$
    where $c(\varepsilon)>0$ is a constant s.t.
    $$
    \int_{\mathbb{R}^d} g_{\varepsilon}(x) \diff x = \int_{B_\varepsilon}g_{\varepsilon}(x)\diff x = 1.
    $$

    Mollifier. For any $a \in D$ and $\varepsilon > 0$ with $\overline{B_{\varepsilon}(a)} \subseteq D$, define 
    $$
    u_{\varepsilon}(a) \triangleq \int_{B_\varepsilon}u(a+x)g_{\varepsilon}(x) \diff x = \int_{\mathbb{R}^d} u(y) g_{\varepsilon}(y-a) \diff  y.
    $$
    Then $u_{\varepsilon}(a) \in C^{\infty}$.

    $$
    \begin{aligned}
    \lim_{\delta\downarrow 0}\frac{u_\varepsilon (a+\delta) - u_\varepsilon(a)}{\delta} &= \lim_{\delta\downarrow 0} \int_{\mathbb{R}^d}u(y) \frac{g_{\varepsilon}(y-a-\delta) - g_{\varepsilon}(y-a)}{\delta} \diff y \\
    &= -\int_{\mathbb{R}^d}u(y)\left( \frac{\diff }{\diff  x}g_{\varepsilon}(x)\Big|_{x=y-a} \right) \diff y.
    \end{aligned}
    $$

    And
    $$
    \begin{aligned}
    u_{\varepsilon}(a) &= c(\varepsilon)\int_{B_\varepsilon}u(a+x)e^{-\frac{1}{\varepsilon^2 - \norm{ x }^2}}\diff x \\
    &=c(\varepsilon)\int_0^\varepsilon S_{\rho}\diff \rho \int_{\partial B_{\rho}} u(a+x) e^{-\frac{1}{\varepsilon^2 - \rho^2}} \mu_\rho(\diff x)\\
    &=c(\varepsilon)\int_0^\varepsilon S_{\rho}\diff \rho \cdot e^{-\frac{1}{\varepsilon^2 - \rho^2}} u (a) \\
    &=u(a)\int_0^\varepsilon S_{\rho}c(\varepsilon)e^{-\frac{1}{\varepsilon^2 -\rho^2}}\diff \rho \\
    &= u(a)\int_{B_\varepsilon}c(\varepsilon)e^{-\frac{1}{\varepsilon^2 - \norm{ x }^2}}\diff x = u(0).
    \end{aligned}
    $$
    
    So $u(a) = u_{\varepsilon}(0)$ for all $a \in D$ with $\overline{B_\varepsilon(a)} \subseteq D$.

    Then $u \in C^\infty(D)$.

    Next, we will prove $\Delta u = 0$. $\forall~ a \in D$, $\exists~\varepsilon>0$ s.t. $\overline{B_\varepsilon(a)} \subseteq D$. By Taylor's expansion at $a$,
    $$
    u(a+y) = u(a) + \sum_{i=1}^d y_i \frac{\partial u}{\partial x_i}(a) + \frac{1}{2}\sum_{i=1}^d \sum_{j=1}^d y_i y_j \frac{\partial ^2 u}{\partial x_i \partial x_j}(a) + o(\norm{ u }^2), \quad y\in \overline{B}_\varepsilon.
    $$

    By symmetry, 
    $$\int_{\partial B_\varepsilon} y_i \mu_{\varepsilon}(\diff y) = 0, \quad \int_{\partial B_\varepsilon} y_i y_j \mu_\varepsilon(\diff y) = 0,\quad \forall~i \neq j.
    $$

    Now integrate $u(a+y)$ over $\partial B_\varepsilon$ to obtain
    $$
    \int_{\partial B_\varepsilon} u(a+y) \mu_\varepsilon (\diff y) = u(a) + 0 + \int_{\partial B_\varepsilon} \frac{1}{2}\Delta u(a)\sum_{i=1}^d y_i^2 \mu_\varepsilon (\diff  y) + o(\varepsilon^2).
    $$
    
    For 
    $$
    \frac{1}{2} \sum_{i=1}^d \int_{\partial B_\varepsilon} \frac{\partial ^2 u}{\partial x_i^2}(a) \cdot y_i^2 \mu_\varepsilon(\diff y).
    $$

    Notice that
    $$
    \int_{\partial B_\varepsilon} y_i^2 \mu_\varepsilon(\diff  y) = \frac{1}{d} \sum_{i=1}^d \int_{\partial B_\varepsilon} y_i^2 \mu_\varepsilon(\diff y) = \frac{1}{d} \int_{\partial B_\varepsilon}\norm{ y }^2 \mu_\varepsilon(\diff  y) = \frac{1}{d}\varepsilon^2 \cdot 1 = \frac{1}{d}\varepsilon^2.
    $$

    Hence 
    $$
    \frac{1}{2} \sum_{i=1}^d \int_{\partial B_\varepsilon} \frac{\partial ^2 u}{\partial x_i^2}(a) \cdot y_i^2 \mu_\varepsilon(\diff y)= \frac{1}{2}\sum_{i=1}^d \frac{\partial^2 u}{\partial x_i^2}(a) \cdot \frac{\varepsilon^2}{d} =\frac{1}{2}\Delta u(a)\cdot \frac{\varepsilon^2}{d}.
    $$

    Now we conclude 
    $$
    u(a) = \int_{\partial B_\varepsilon}u(a+y)\mu_\varepsilon(\diff  y) = u(a) + \frac{\Delta}{2} u(a)\cdot \frac{\varepsilon^k}{d} + o(\varepsilon^2).
    $$

    So $ 0 = \frac{\Delta}{2} u(a) \cdot \frac{1}{d} + o(1)$, be letting $\varepsilon \downarrow 0$ to get $\Delta u(a) = 0.$
    
}

\subsection{Dirichlet Problem}
Given an open set $D \subseteq \mathbb{R}^d$, and a continuous function $f$ on $\partial D $, the Dirichlet probability (D.f) is to find $u \in C(\overline{D})$ s.t.
$$
    \begin{cases}
        \Delta u = 0, & \text{in } D \\
        u =f ,&\text{on } \partial D.
    \end{cases}
$$

$$
\frac{\partial u(t,x)}{\partial t} = \frac{\Delta}{2} u (t,x).
$$

Heat Equation $u(t,x) = $temperature at $x \in D$ at time $t$.

\textbf{Existence.}

Let $u(x) = E^x[f(W_{\tau_D})]$ for all $x\in \overline{D}$, provided that $E^x \abs{ f(W_{\tau_D}) }< \infty$, $\forall~ x\in D$. Then for $x\in \partial D$, we get 
$$
\mathsf{P}^x(\tau_D = 0) = 1
$$
by $\tau_D = \inf \{t \geq 0: W_t \notin D\}$. Then 
$$
\mathsf{P}^x(W_{\tau_D} = x) = 1,
$$
which implies $E^x[f(W_{\tau_D})] = f(x)$.

So $u(x) = f(x)$, $\forall~ x \in \partial D$.

Next, we will prove $\Delta u = 0$, $\forall~ x \in D \Leftrightarrow u$ has MVP.

For all $a \in D$ and $r > 0$ s.t. $\overline{B_r(a)} \subseteq D$, we have 
$$
\begin{aligned}
u(a)= E^a[f(W_{\tau_D})] &= E^a\left[ E^a\left[f(W_{\tau_D})~\middle|~ \mathcal{F}_{\tau_{B_r(a)}}\right] \right]\\
(\tau_{B_r(a)} \leq \tau_D)~ &= E^a\left[ E^{W_{\tau_{B_r(a)}}}[f(W_{\tau_D})] \right]\\
&= E^a \left[ u(W_{\tau_{B_r(a)}}) \right]\\
&= \int_{\partial B_r} u(a+y) \mu_r(\diff y).
\end{aligned}
$$

Therefore $u$ has MVP. It remains to prove $u \in C(\overline{D}) $ i.e. $\forall~ a \in \partial D$
$$
    \lim_{\substack{x\to a\\ x\in D}}u(x) = u(a) =f(a).
$$

$\forall~ a \in \partial D$, we need that 
$$
\lim_{\substack{x\to a\\ x\in D}} E^x[f(W_{\tau_D})] = f(a).
$$








\textbf{Uniqueness.}

If $f$ is bdd and $\mathsf{P}^a(\tau_D < \infty) = 1$, $\forall~ a \in D$. Then any bounded solution $u$ to Dirichlet Problem (D,f) is given by $u(a) = E^a[f(W_{\tau_D})]$.

\pf{
    Let $u \in C(\overline{D})$ be bounded and
    $$
    \begin{cases}
        \Delta u = 0, & \text{in } D\\
        u=f,&\text{on } \partial D .
    \end{cases}
    $$
    For each $n \geq 1$, we let 
    $$
    \begin{aligned}
        D_n &= \{x\in D : \inf_{y\in \partial D}\norm{ x-y } > \frac{1}{n}\}\\
        &=\{x\in D : d(x,\partial D) > \frac{1}{n}\}.
    \end{aligned}
    $$

    By It\^o's formula
    $$
    \begin{aligned}
        u(W_{t\wedge \tau_{D_n} \wedge \tau_{B_n}}) = u(W_0) + \sum_{i=1}^d \int_0^{t \wedge \tau_{D_n} \wedge \tau_{B_n}} \frac{\partial u}{\partial x_i}(W_s) \diff W_s^{(i)}.
    \end{aligned}
    $$

    Since $\frac{\partial u}{\partial x_i}(W_s)$ is bounded for all $0 \leq s \leq t \wedge \tau_{B_n} \wedge \tau_{D_n}$, we get 
    $$
    E\left[ \int_0^{t\wedge \tau_{D_n} \wedge \tau_{B_n}} \frac{\partial u}{\partial x_i}(W_s)\diff W_s^{(i)} \right] = 0.
    $$

    Then 
    $$
    u(a) = E^a \left[ u(W_{t \wedge \tau_{B_n} \wedge \tau_{D_n}}) \right].
    $$

    By $\mathsf{P}^a(\tau_D < \infty) = 1$, we have ($t\to \infty, n \to \infty$)
    $$
    u(W_{t \wedge \tau_{B_n} \wedge \tau_{D_n}}) \xlongrightarrow{a.s.} u(W_{\tau_D}) = f(W_{\tau_D}).
    $$

    So by bounded convergence theorem, 
    $$
    u(a) = \lim_{n\to \infty} \lim_{t\to \infty} E^a\left[ u(W_{t \wedge \tau_{B_n} \wedge \tau_{D_n}}) \right] = E^a(f(W_{\tau_D})).
    $$

    
}

    To prove $\forall~ a \in \partial D$,
    $$
    \lim_{\substack{x\to a \\ x\in D}} E^x\left[ f(W_{\tau_{D}}) \right] = f(a.)
    $$

    Define $\sigma_D = \inf\{t > 0 \colon W_t \notin D\}$. It is obvious that $\tau_D \leq \sigma_D$. 
    
    \defn{regular point}{
        We say a point $a \in \partial D$ is regular for $D$ if $\mathsf{P}^a(\sigma_D = 0) = 1$.
    }
    
    $\sigma_D = \frac{1}{2} \Leftrightarrow \forall~ 0 < t < \frac{1}{2}$, $W_t \in D$.

    $\sigma_D = 0 \Leftrightarrow \exists~ t_n \downarrow 0$ s.t. $W_{t_n} \notin D$.

    When $d=1$, $D =(a,b)$. Then $\partial D = \{a,b\}$. By LIL, $a$ and $b$ are both regular points.
    $$
    \limsup_{t \downarrow 0} \frac{B_{t} - B_{0}^{n}}{\sqrt{2t \log\log \frac{1}{t}}} = 1 \quad \text{a.s.}
    $$
    $\implies \exists~t_{n}\downarrow 0 $ s.t. 
    $$
    B_{t} - b \geq (\sqrt{2} - \varepsilon) \sqrt{t_{n} \log\log \frac{1}{t_{n}}} > 0.
    $$

    \defn{irrgular point}{
        A point $a \in \partial D$ is irrgular if $\mathsf{P}^{a}(\sigma_{D} = 0) < 1 \Leftrightarrow \mathsf{P}^{a}(\sigma_{D}>0)>0$. 
    }
    By Blumenthal 0-1 law, $\mathsf{P}^{a}(\sigma_{D}>0) \in \{0,1\}$. So if $a \in \partial D$ is irrgular, then $\mathsf{P}^{a}(\sigma_{D} > 0) = 1$.

\thmnn{
    ($d \geq 2$) and fix $a \in \partial D$. The following are equivalent
    \begin{itemize}
        \item[(i)] $$
        \lim_{\substack{x \to a \\ x \in D}} E^{x}\left[ f(W_{\tau_{D}}) \right] = f(a)
        $$
        holds for every bounded, measurable function $f$ on $\partial D$ s.t. $f$ is continuous at $a$.
        \item[(ii)] $a$ is regular.
        \item[(iii)] $\forall~ \varepsilon > 0$, 
        $$
        \lim_{\substack{x\to a \\ x \in D}} \mathsf{P}^{x}(\tau_{D} > \varepsilon) = 0.
        $$ 
        ($\tau_{D}\xlongrightarrow{\mathsf{P}^{x}}0$ as $x \to a$)
    \end{itemize}
}

\pf{
    (iii)$\implies$(i) By continuity of Brownian path,
    $$
    \mathsf{P}^{x}\left( \max_{0 \leq t \leq \varepsilon}\norm{ W_{t} - W_{0} } < r \right) \to 1
    $$
    as $\varepsilon \downarrow 0$,
    and
    $$
    \mathsf{P}^{0}\left( \max_{0 \leq t \leq \varepsilon}\norm{ W_{t}} < r \right) \to 1.
    $$
    Then 
    $$
    \begin{aligned}
    \mathsf{P}^{x}\left( \norm{ W_{\tau_{D}} - W_{0} } < r  \right) &\geq \mathsf{P}^{x}\left( \max_{0 \leq t \leq \varepsilon} \norm{ W_{t} - W_{0} }< r,~ \tau_{D} \leq \varepsilon \right) \\
    &\geq \mathsf{P}^{x}\left( \max_{0 \leq t \leq \varepsilon}\norm{ W_{t}- W_{0} }<r \right)- \mathsf{P}^{x}(\tau_{D} > \varepsilon).
    \end{aligned}
    $$

    Let $x\in D \to a$ to get 
    $$
    \lim_{\substack{x \to a \\ x\in D}} \mathsf{P}^{x}\left( \norm{ W_{\tau_{D}} - W_{0} }<r \right) \geq \mathsf{P}^{0}(\max_{0 \leq t \leq \varepsilon} \norm{ W_{t} }<r) - 0.
    $$

    $\forall~ \varepsilon > 0$, $\exists~ r > 0$ s.t. $\abs{ f(x) - f(a) } < \varepsilon$ for all $\abs{ x-a }<r$,
    $$
    E^{x}\abs{ f(W_{\tau_{D}}) - f(a) } \leq \norm{ f }_{\infty} \cdot E^{x}\left( \mathsf{1}_{\norm{ W_{\tau_{D}} - x }> \frac{r}{2}} \right) + E^{x}\left( \abs{ f(W_{\tau_{D}}) - f(a) } \cdot \mathsf{1}_{\norm{ W_{\tau_{D} - x} } < \frac{r}{2}} \right).
    $$
    If we let $\abs{ x-a } < \frac{r}{2}$, then 
    $$
    \norm{ W_{\tau_{D}} - a  } \leq \norm{  W_{\tau_{D}} - x } + \norm{ x-a } < r.
    $$ 
    $$
    \implies \abs{ f(W_{\tau_{D}}) - f(a) } < \varepsilon.
    $$

    (ii)$\implies$(iii) Assume $a = 0$ .
    $$
    g(x) = \mathsf{P}^x(\sigma_D > \varepsilon).
    $$
    Because $\tau_D \leq \sigma_D$, we get 
    $$
    \limsup_{\substack{x\to 0 \\ x\in D}} \mathsf{P}^x(\tau_D > \varepsilon) \leq \limsup_{\substack{x\to 0 \\ x \in D}} \mathsf{P}^x(\sigma_D > \varepsilon) \leq \limsup_{\substack{x\to 0 \\ x\in D}} g(x).
    $$
    Because $g(0) = 0$, it suffices to prove $g$ is upper semicontinuos, i.e. $\forall~ x_0 \in \mathbb{R}$,
    $$
    \limsup_{x\to x_0} g(x) \leq g(x_0).
    $$

    Note that 
    $$
    g(x) = \mathsf{P}^x(\sigma_D > \varepsilon) = \mathsf{P}^x\left( W_s \in D, ~\forall~ 0 < s \leq \varepsilon \right).
    $$
    $\forall~ 0 < \delta < \varepsilon$, $$g_\delta(x) = \mathsf{P}^x(W_s \in D,~\forall~ \delta \leq s \leq \varepsilon).$$

    Then (1)$\forall~ x \in \mathbb{R}$, $g_{\delta(x)}\downarrow g(x)$. Next, we will show that (2) $x\mapsto g_\delta(x)$ is continuous.

    Then (1)+(2) $\implies g(x)$ is upper-semicontinuos. 

    Fix $\delta \in (0,\varepsilon)$, we have 
    $$
    \begin{aligned}
    g_\delta(x) &= \mathsf{P}^x\left( W_s\in D,~ \delta \leq s \leq \varepsilon \right) \\
    &= E^x\left[ \mathsf{P}^{W_{\delta}}\left( W_s \in D,~ 0 \leq s \leq \varepsilon - \delta  \right) \right] \\
    &= E^x\left[ \mathsf{P}^{W_\delta}\left( \tau_D > \varepsilon - \delta \right) \right]\\
    &= \int_{\mathbb{R}^d} \mathsf{P}_{\delta}(x,y)\mathsf{P}^y \left( \tau_D > \varepsilon - \delta \right) \diff y.
    \end{aligned}
    $$

    So $x\mapsto \int_{\mathbb{R}^d} \mathsf{P}_{\delta}(x,y)\mathsf{P}^y(\tau_D > \varepsilon - \delta) \diff y$ is continuous.

    (i)$\implies$(ii) Assume 0 is not regular, i.e.
    $$
    \mathsf{P}^0(\sigma_D = 0 ) = 0 \implies \mathsf{P}^0 (\sigma_D > 0) = 1.
    $$
    Because starting from 0,
    $$
    \mathsf{P}^0\left( W_t \neq 0,~\forall~t > 0 \right) = 1 \implies \mathsf{P}^0 \left( W_{\sigma_D} = 0 \right) = 0.
    $$
    Hence 
    $$
    \lim_{r \downarrow 0}\mathsf{P}^0\left( W_{\sigma_D} \in B_r \right) = 0.
    $$

    Fix $r > 0$ small s.t.
    $$
    \mathsf{P}^0\left( W_{\sigma_D} \in B_r \right) < \frac{1}{4}.
    $$
    Let $\delta_n \downarrow 0$, $0 < \delta_n < r$. 

    Define
    $$
    \tau_n = \inf \{t \geq 0 \colon \norm{ W_t } \geq \delta_n\}.
    $$

    Then 
    $$
    \mathsf{P}^0\left( \tau_n \downarrow 0 \right) = 1.
    $$

    By 
    $$
    \mathsf{P}^0(\sigma_D > 0) = 1,
    $$
    we have 
    $$
    \lim_{n\to \infty}\mathsf{P}^0(\tau_n < \sigma_D) = 1.
    $$
    Since $\tau_n \downarrow 0$, $\sigma_D > 0$, we have $\{\tau_n < \sigma_D\} \subseteq \{\tau_{n+1} < \sigma_D\}$.

    $$
    \mathsf{P}\left( \bigcup_{n=1}^\infty \{\tau_n < \sigma_D\} \right) = \lim_{n\to \infty}\mathsf{P}(\tau_n < \sigma_D).
    $$

    $$
    \mathsf{P}^0\left( \bigcup_{n=1}^\infty \{\tau_n < \sigma _D\} \right) =1 = \lim_{n\to \infty}\mathsf{P}^0(\tau_n < \sigma_D) = 1.
    $$

    So $\exists~ N \geq 1$ s.t. $\forall~ n \geq N$,
    $$
    \mathsf{P}^0(\tau_n < \sigma_D) \geq \frac{1}{2}.
    $$
}